% Leçon 6 — Grammaire : 如果……，那么……；les nombres ；有的……有的……；是……的 — Chinois 1ère A — Cours signé OnBuch+
% Programme officiel MINESEC. Compiler avec : tectonic ZHA06-grammaire-les-nombres-les-nombres.tex
\newcommand{\DOCMATIERE}{Chinois}
\newcommand{\DOCNIVEAU}{1ère A}
\newcommand{\DOCMODULE}{Module 1 : Vie familiale et intégration sociale}
\newcommand{\DOCLECON}{ZHA06}
\newcommand{\DOCTITRE}{Leçon 6 — Grammaire : 如果……，那么……；les nombres ；有的……有的……；是……的}
% OnBuch+ — préambule « cours signés » ENRICHI (compilé avec tectonic / XeLaTeX)
% Système visuel « Pop » d'OnBuch+ : fond crème (jamais blanc), bordures encre
% épaisses, ombres « dures » décalées, titres Archivo Black en capitales.
% Version MATHÉMATIQUES : graphes (pgfplots/tikz), boîtes pédagogiques
% (définition, propriété, méthode, attention, le savais-tu, exercice résolu,
% exercice, TP, bilan), page de garde et signature OnBuch+ ancrée en bas.
%
% Macros attendues dans main.tex AVANT \input{preamble} :
%   \DOCMATIERE  \DOCNIVEAU  \DOCMODULE  \DOCLECON (numéro)  \DOCTITRE
\documentclass[11pt]{article}

\usepackage{fontspec}
\usepackage[french]{babel} % français = langue principale ; le chinois via \zh{..} / textebox (xeCJK)
\usepackage{xeCJK}
\usepackage{pifont} % \ding{51} \ding{55} utilisés par les modèles
\usepackage[a4paper,margin=2cm,top=3.4cm,bottom=2.8cm,headheight=1.4cm,footskip=1.3cm]{geometry}
\usepackage{amsmath,amssymb,amsfonts}
\usepackage{mathtools} % \xmapsto, \coloneqq... souvent utilisés par les modèles
\usepackage{cancel} % simplifications barrées (bilans de forces)
% \abs{z} (module, valeur absolue) et \norm{u} (norme), utilisés par les modèles
\providecommand{\abs}{}\let\abs\relax\DeclarePairedDelimiter{\abs}{\lvert}{\rvert}
\providecommand{\norm}{}\let\norm\relax\DeclarePairedDelimiter{\norm}{\lVert}{\rVert}
\usepackage[table]{xcolor} % [table] : fournit aussi \rowcolors (alternance de lignes)
\usepackage{pagecolor}
\usepackage{tikz}
\usetikzlibrary{arrows.meta,positioning,calc,shapes.geometric,shapes.misc,shapes.arrows,decorations.pathmorphing,decorations.pathreplacing,decorations.markings,patterns,shadows,mindmap,trees,fit,backgrounds,matrix,intersections,angles,quotes}
\usepackage{pgfplots}
\usepackage[version=4]{mhchem} % équations nucléaires \ce{^{235}_{92}U}
\usepackage[european,siunitx]{circuitikz} % schémas électriques
\pgfplotsset{compat=1.18}
\usepgfplotslibrary{fillbetween,groupplots} % aires entre courbes (calcul intégral)
\usepackage{enumitem}
\usepackage{fancyhdr}
% [most] charge toutes les bibliothèques tcolorbox (listings, minted, raster,
% documentation...) et gonfle inutilement la mémoire TeX sur un document long
% (~300 boîtes) : on ne charge que ce qui est réellement utilisé.
\usepackage[skins,breakable]{tcolorbox}
\usepackage{graphicx}
\usepackage{colortbl}
\usepackage{booktabs}
\usepackage{tabularx}
\usepackage{diagbox} % case d'en-tête diagonale des tableaux à double entrée
% Colonnes extensibles centrée / gauche / droite (C, L, R), souvent utilisées par les modèles
\newcolumntype{C}{>{\centering\arraybackslash}X}
\newcolumntype{L}{>{\raggedright\arraybackslash}X}
\newcolumntype{R}{>{\raggedleft\arraybackslash}X}
\usepackage{array}
\usepackage{multicol}
\usepackage{multirow}
\usepackage{siunitx}
% parse-numbers=false : accepte aussi \SI{2,5 \times 10^{-3}}{...}
\sisetup{locale=FR,output-decimal-marker={,},group-minimum-digits=4,parse-numbers=false}
\DeclareSIUnit\ohm{\ensuremath{\symup{\Omega}}} % U+2126 absent de la police
\DeclareSIUnit\division{div} % réglages d'oscilloscope (ms/div, V/div)
\DeclareSIUnit\tr{tr} % tours (vitesse de rotation, ex. \SI{1500}{\tr\per\minute})
\DeclareSIUnit\molar{M} % concentration molaire (ex. \SI{3}{\molar}, \SI{1}{\micro\molar})
\DeclareSIUnit\an{an} % année (le modèle confond parfois avec \year, un compteur TeX) — ex. \SI{5}{\kilo\meter\per\an}
\DeclareSIUnit\FCFA{FCFA} % franc CFA (ex. \SI{500}{\FCFA\per\kilo\gram})
\DeclareSIUnit\degreeBrix{\textdegree Brix} % teneur en sucre d'un sirop/jus (réfractométrie)
\DeclareSIUnit\month{mois} % le modèle utilise parfois \month, un compteur TeX, comme unité de durée
\usepackage{qrcode}
% \degree : commande fréquemment utilisée par les modèles pour un angle ou une
% température en dehors de \SI{}{} (non fournie par les paquets chargés).
\providecommand{\degree}{^{\circ}}
% \qed (fin de démonstration), utilisé par les modèles sans amsthm
\providecommand{\qed}{\hfill$\square$}
% \barre : double barre des valeurs interdites dans un tableau de variations (tkz-tab)
\providecommand{\barre}{\ensuremath{\|}}
% environnement proof (amsthm non chargé) utilisé par les modèles
\ifdefined\proof\else\newenvironment{proof}[1][Démonstration]{\par\noindent\textit{#1.}\ }{\qed\par}\fi
\usepackage{lastpage}
\usepackage{hyperref}
\hypersetup{hidelinks}

% ============================================================
% Palette « Pop » OnBuch+ — mêmes hex que la classe P (plus_ui.dart)
% ============================================================
\definecolor{poporange}{HTML}{F59321}
\definecolor{popdark}{HTML}{E07A0C}
\definecolor{popcream}{HTML}{FFF6E8}
\definecolor{popink}{HTML}{1C1714}
\definecolor{poppurple}{HTML}{7A5AE0}
\definecolor{popgreen}{HTML}{2E9E6A}
\definecolor{poppink}{HTML}{E0457B}
\definecolor{popblue}{HTML}{2F7DE1}
\definecolor{popgold}{HTML}{C98A12}
\definecolor{popmuted}{HTML}{8A7F76}
\definecolor{popline}{HTML}{EADFD2}
\definecolor{popgray}{HTML}{8A7F76}
\colorlet{popgrey}{popgray}
% Alias tolérants (le modèle nomme parfois les couleurs différemment)
\colorlet{popwhite}{white}
\colorlet{popblack}{popink}
\colorlet{popred}{poppink}
\colorlet{popyellow}{popgold}
% Teintes claires (fonds de tableaux/encadrés) — le modèle nomme parfois
% les couleurs avec un suffixe L (ex. popblueL) sans les avoir définies.
\colorlet{poporangeL}{poporange!20}
\colorlet{popdarkL}{popdark!20}
\colorlet{poppurpleL}{poppurple!20}
\colorlet{popgreenL}{popgreen!20}
\colorlet{poppinkL}{poppink!20}
\colorlet{popblueL}{popblue!20}
\colorlet{popgoldL}{popgold!20}
\colorlet{popredL}{popred!20}
\colorlet{popgrayL}{popgray!20}
% Teintes claires pour fonds de tableaux / zones
\colorlet{popblueL}{popblue!10!white}
\colorlet{poporangeL}{poporange!14!white}
\colorlet{popgreenL}{popgreen!12!white}
\colorlet{poppurpleL}{poppurple!12!white}
\colorlet{poppinkL}{poppink!12!white}
\colorlet{popgoldL}{popgold!14!white}

\pagecolor{popcream}

% ============================================================
% Typographie
% ============================================================
\defaultfontfeatures{Ligatures=TeX}
\setmainfont{PlusJakartaSans-Medium.ttf}[Path=../../../fonts/,BoldFont=PlusJakartaSans-Bold.ttf]
% Symboles, API et emojis absents de la police principale : repli sur DejaVu Sans (caractères actifs)
\newfontface\symfont{DejaVuSans.ttf}[Path=../../../fonts/]
\makeatletter
\newcommand\symactive[1]{\catcode#1=13 \begingroup\lccode`\~=#1 \lowercase{\endgroup\def~}{{\symfont\char#1}}}
\newcommand\symblank[1]{\catcode#1=13 \begingroup\lccode`\~=#1 \lowercase{\endgroup\def~}{}}
\newcommand\symhyph[1]{\catcode#1=13 \begingroup\lccode`\~=#1 \lowercase{\endgroup\def~}{\mbox{-}}}
\newcount\symc
\newcommand\symrange[2]{\symc=#1 \@whilenum\symc<#2 \do{\expandafter\symactive\expandafter{\the\symc}\advance\symc by 1 }}
\newcommand\symblankrange[2]{\symc=#1 \@whilenum\symc<#2 \do{\expandafter\symblank\expandafter{\the\symc}\advance\symc by 1 }}
\makeatother
\symrange{"0250}{"02FF}\symactive{"2713}\symactive{"2714}\symactive{"2717}\symactive{"2718}\symactive{"2610}\symactive{"2611}\symactive{"2612}
\symrange{"2600}{"27BF}
\catcode"2705=13 \begingroup\lccode`\~="2705 \lowercase{\endgroup\def~}{{\symfont\char"2713}}\symblank{"26BD}\symblank{"26A1}
\symrange{"2460}{"257F}\symactive{"223C}\symactive{"2248}\symactive{"2260}
\symhyph{"2011}\symhyph{"2E1A}\symblankrange{"1F300}{"1FAFF}\symblank{"FE0F}
% Caractères chinois : WenQuanYi Zen Hei (sous-ensemble GB2312 + pinyin), gras/italique simulés
\setCJKmainfont{WenQuanYiZenHei-subset.ttf}[Path=../../../fonts/,AutoFakeBold=3,AutoFakeSlant=0.2]
\xeCJKsetup{CJKspace=false}
% Pinyin : voyelles à caron (ǎ ǐ ǒ ǔ ǖ ǘ ǚ ǜ) absentes de la police principale -> caractères actifs
\newfontface\pyfont{WenQuanYiZenHei-subset.ttf}[Path=../../../fonts/]
\newcommand\pyactive[1]{\catcode#1=13 \begingroup\lccode`\~=#1 \lowercase{\endgroup\def~}{{\pyfont\char#1}}}
\pyactive{"01CD}\pyactive{"01CE}\pyactive{"01CF}\pyactive{"01D0}\pyactive{"01D1}\pyactive{"01D2}\pyactive{"01D3}\pyactive{"01D4}\pyactive{"01D5}\pyactive{"01D6}\pyactive{"01D7}\pyactive{"01D8}\pyactive{"01D9}\pyactive{"01DA}\pyactive{"01DB}\pyactive{"01DC}
% Maths en sans-serif (Fira Math) avec chiffres et lettres droites de Plus
% Jakarta, pour que les formules se fondent dans le texte.
\usepackage[math-style=ISO,bold-style=ISO]{unicode-math}
\setmathfont{FiraMath-Regular.otf}[Path=../../../fonts/]
\setmathfont{PlusJakartaSans-Medium.ttf}[Path=../../../fonts/,range={up/num,up/{Latin,latin},\mathup}]
\usepackage{icomma} % virgule décimale sans espace en mode math
% Caractères mathématiques Unicode absents des polices de texte (Plus Jakarta,
% Archivo Black) : rendus par leur équivalent mathématique, en texte comme en
% formule — sinon ils s'affichent en carré vide (ex. « Arithmétique dans ℤ »).
\usepackage{newunicodechar}
\newunicodechar{ℝ}{\ensuremath{\mathbb{R}}}
\newunicodechar{ℤ}{\ensuremath{\mathbb{Z}}}
\newunicodechar{ℂ}{\ensuremath{\mathbb{C}}}
\newunicodechar{ℕ}{\ensuremath{\mathbb{N}}}
\newunicodechar{ℚ}{\ensuremath{\mathbb{Q}}}
\newunicodechar{𝔻}{\ensuremath{\mathbb{D}}}
\newunicodechar{α}{\ensuremath{\alpha}}
\newunicodechar{β}{\ensuremath{\beta}}
\newunicodechar{γ}{\ensuremath{\gamma}}
\newunicodechar{δ}{\ensuremath{\delta}}
\newunicodechar{ε}{\ensuremath{\varepsilon}}
\newunicodechar{θ}{\ensuremath{\theta}}
\newunicodechar{λ}{\ensuremath{\lambda}}
\newunicodechar{μ}{\ensuremath{\mu}}
\newunicodechar{σ}{\ensuremath{\sigma}}
\newunicodechar{τ}{\ensuremath{\tau}}
\newunicodechar{φ}{\ensuremath{\varphi}}
\newunicodechar{ψ}{\ensuremath{\psi}}
\newunicodechar{ω}{\ensuremath{\omega}}
\newunicodechar{Σ}{\ensuremath{\Sigma}}
% majuscules produites par \MakeUppercase dans les titres (α-aminés -> Α)
\newunicodechar{Α}{\ensuremath{\alpha}}
\newunicodechar{Β}{\ensuremath{\beta}}
\newunicodechar{Γ}{\ensuremath{\Gamma}}
\newunicodechar{Δ}{\ensuremath{\Delta}}
\newunicodechar{Ε}{\ensuremath{\varepsilon}}
\newunicodechar{Θ}{\ensuremath{\Theta}}
\newunicodechar{Λ}{\ensuremath{\Lambda}}
\newunicodechar{Μ}{\ensuremath{\mu}}
\newunicodechar{Τ}{\ensuremath{\tau}}
\newunicodechar{Φ}{\ensuremath{\Phi}}
\newunicodechar{Ψ}{\ensuremath{\Psi}}
\newunicodechar{Ω}{\ensuremath{\Omega}}
\newunicodechar{↦}{\ensuremath{\mapsto}}
\newunicodechar{∈}{\ensuremath{\in}}
\newunicodechar{∉}{\ensuremath{\notin}}
\newunicodechar{∧}{\ensuremath{\wedge}}
\newunicodechar{⇔}{\ensuremath{\Leftrightarrow}}
\newunicodechar{⟺}{\ensuremath{\Longleftrightarrow}}
\newunicodechar{⇒}{\ensuremath{\Rightarrow}}
\newunicodechar{⟹}{\ensuremath{\Longrightarrow}}
\newunicodechar{∘}{\ensuremath{\circ}}
\newunicodechar{∪}{\ensuremath{\cup}}
\newunicodechar{∩}{\ensuremath{\cap}}
\newunicodechar{≡}{\ensuremath{\equiv}}
\newunicodechar{⊂}{\ensuremath{\subset}}
\newunicodechar{⊥}{\ensuremath{\perp}}
\newunicodechar{∃}{\ensuremath{\exists}}
\newunicodechar{∀}{\ensuremath{\forall}}
\newunicodechar{∛}{\ensuremath{\sqrt[3]{\,}}}
\newunicodechar{∜}{\ensuremath{\sqrt[4]{\,}}}
\newunicodechar{⃗}{\ensuremath{{}^{\to}}}
\newunicodechar{ⁿ}{\ifmmode^{n}\else\textsuperscript{\ensuremath{n}}\fi}
\newunicodechar{ˣ}{\ifmmode^{x}\else\textsuperscript{\ensuremath{x}}\fi}
\newunicodechar{ᵇ}{\ifmmode^{b}\else\textsuperscript{\ensuremath{b}}\fi}
\newunicodechar{ʳ}{\ifmmode^{r}\else\textsuperscript{\ensuremath{r}}\fi}
\newunicodechar{ᵘ}{\ifmmode^{u}\else\textsuperscript{\ensuremath{u}}\fi}
\newunicodechar{ʸ}{\ifmmode^{y}\else\textsuperscript{\ensuremath{y}}\fi}
\newunicodechar{ᵃ}{\ifmmode^{a}\else\textsuperscript{\ensuremath{a}}\fi}
\newunicodechar{ᵉ}{\ifmmode^{e}\else\textsuperscript{\ensuremath{e}}\fi}
\newunicodechar{ᵏ}{\ifmmode^{k}\else\textsuperscript{\ensuremath{k}}\fi}
\newunicodechar{ᵖ}{\ifmmode^{p}\else\textsuperscript{\ensuremath{p}}\fi}
\newunicodechar{ᴺ}{\ifmmode^{N}\else\textsuperscript{\ensuremath{N}}\fi}
\newunicodechar{ᶜ}{\ifmmode^{c}\else\textsuperscript{\ensuremath{c}}\fi}
\newunicodechar{ʰ}{\ifmmode^{h}\else\textsuperscript{\ensuremath{h}}\fi}
\newunicodechar{ᵠ}{\ifmmode^{q}\else\textsuperscript{\ensuremath{q}}\fi}
\newunicodechar{ₙ}{\ifmmode_{n}\else\textsubscript{\ensuremath{n}}\fi}
\newunicodechar{ₐ}{\ifmmode_{a}\else\textsubscript{\ensuremath{a}}\fi}
\newunicodechar{ᵢ}{\ifmmode_{i}\else\textsubscript{\ensuremath{i}}\fi}
\newunicodechar{ᵣ}{\ifmmode_{r}\else\textsubscript{\ensuremath{r}}\fi}
\newunicodechar{ₖ}{\ifmmode_{k}\else\textsubscript{\ensuremath{k}}\fi}
\newunicodechar{ᵦ}{\ifmmode_{\beta}\else\textsubscript{\ensuremath{\beta}}\fi}
\newunicodechar{ₓ}{\ifmmode_{x}\else\textsubscript{\ensuremath{x}}\fi}
\newfontfamily\popdisplay{ArchivoBlack-Regular.ttf}[Path=../../../fonts/]
\newfontfamily\poplogo{SpaceGrotesk-Bold.ttf}[Path=../../../fonts/]
\newfontfamily\popsemi{PlusJakartaSans-SemiBold.ttf}[Path=../../../fonts/]
\newfontfamily\popxbold{PlusJakartaSans-ExtraBold.ttf}[Path=../../../fonts/]
\newfontfamily\popxboldls{PlusJakartaSans-ExtraBold.ttf}[Path=../../../fonts/,LetterSpace=3.0]

\setlist[enumerate]{leftmargin=1.7em,itemsep=5pt,topsep=4pt}
\setlist[itemize]{leftmargin=1.7em,itemsep=4pt,topsep=4pt,label={\color{poporange}\textbullet}}
\setlength{\parindent}{0pt}
\setlength{\parskip}{5pt}
\renewcommand{\arraystretch}{1.25}

% pgfplots : style Pop par défaut
\pgfplotsset{
  every axis/.append style={axis line style={popink,line width=1pt},tick style={popink},
    grid style={popline,line width=0.5pt},label style={font=\small},tick label style={font=\footnotesize}},
  popcurve/.style={poporange,line width=1.8pt,no markers,smooth},
}
% Même style disponible dans un \draw TikZ simple (hors pgfplots)
\tikzset{popcurve/.style={poporange,line width=1.8pt}}
\tikzset{popgrid/.style={popline,very thin}}
\colorlet{popcurve}{poporange} % le nom du style est parfois utilisé comme couleur

% ============================================================
% Pill / squiggle
% ============================================================
\newcommand{\poppill}[3]{%
  \tikz[baseline=-0.55ex]\node[draw=popink,line width=1.2pt,rounded corners=2.6pt,%
    fill=#2,inner xsep=6.5pt,inner ysep=3pt]{%
    \popxboldls\fontsize{7}{7}\selectfont\color{#3}\MakeUppercase{#1}};%
}
\newcommand{\popsquiggle}[1]{%
  \tikz[baseline=-0.4ex]{
    \draw[#1,line width=2.6pt,line cap=round]
      plot [smooth] coordinates {(0,0.09) (0.34,0.01) (0.68,0.21) (1.02,0.01) (1.36,0.10)};
  }%
}
% Mot-clé surligné (marqueur orange sous le texte)
\newcommand{\cle}[1]{{\popxbold\color{popdark}#1}}

% ============================================================
% En-tête / pied
% ============================================================
\pagestyle{fancy}
\fancyhf{}
\renewcommand{\headrulewidth}{0pt}
\renewcommand{\footrulewidth}{0pt}
\lhead{\raisebox{-0.15ex}{\poplogo\fontsize{13}{13}\selectfont\color{popink}OnBuch\color{poporange}+}}
\rhead{\poppill{\DOCMATIERE\ \textbullet\ \DOCNIVEAU\ \textbullet\ Leçon \DOCLECON}{white}{popink}}
\lfoot{\popxbold\fontsize{7.6}{7.6}\selectfont\color{popmuted}\MakeUppercase{Cours signé OnBuch+}\ \textbullet\ {\popsemi\DOCTITRE}}
\rfoot{\tikz[baseline=-0.6ex]\node[rounded corners=8.5pt,fill=popink,minimum height=17pt,minimum width=17pt,inner xsep=5pt,inner sep=0]{\popxbold\fontsize{8}{8}\selectfont\color{popcream}\thepage\,/\,\pageref*{LastPage}};}

% ============================================================
% Titres
% ============================================================
\newcommand{\chaptitre}[2]{%
  \begin{tcolorbox}[enhanced,colback=poporange,colframe=popink,boxrule=2.6pt,arc=11pt,
      drop shadow=popink,left=15pt,right=15pt,top=12pt,bottom=11pt,before skip=3pt,after skip=18pt]
    \poppill{#2}{popink}{popcream}\par\vspace{9pt}
    {\raggedright\hyphenpenalty=10000\exhyphenpenalty=10000\popdisplay\fontsize{22}{24}\selectfont\color{popcream}\MakeUppercase{#1}\par}
  \end{tcolorbox}
}
% Section numérotée : pastille orange avec le numéro + titre Archivo
\newcounter{popsec}
\newcommand{\coursec}[1]{%
  \stepcounter{popsec}\setcounter{popsub}{0}%
  \par\vspace{16pt}\noindent
  \begin{minipage}{\linewidth}
  \tikz[baseline=-0.7ex]\node[draw=popink,line width=1.6pt,fill=poporange,rounded corners=4pt,minimum size=22pt,inner sep=0]{\popdisplay\fontsize{12}{12}\selectfont\color{popcream}\thepopsec};%
  \hspace{8pt}{\popdisplay\fontsize{15}{17}\selectfont\color{popink}\MakeUppercase{#1}}\par
  \vspace{4pt}\noindent{\color{poporange}\rule{36pt}{3.6pt}}
  \end{minipage}\par\nopagebreak\vspace{8pt}
}
\newcounter{popsub}
\newcommand{\courssub}[1]{%
  \stepcounter{popsub}%
  \par\vspace{9pt}\noindent{\popxbold\fontsize{11.5}{13}\selectfont\color{popink}{\color{poporange}\thepopsec.\thepopsub}\hspace{6pt}#1}\par\nopagebreak\vspace{4pt}
}

% ============================================================
% Boîtes pédagogiques — même recette : fond blanc, bordure épaisse colorée,
% ombre dure encre, badge plein. Titre optionnel [#1].
% ============================================================
\tcbset{popbase/.style={breakable,enhanced,colback=white,boxrule=2.3pt,arc=9pt,
  shadow={3pt}{-3pt}{0pt}{popink},left=14pt,right=14pt,top=11pt,bottom=11pt,before skip=11pt,after skip=15pt,
  fontupper=\popsemi\fontsize{10.6}{14.5}\selectfont\color{popink},
  fontlower=\popsemi\fontsize{10.6}{14.5}\selectfont\color{popink}}}
% Coche dessinée (le glyphe ✓ n'existe pas dans les polices du document)
\newcommand{\popcheck}[1]{\tikz[baseline=-0.5ex]\draw[#1,line width=1.6pt,line cap=round,line join=round](-0.13,0.0)--(-0.04,-0.09)--(0.13,0.1);}
\providecommand{\coche}{\popcheck{popgreen}}
\providecommand{\resultat}[1]{\par\smallskip\noindent\fcolorbox{popgreen}{popgreenL}{\parbox{\dimexpr\linewidth-2\fboxsep-2\fboxrule\relax}{\textbf{Résultat.} #1}}\par\smallskip}
\newcommand{\popboxhead}[3]{\poppill{#1}{#2}{white}\ifx\relax#3\relax\else\hspace{6pt}{\popxbold\fontsize{10.6}{12}\selectfont\color{popink}#3}\fi\par\vspace{7pt}}

\newtcolorbox{aretenir}[1][]{popbase,colframe=popgreen,before upper={\popboxhead{À retenir}{popgreen}{#1}}}
% Texte en chinois (document de lecture, dialogue, extrait) : cadre
% d'étude. Un titre optionnel (source, auteur) s'affiche dans l'en-tête.
\newtcolorbox{textebox}[1][]{popbase,colframe=popblue,before upper={\popboxhead{文本 · Texte}{popblue}{#1}},after upper={}}
% Marques ✓ / ✗ (forme correcte / fautive) souvent utilisées par les modèles
\providecommand{\cross}{\textcolor{poppink}{\ensuremath{\times}}}
\providecommand{\xmark}{\cross}\providecommand{\crossmark}{\cross}\providecommand{\cmark}{\ensuremath{\checkmark}}
% Ligne de réponse / justification à compléter (inventée par les modèles)
\providecommand{\justification}[1]{\par\noindent\textit{Justification~:} \dotfill\par}
% \textipa (package tipa non chargé) : la police ne couvre pas l'API -> simple passage
\providecommand{\textipa}[1]{#1}
% Soulignement (ulem non chargé) : \uline, \ul
\providecommand{\uline}[1]{\underline{#1}}\providecommand{\ul}[1]{\underline{#1}}
% \glossaire{\item ...} inventée par les modèles : liste de glossaire
\providecommand{\glossaire}[1]{\begin{itemize}#1\end{itemize}}
% \alert (beamer) inventée par les modèles : texte mis en évidence
\providecommand{\alert}[1]{\textbf{\textcolor{poppink}{#1}}}
% variantes ASCII d'ordinaux écrites en macro
\providecommand{\iere}{\textsuperscript{re}}\providecommand{\ieme}{\textsuperscript{e}}\providecommand{\ere}{\textsuperscript{re}}\providecommand{\eme}{\textsuperscript{e}}\providecommand{\er}{\textsuperscript{er}}\providecommand{\ie}{\textsuperscript{e}}
% Ordinaux français écrits en macro par les modèles : 1\ière, 3\ième
\newcommand{\ière}{\textsuperscript{re}}\newcommand{\ième}{\textsuperscript{e}}
% \exemple (inventée par les modèles) : étiquette « Exemple : »
\providecommand{\exemple}{\textbf{Exemple~:}\ }
% \square utilisable aussi hors mode math (cases à cocher)
\AtBeginDocument{\let\squareMath\square\renewcommand{\square}{\ensuremath{\squareMath}}}
% Mot ou phrase en chinois à l'intérieur d'un paragraphe français
\newcommand{\zh}[1]{\ifmmode\text{#1}\else{#1}\fi}
\let\eng\zh \let\esp\zh \let\all\zh % alias tolérés
% flèches d'intonation inventées par les modèles
\providecommand{\textsearrow}{}\renewcommand{\textsearrow}{\ensuremath{\searrow}}\providecommand{\textnearrow}{}\renewcommand{\textnearrow}{\ensuremath{\nearrow}}
% \fr{..} : mot français (inventé par les modèles)
\providecommand{\fr}[1]{\textit{#1}}
% zhbox : encadré de texte chinois inventé par les modèles
\newenvironment{zhbox}{\begin{quote}}{\end{quote}}
% \vs : « versus » inventé par les modèles
\providecommand{\vs}{\textit{vs}\ }
% \blank : case à compléter inventée par les modèles
\providecommand{\blank}[1][]{\underline{\hspace{1.6em}}}
\newtcolorbox{exemplebox}[1][]{popbase,colframe=poporange,before upper={\popboxhead{Exemple}{poporange}{#1}}}
\newtcolorbox{definition}[1][]{popbase,colframe=popblue,before upper={\popboxhead{Définition}{popblue}{#1}}}
\newtcolorbox{propriete}[1][]{popbase,colframe=poppurple,before upper={\popboxhead{Propriété}{poppurple}{#1}}}
\newtcolorbox{methode}[1][]{popbase,colframe=popgold,before upper={\popboxhead{Méthode}{popgold}{#1}}}
\newtcolorbox{attention}[1][]{popbase,colframe=poppink,before upper={\popboxhead{Attention}{poppink}{#1}}}
\newtcolorbox{experience}[1][]{popbase,colframe=popblue,colback=popblueL,before upper={\popboxhead{Expérience}{popblue}{#1}}}
% « Le savais-tu ? » : carte violette pleine (contexte, culture, Cameroun)
\newtcolorbox{savaistu}[1][]{popbase,colback=poppurple,colframe=popink,
  fontupper=\popsemi\fontsize{10.4}{14.2}\selectfont\color{white},
  before upper={\poppill{Le savais-tu\,?}{white}{poppurple}\ifx\relax#1\relax\else\hspace{6pt}{\popxbold\color{white}#1}\fi\par\vspace{7pt}}}
% Exercice résolu : énoncé en haut, \tcblower puis la solution sur fond crème
\newcounter{popexo}\newcounter{popexr}
\newtcolorbox{exoresolu}[1][]{popbase,colframe=poporange,colbacklower=popcream,
  segmentation style={popink,line width=1.2pt,dashed},
  before upper={\stepcounter{popexr}\popboxhead{Exercice résolu \thepopexr}{poporange}{#1}},
  before lower={\poppill{Solution}{popink}{popcream}\par\vspace{6pt}}}
% Exercice (non résolu en place) — la correction est dans la partie Corrigés
\newtcolorbox{exercice}[1][]{popbase,colframe=popink,
  before upper={\stepcounter{popexo}\popboxhead{Exercice \thepopexo}{popink}{#1}}}
% Corrigé
\newtcolorbox{corrige}[1][]{popbase,colframe=popgreen,colback=popgreenL,
  before upper={\popboxhead{Corrigé}{popgreen}{#1}}}
% Objectifs / prérequis / bilan
\newtcolorbox{objectifs}{popbase,colframe=popink,colback=poporangeL,
  before upper={\popboxhead{Objectifs de la leçon}{popink}{}}}
\newtcolorbox{prerequis}{popbase,colframe=popink,colback=popblueL,
  before upper={\popboxhead{Prérequis}{popblue}{}}}
\newtcolorbox{bilan}{popbase,colframe=popink,colback=white,boxrule=2.8pt,
  before upper={\popboxhead{Fiche bilan}{poporange}{L'essentiel à connaître pour l'examen}}}
% Figure encadrée avec légende
\newtcolorbox{popfigure}[1][]{enhanced,breakable=false,colback=white,colframe=popline,boxrule=1.4pt,arc=8pt,
  left=10pt,right=10pt,top=10pt,bottom=8pt,before skip=10pt,after skip=12pt,halign=center,
  lower separated=false,halign lower=center,
  fontlower=\popsemi\fontsize{9}{11.5}\selectfont\color{popmuted},
  #1}
\newcommand{\legende}[1]{\par\vspace{6pt}{\popsemi\fontsize{9}{11.5}\selectfont\color{popmuted}#1}}

% ============================================================
% Signature OnBuch+
% ============================================================
\newcommand{\poplogoinline}[1]{{\poplogo\fontsize{#1}{#1}\selectfont\color{popink}OnBuch\color{poporange}+}}
\newcommand{\signatureonbuch}{%
  \begin{tcolorbox}[enhanced,colback=popink,colframe=popink,boxrule=2.2pt,arc=10pt,
      drop shadow=poporange,left=14pt,right=14pt,top=10pt,bottom=10pt,before skip=0pt,after skip=0pt]
    \tikz[baseline=-0.7ex]\node[circle,fill=popgreen,draw=popcream,line width=1.3pt,minimum size=22pt,inner sep=0]{\popcheck{white}};%
    \hspace{9pt}\begin{minipage}[c]{0.62\linewidth}
      {\popxbold\fontsize{12}{13}\selectfont\color{popcream}Signé\ \textbullet\ {\poplogo OnBuch\color{poporange}+}}\par\vspace{2pt}
      {\raggedright\popsemi\fontsize{8}{10.5}\selectfont\color{popline}\MakeUppercase{Équipe pédagogique OnBuch+}\\\MakeUppercase{Conforme au programme officiel MINESEC}\par}
    \end{minipage}\hfill
    {\poplogo\fontsize{22}{22}\selectfont\color{popcream}OnBuch\color{poporange}+}
  \end{tcolorbox}%
}

% ============================================================
% Page de garde
% #1 titre · #2 matière · #3 niveau · #4 module · #5 n° leçon · #6 durée ·
% #7 description / objectifs (texte ou itemize)
% ============================================================
\newcommand{\pagegarde}[7]{%
  \thispagestyle{empty}
  \newgeometry{a4paper,margin=1.7cm,top=1.6cm,bottom=1.4cm}%
  \begin{tikzpicture}[remember picture,overlay]
    \fill[poporange!18] ([xshift=-3.2cm,yshift=-3.6cm]current page.north east) circle (4.6cm);
    \fill[poppurple!14] ([xshift=2.4cm,yshift=7.6cm]current page.south west) circle (3.8cm);
  \end{tikzpicture}%
  {\poplogo\fontsize{30}{30}\selectfont\color{popink}OnBuch\color{poporange}+}\hfill
  \raisebox{0.6ex}{\poppill{Cours signé \textbullet\ Premium}{popink}{popcream}}\par
  \vspace{1pt}\popsquiggle{poppurple}\par
  \vspace{8pt}
  \begin{tcolorbox}[enhanced,colback=poporange,colframe=popink,boxrule=2.8pt,arc=13pt,
      drop shadow=popink,left=18pt,right=18pt,top=12pt,bottom=13pt,after skip=10pt]
    \poppill{#2 \textbullet\ #3}{popink}{popcream}\hspace{5pt}\poppill{#4}{white}{popink}\par\vspace{8pt}
    {\popdisplay\fontsize{14}{14}\selectfont\color{popink}LEÇON #5}\par\vspace{4pt}
    {\raggedright\hyphenpenalty=10000\exhyphenpenalty=10000\popdisplay\fontsize{25}{27}\selectfont\color{popcream}\MakeUppercase{#1}\par}
    \vspace{8pt}
    {\popxbold\fontsize{9.5}{9.5}\selectfont\color{popink}\MakeUppercase{Durée conseillée : #6}}
  \end{tcolorbox}
  \begin{tcolorbox}[enhanced,colback=white,colframe=popink,boxrule=2.2pt,arc=10pt,
      drop shadow=popink,left=16pt,right=16pt,top=10pt,bottom=10pt,after skip=8pt]
    \poppill{Dans cette leçon}{poppurple}{white}\par\vspace{5pt}
    \popsemi\fontsize{9.3}{12.6}\selectfont\color{popink}#7
  \end{tcolorbox}
  \noindent\begin{minipage}[t]{0.487\linewidth}
    \begin{tcolorbox}[enhanced,colback=popgreen,colframe=popink,boxrule=2.2pt,arc=10pt,
        drop shadow=popink,left=11pt,right=11pt,top=10pt,bottom=10pt,height=3.1cm,valign=center]
      \tcbox[colback=white,colframe=popink,boxrule=1.3pt,arc=5pt,boxsep=0pt,nobeforeafter,
        left=3pt,right=3pt,top=3pt,bottom=3pt,valign=center]{\qrcode[height=1.9cm]{https://play.google.com/store/apps/details?id=com.luvvix.onbuch&hl=fr}}%
      \hspace{8pt}\begin{minipage}[c]{0.5\linewidth}
        {\popdisplay\fontsize{11}{12.5}\selectfont\color{white}\MakeUppercase{Télécharge OnBuch}}\par\vspace{4pt}
        {\raggedright\popsemi\fontsize{8.4}{11}\selectfont\color{white}Gratuit \textbullet\ Google Play\\Tuteur IA, annales, concours, résultats\par}
      \end{minipage}
    \end{tcolorbox}
  \end{minipage}\hfill
  \begin{minipage}[t]{0.487\linewidth}
    \begin{tcolorbox}[enhanced,colback=poppurple,colframe=popink,boxrule=2.2pt,arc=10pt,
        drop shadow=popink,left=11pt,right=11pt,top=10pt,bottom=10pt,height=3.1cm,valign=center]
      \tikz[baseline=-0.7ex]\node[circle,fill=white,draw=popink,line width=1.3pt,minimum size=1.6cm,inner sep=0]{\popdisplay\fontsize{24}{24}\selectfont\color{poppurple}+};%
      \hspace{8pt}\begin{minipage}[c]{0.6\linewidth}
        {\popdisplay\fontsize{11}{12.5}\selectfont\color{white}\MakeUppercase{Passe à OnBuch+}}\par\vspace{4pt}
        {\raggedright\popsemi\fontsize{8.4}{11}\selectfont\color{white}Tous les cours signés, Tuteur IA illimité, fiches PDF — onglet OnBuch+ de l'app\par}
      \end{minipage}
    \end{tcolorbox}
  \end{minipage}
  \vspace{8pt}
  \popcontenu
  \vfill
  \signatureonbuch
  \restoregeometry
  \newpage
}

% Rangée « Ce cours contient » (page de garde)
\newcommand{\poptile}[3]{% #1 couleur #2 gros texte #3 légende
  \begin{tcolorbox}[enhanced,colback=#1,colframe=popink,boxrule=2pt,arc=9pt,shadow={2.5pt}{-2.5pt}{0pt}{popink},
      left=6pt,right=6pt,top=9pt,bottom=9pt,halign=center,valign=center,height=1.75cm,width=\linewidth,nobeforeafter]
    {\popdisplay\fontsize{12.5}{13}\selectfont\color{white}\MakeUppercase{#2}}\par\vspace{3pt}
    {\centering\popsemi\fontsize{7.8}{9.5}\selectfont\color{white}#3\par}
  \end{tcolorbox}}
\newcommand{\popcontenu}{%
  \par\noindent{\popxboldls\fontsize{7.5}{7.5}\selectfont\color{popmuted}CE COURS SIGNÉ CONTIENT}\par\vspace{7pt}
  \noindent\begin{minipage}[t]{0.235\linewidth}\poptile{popblue}{Cours}{complet, illustré, pas à pas}\end{minipage}\hfill
  \begin{minipage}[t]{0.235\linewidth}\poptile{poporange}{Méthodes}{et exercices résolus}\end{minipage}\hfill
  \begin{minipage}[t]{0.235\linewidth}\poptile{poppink}{Exercices}{type examen + corrigés}\end{minipage}\hfill
  \begin{minipage}[t]{0.235\linewidth}\poptile{popgreen}{Fiche bilan}{l'essentiel à retenir}\end{minipage}\par}

% Sommaire
\newtcolorbox{sommaire}{enhanced,colback=white,colframe=popink,boxrule=2.2pt,arc=10pt,
  drop shadow=popink,left=18pt,right=18pt,top=13pt,bottom=13pt,before skip=6pt,after skip=20pt,
  before upper={\poppill{Sommaire}{poppurple}{white}\par\vspace{9pt}\popsemi\fontsize{10.6}{14.5}\selectfont\color{popink}}}

% Signature de fin de cours — poussée tout en bas de la dernière page
\newcommand{\findecours}{%
  \par\vfill\nopagebreak
  \begin{center}\popsquiggle{poporange}\end{center}\vspace{4pt}
  \signatureonbuch
}
\AtBeginDocument{\def\llbracket{\mathopen{[\mkern-3.5mu[}}\def\rrbracket{\mathclose{]\mkern-3.5mu]}}} % intervalles d'entiers (glyphe ⟦ absent de la police)
% Glyphes absents de Fira Math : redéfinis à partir de glyphes disponibles
\AtBeginDocument{%
  \def\setminus{\mathbin{\backslash}}\let\smallsetminus\setminus
  \def\vdots{\mathord{\vbox{\baselineskip3.2pt\lineskiplimit0pt\kern4pt\hbox{.}\hbox{.}\hbox{.}}}}%
  \def\ddots{\mathinner{\mkern1mu\raise6.5pt\hbox{.}\mkern2mu\raise3.6pt\hbox{.}\mkern2mu\raise0.7pt\hbox{.}\mkern1mu}}%
  \def\triangle{\mathord{\scalebox{0.9}{$\symup{\Delta}$}}}%
  \def\nabla{\mathord{\raisebox{\depth}{\scalebox{1}[-1]{$\symup{\Delta}$}}}}%
  \def\checkmark{\popcheck{popdark}}%
  \def\star{\mathbin{\ast}}\def\bigstar{\mathord{\ast}}%
  \def\diamond{\mathbin{\scalebox{0.6}{$\Diamond$}}}%
  \def\bigcup{\mathop{\mathchoice{\raisebox{-0.3em}{\scalebox{1.7}{$\cup$}}}{\scalebox{1.25}{$\cup$}}{\cup}{\cup}}\displaylimits}%
  \def\bigcap{\mathop{\mathchoice{\raisebox{-0.3em}{\scalebox{1.7}{$\cap$}}}{\scalebox{1.25}{$\cap$}}{\cap}{\cap}}\displaylimits}%
  \def\lesseqqgtr{\mathrel{\lessgtr}}\def\bigoplus{\mathop{\oplus}\displaylimits}%
}
\providecommand{\ssi}{si et seulement si} % abréviation fréquente
\AtBeginDocument{\providecommand{\wideparen}{\overparen}} % arc de cercle

%% FIGFIX-BEGIN v5 — correctifs globaux des figures (docs/onbuch-generation/tools/figfix)
\usepackage{adjustbox}\usepackage{etoolbox}\tcbuselibrary{raster}
% toute figure TikZ/pgfplots plus large que la ligne est réduite à la largeur disponible
\BeforeBeginEnvironment{tikzpicture}{\begin{adjustbox}{max width=\linewidth}}
\AfterEndEnvironment{tikzpicture}{\end{adjustbox}}
% légendes pgfplots lisibles par-dessus les courbes
\pgfplotsset{every axis/.append style={legend style={fill=white,fill opacity=0.92,text opacity=1,draw=black!15,inner sep=3pt}}}
% étiquettes d'arêtes (syntaxe edge["..."]) avec fond blanc
\tikzset{every edge quotes/.append style={fill=white,inner sep=1.2pt}}
%% FIGFIX-END
\begin{document}

\pagegarde{Leçon 6 — Grammaire : 如果……，那么……；les nombres ；有的……有的……；是……的}{Chinois}{1ère A}{Module 1 : Vie familiale et intégration sociale}{ZHA06}{2 h}{Cette leçon de grammaire présente quatre outils essentiels de la phrase chinoise : l'hypothèse avec 如果……那么……, la formation et l'emploi des nombres, la répartition avec 有的……有的…… et la mise en relief avec 是……的. Chaque structure est observée, schématisée, comparée au français et entraînée par des exercices d'application et de transformation.
\vspace{4pt}
\begin{multicols}{2}\small
\begin{itemize}
\item À la fin de cette leçon, je sais construire une phrase hypothétique avec 如果……，就/那么……
\item À la fin de cette leçon, je sais lire et écrire les nombres chinois de 0 à 10 000 et les utiliser dans des phrases
\item À la fin de cette leçon, je sais décrire une scène avec 有的……有的…… (certains… d'autres…)
\item À la fin de cette leçon, je sais mettre en relief le moment, le lieu ou la manière d'une action passée avec 是……的
\item À la fin de cette leçon, je sais placer correctement les mots-outils 就, 那么, 的 dans la phrase
\item À la fin de cette leçon, je sais comparer ces structures avec le français et éviter les erreurs typiques des francophones
\end{itemize}\end{multicols}}

\begin{sommaire}
\begin{multicols}{2}
\begin{enumerate}
\item Observation d'exemples en contexte
\item L'hypothèse : 如果……，那么/就……
\item Les nombres
\item La répartition : 有的……有的……
\item La mise en relief : 是……的
\item Synthèse, comparaison et entraînement
\item Activité d'intégration
\item Méthodes et exercices résolus
\item Exercices
\item Corrigés des exercices
\item Fiche bilan
\end{enumerate}
\end{multicols}
\end{sommaire}

\begin{objectifs}
\begin{itemize}
\item À la fin de cette leçon, je sais construire une phrase hypothétique avec 如果……，就/那么……
\item À la fin de cette leçon, je sais lire et écrire les nombres chinois de 0 à 10 000 et les utiliser dans des phrases
\item À la fin de cette leçon, je sais décrire une scène avec 有的……有的…… (certains… d'autres…)
\item À la fin de cette leçon, je sais mettre en relief le moment, le lieu ou la manière d'une action passée avec 是……的
\item À la fin de cette leçon, je sais placer correctement les mots-outils 就, 那么, 的 dans la phrase
\item À la fin de cette leçon, je sais comparer ces structures avec le français et éviter les erreurs typiques des francophones
\item À la fin de cette leçon, je sais transformer des phrases simples en phrases utilisant ces quatre structures
\item À la fin de cette leçon, je sais employer ces structures dans un court texte cohérent sur ma vie quotidienne
\end{itemize}
\end{objectifs}

\begin{prerequis}
\begin{itemize}
\item Structure de base de la phrase chinoise : Sujet + Verbe + Complément (Seconde)
\item Pronoms personnels et verbes courants (去, 来, 买, 学习, 看)
\item Marqueurs de temps et compléments de temps (今天, 昨天, 明天)
\item Nombres de 1 à 10 déjà rencontrés en Seconde
\item Négation avec 不 et 没(有)
\end{itemize}
\end{prerequis}

\newpage

\coursec{Observation d'exemples en contexte}

C'est la saison des pluies à Yaoundé. Le ciel se couvre, les gouttes commencent à tomber sur les toits de tôle, et Amina dit à son petit frère : « S'il pleut demain, alors je prendrai un taxi pour aller au marché Mokolo. » Plus tard, assise au parc, elle observe autour d'elle : « Certains enfants jouent au foot, d'autres mangent des beignets. » Le soir, de retour à la maison, elle montre son achat à sa mère : « C'est hier que j'ai acheté ce pagne. »

Trois phrases de la vie quotidienne camerounaise, trois structures chinoises indispensables : \zh{如果……，那么……}, \zh{有的……有的……} et \zh{是……的}. Ajoutons-y les nombres, qui servent à compter les élèves d'une classe, les beignets du marché ou les minutes d'un match. Comment dire tout cela en chinois ? C'est toute cette leçon.

\textbf{Problématique :} comment le chinois exprime-t-il l'hypothèse, la quantité, la répartition et la mise en relief, et où se placent exactement les mots-outils dans la phrase ?

\courssub{Lecture guidée : un texte pour tout observer}

Avant d'étudier les règles une par une, lisons un court texte qui les réunit toutes. Lis d'abord le chinois en silence, puis à voix haute, en suivant le pinyin.

\begin{textebox}[Au parc, un dimanche matin]
如果天气好，我们就去公园。公园里，有的人跳舞，有的人下棋。我们是早上八点到的。\\
\emph{Rúguǒ tiānqì hǎo, wǒmen jiù qù gōngyuán. Gōngyuán lǐ, yǒude rén tiàowǔ, yǒude rén xiàqí. Wǒmen shì zǎoshang bā diǎn dào de.}\\
« S'il fait beau, nous irons au parc. Au parc, certains dansent, d'autres jouent aux échecs. C'est à huit heures du matin que nous sommes arrivés. »
\end{textebox}

\textbf{Glossaire du texte :}

\begin{center}
\begin{tabularx}{\linewidth}{|X|X|X|X|}
\hline
\rowcolor{popblueL} Caractères & Pinyin & Nature & Traduction \\
\hline
天气 & tiānqì & nom & le temps (qu'il fait) \\
\hline
公园 & gōngyuán & nom & le parc \\
\hline
跳舞 & tiàowǔ & verbe & danser \\
\hline
下棋 & xiàqí & verbe & jouer aux échecs \\
\hline
早上 & zǎoshang & nom (temps) & le matin \\
\hline
到 & dào & verbe & arriver \\
\hline
\end{tabularx}
\end{center}

\textbf{Questions d'observation :}
\begin{enumerate}
\item Quels mots se répètent ou forment des paires dans ce texte ? Souligne-les.
\item Dans la première phrase, quel mot ouvre la proposition de condition ? Quel mot introduit la conséquence ?
\item Dans la dernière phrase, où se placent \zh{是} et \zh{的} par rapport à \zh{早上八点} ?
\end{enumerate}

\courssub{Les quatre phrases modèles}

Voici maintenant les quatre phrases modèles de la leçon, une par structure. Pour chacune, observe, souligne les mots-outils, puis essaie de deviner leur fonction avant de lire l'analyse.

\begin{exemplebox}[Quatre phrases, quatre structures]
\textbf{Exemple 1 — L'hypothèse :}\\
\zh{如果明天下雨，那么我就不去市场。}\\
\emph{Rúguǒ míngtiān xià yǔ, nàme wǒ jiù bù qù shìchǎng.}\\
« S'il pleut demain, je n'irai pas au marché. »\\[0.3em]
\textbf{Exemple 2 — Les nombres :}\\
\zh{我们班有三十五个学生。}\\
\emph{Wǒmen bān yǒu sānshíwǔ ge xuésheng.}\\
« Notre classe compte trente-cinq élèves. »\\[0.3em]
\textbf{Exemple 3 — La répartition :}\\
\zh{公园里，有的人跑步，有的人唱歌。}\\
\emph{Gōngyuán lǐ, yǒude rén pǎobù, yǒude rén chànggē.}\\
« Au parc, certains courent, d'autres chantent. »\\[0.3em]
\textbf{Exemple 4 — La mise en relief :}\\
\zh{我是昨天买的这本书。}\\
\emph{Wǒ shì zuótiān mǎi de zhè běn shū.}\\
« C'est hier que j'ai acheté ce livre. »
\end{exemplebox}

\courssub{Repérage des mots-outils}

Dans chaque phrase modèle, certains petits mots ne se traduisent pas toujours directement, mais ils \emph{construisent} la phrase : ce sont les \cle{mots-outils}. Repérons-les.

\begin{definition}[Mot-outil]
Un \cle{mot-outil} est un mot grammatical qui n'a pas de sens concret propre (ce n'est ni un objet, ni une action), mais qui indique la \textbf{relation} entre les parties de la phrase : condition, répartition, mise en relief. En chinois, les mots-outils ont une \textbf{position fixe} qu'il faut respecter absolument.
\end{definition}

\begin{center}
\begin{tabularx}{\linewidth}{|X|X|X|X|}
\hline
\rowcolor{popblueL} Structure & Mots-outils & Fonction & Équivalent français \\
\hline
\zh{如果……，那么/就……} & \zh{如果}, \zh{那么}, \zh{就} & exprimer une hypothèse et sa conséquence & « si…, alors… » \\
\hline
Les nombres & \zh{十}, \zh{百}, \zh{千}, \zh{个} & compter, quantifier & « trente-cinq élèves » \\
\hline
\zh{有的……有的……} & \zh{有的} (répété) & répartir un groupe en sous-groupes & « certains…, d'autres… » \\
\hline
\zh{是……的} & \zh{是} + \zh{的} & mettre en relief quand, où, comment & « c'est… que » \\
\hline
\end{tabularx}
\end{center}

\begin{methode}[Observer une structure grammaticale en trois étapes]
\begin{enumerate}
\item \textbf{Repérer les mots-outils} : souligne les petits mots qui se répètent ou qui encadrent une partie de la phrase (\zh{如果}, \zh{那么}, \zh{就}, \zh{有的}, \zh{是…的}).
\item \textbf{Noter leur position} : avant le sujet ou après ? Au début de la proposition ou à la fin ? En chinois, la position est fixe et ne change jamais.
\item \textbf{Comparer avec la traduction française} : qu'est-ce que le français exprime avec des mots, et qu'est-ce que le chinois exprime avec la position ? Par exemple, le français « c'est hier \emph{que} » correspond au cadre chinois \zh{是……的}.
\end{enumerate}
\end{methode}

Appliquons cette méthode à la phrase 1 : \zh{如果明天下雨，那么我就不去市场。}
\begin{itemize}
\item Étape 1 : les mots-outils sont \zh{如果} et \zh{那么} (avec \zh{就} juste après).
\item Étape 2 : \zh{如果} ouvre la proposition de condition ; \zh{那么} et \zh{就} se placent au début de la conséquence, \zh{就} se met juste avant le verbe.
\item Étape 3 : le français dit « si… alors… » ; le chinois dit \zh{如果……，那么……就……}. La correspondance est presque mot à mot, mais attention : \zh{就} n'a pas d'équivalent français, il marque simplement l'enchaînement logique.
\end{itemize}

\courssub{Carte mentale des quatre structures}

\begin{popfigure}
\begin{center}\tcbox[colback=popblue,colframe=popblue,coltext=white,arc=9pt,boxsep=4pt,left=6pt,right=6pt]{\bfseries\small Leçon 6 Quatre structures}\end{center}
\vspace{-2pt}
\begin{tcbraster}[raster columns=2,raster equal height=rows,raster column skip=6pt,raster row skip=6pt,enhanced,blankest]
\begin{tcolorbox}[enhanced,colback=poporange,colframe=poporange,coltext=white,boxrule=0.9pt,arc=3pt,left=4pt,right=4pt,top=3pt,bottom=3pt,boxsep=1pt,halign=left]\bfseries\scriptsize 如果…那么… \footnotesize 如果明天下雨， 那么我就不去市场。\end{tcolorbox}
\begin{tcolorbox}[enhanced,colback=popgreen,colframe=popgreen,coltext=white,boxrule=0.9pt,arc=3pt,left=4pt,right=4pt,top=3pt,bottom=3pt,boxsep=1pt,halign=left]\bfseries\scriptsize Les nombres \footnotesize 我们班有 三十五个学生。\end{tcolorbox}
\begin{tcolorbox}[enhanced,colback=poppurple,colframe=poppurple,coltext=white,boxrule=0.9pt,arc=3pt,left=4pt,right=4pt,top=3pt,bottom=3pt,boxsep=1pt,halign=left]\bfseries\scriptsize 有的…有的… \footnotesize 有的人跑步， 有的人唱歌。\end{tcolorbox}
\begin{tcolorbox}[enhanced,colback=poppink,colframe=poppink,coltext=white,boxrule=0.9pt,arc=3pt,left=4pt,right=4pt,top=3pt,bottom=3pt,boxsep=1pt,halign=left]\bfseries\scriptsize 是……的 \footnotesize 我是昨天 买的这本书。\end{tcolorbox}
\end{tcbraster}

\legende{Carte mentale des quatre structures de la leçon, avec un exemple sous chaque branche.}
\end{popfigure}

\courssub{Premières hypothèses : à toi de deviner la règle}

Avant la leçon magistrale, formule tes propres hypothèses. Voici celles qu'expriment souvent les élèves de Première :

\begin{itemize}
\item « \zh{如果} veut dire "si", donc je peux peut-être dire la condition sans \zh{那么}. » — \textbf{Bonne intuition !} En effet, \zh{那么} est souvent omis, mais \zh{就} reste presque toujours devant le verbe de la conséquence : \zh{如果明天下雨，我就不去市场。}
\item « \zh{有的} ressemble à \zh{有} ("il y a"), donc \zh{有的人} doit vouloir dire "il y a des gens qui…". » — \textbf{Exact}, et c'est précisément l'origine de la structure : « il y a des gens qui courent, il y a des gens qui chantent » devient « certains courent, d'autres chantent ».
\item « \zh{是……的} contient \zh{是} ("être"), donc la phrase doit être au présent. » — \textbf{Piège !} \zh{是……的} sert au contraire à parler d'une action \emph{déjà accomplie}, en insistant sur le moment, le lieu ou la manière.
\end{itemize}

\begin{attention}[Ne traduis jamais mot à mot depuis le français]
Le réflexe du francophone est de construire la phrase chinoise en calquant le français. Deux dangers majeurs :
\begin{itemize}
\item En français, « si » peut se placer au milieu de la phrase (« Je n'irai pas au marché \emph{s'il pleut} »). En chinois, la proposition avec \zh{如果} vient \textbf{toujours en premier} : jamais \zh{我就不去市场，如果明天下雨}.
\item Le français dit « c'est hier \emph{que} j'ai acheté » avec « que » au milieu ; le chinois encadre l'information avec \zh{是} au début et \zh{的} à la fin : \zh{我是昨天买的}. Oublier le \zh{的} final est l'erreur la plus fréquente des élèves camerounais.
\end{itemize}
Retiens : les mots-outils chinois ont des \textbf{positions fixes}. On ne les déplace pas, on ne les oublie pas.
\end{attention}

\begin{exoresolu}[Repérage guidé]
\textbf{Énoncé.} Dans les phrases suivantes, souligne les mots-outils et indique à quelle structure (1 à 4) appartient chaque phrase :\\
a) \zh{如果我有时间，我就去看你。}\\
b) \zh{食堂里，有的同学吃米饭，有的同学吃面包。}\\
c) \zh{他家有五口人。}\\
d) \zh{她是在杜阿拉出生的。}
\tcblower
\textbf{Solution détaillée.}\\
a) Mots-outils : \zh{如果} et \zh{就}. Structure 1 (hypothèse) : « Si j'ai le temps, j'irai te voir. » On remarque que \zh{那么} est omis, ce qui est très fréquent.\\
b) Mots-outils : \zh{有的} répété deux fois. Structure 3 (répartition) : « À la cantine, certains élèves mangent du riz, d'autres mangent du pain. »\\
c) Mot-outil : le nombre \zh{五} suivi du classificateur \zh{口} (pour les membres d'une famille). Structure 2 (les nombres) : « Sa famille compte cinq personnes. »\\
d) Mots-outils : \zh{是} et \zh{的} qui encadrent \zh{在杜阿拉出生}. Structure 4 (mise en relief) : « C'est à Douala qu'elle est née. » On insiste sur le lieu de naissance.
\end{exoresolu}

\begin{exercice}[À ton tour]
Observe le texte du parc (\emph{textebox} ci-dessus) et réponds par écrit :
\begin{enumerate}
\item Recopie la phrase d'hypothèse et souligne ses deux mots-outils.
\item Trouve les deux verbes d'action dans la phrase de répartition. Que font « certains » ? Que font « les autres » ?
\item Dans \zh{我们是早上八点到的}, quelle information est mise en relief : le moment, le lieu ou la manière ?
\item Reformule en français la phrase \zh{我们是早上八点到的} sans utiliser « c'est… que ». Laquelle des deux versions insiste le plus sur l'heure ?
\end{enumerate}
\end{exercice}

\begin{corrige}[Exercice « À ton tour »]
\begin{enumerate}
\item \zh{如果天气好，我们就去公园。} — mots-outils : \zh{如果} et \zh{就}.
\item Les verbes sont \zh{跳舞} (danser) et \zh{下棋} (jouer aux échecs). Certains dansent, les autres jouent aux échecs.
\item C'est le \textbf{moment} qui est mis en relief : \zh{早上八点} (huit heures du matin).
\item « Nous sommes arrivés à huit heures du matin. » La version avec « c'est… que » insiste davantage sur l'heure, exactement comme la structure chinoise \zh{是……的}.
\end{enumerate}
\end{corrige}

\begin{aretenir}[L'essentiel de cette première étape]
\begin{itemize}
\item Quatre structures à maîtriser dans cette leçon : \zh{如果……，那么/就……} (hypothèse), les nombres (quantité), \zh{有的……有的……} (répartition), \zh{是……的} (mise en relief d'une action passée).
\item Chaque structure repose sur des \cle{mots-outils} placés à des \textbf{positions fixes}.
\item La méthode d'observation : repérer les mots-outils, noter leur position, comparer avec le français.
\item Ne jamais calquer l'ordre des mots français : la condition avec \zh{如果} vient toujours en premier, et le cadre \zh{是……的} exige les deux éléments.
\end{itemize}
\end{aretenir}

\coursec{L'hypothèse : 如果……，那么/就……}

Dans la section précédente, nous avons observé plusieurs phrases en contexte. Vous avez remarqué que certaines d'entre elles exprimaient une \emph{condition} : « si cela arrive, alors ceci se produira ». C'est exactement le rôle de la structure \zh{如果……，那么/就……}, la toute première structure hypothétique du programme de Première. Nous allons maintenant l'étudier méthodiquement : observation, règle, comparaison avec le français, exemples, puis entraînement.

\courssub{Observation : lire avant de comprendre la règle}

Lisons d'abord un court dialogue entre deux élèves d'un lycée de Yaoundé, la veille d'une sortie scolaire.

\begin{textebox}[Dialogue : la sortie à Douala]
A : 明天我们去杜阿拉，你去吗？\\
\emph{Míngtiān wǒmen qù Dù'ālā, nǐ qù ma ?}\\
« Demain nous allons à Douala, tu viens ? »

B : 如果明天不下雨，我就去。\\
\emph{Rúguǒ míngtiān bù xià yǔ, wǒ jiù qù.}\\
« S'il ne pleut pas demain, j'irai. »

A : 如果下雨呢？\\
\emph{Rúguǒ xià yǔ ne ?}\\
« Et s'il pleut ? »

B : 如果下雨，那么我们就在学校上课。\\
\emph{Rúguǒ xià yǔ, nàme wǒmen jiù zài xuéxiào shàng kè.}\\
« S'il pleut, alors nous resterons en classe à l'école. »

A : 好！如果你去，我就很高兴。\\
\emph{Hǎo ! Rúguǒ nǐ qù, wǒ jiù hěn gāoxìng.}\\
« Bien ! Si tu viens, je serai très content. »
\end{textebox}

\begin{center}
\begin{tabularx}{\linewidth}{|X|X|X|X|}
\hline
\rowcolor{popblueL} Caractères & Pinyin & Nature & Traduction \\
\hline
如果 & rúguǒ & conjonction & si \\
\hline
那么 & nàme & conjonction & alors, dans ce cas \\
\hline
就 & jiù & adverbe & alors, aussitôt (conséquence) \\
\hline
下雨 & xià yǔ & verbe & pleuvoir \\
\hline
上课 & shàng kè & verbe & suivre les cours, être en classe \\
\hline
高兴 & gāoxìng & adjectif & content, joyeux \\
\hline
\end{tabularx}
\end{center}

\textbf{Questions d'observation :}
\begin{enumerate}
\item Quel mot introduit la condition dans chaque phrase ? Réponse : \zh{如果}.
\item Quels mots introduisent la conséquence ? Réponse : \zh{就} ou \zh{那么}.
\item Où se placent \zh{就} et \zh{那么} par rapport au sujet et au verbe de la conséquence ?
\item La condition vient-elle avant ou après la conséquence ?
\end{enumerate}

\courssub{La règle : forme et emploi}

\begin{propriete}[Structure de l'hypothèse]
La structure hypothétique chinoise suit le schéma :

\textbf{如果 + condition ，（那么 / 就）+ conséquence}

\begin{itemize}
\item \zh{如果} (rúguǒ) introduit la \cle{condition} et se place en tête de la première proposition (avant ou juste après son sujet).
\item \zh{就} (jiù) est un \textbf{adverbe} : il introduit la conséquence et se place \textbf{juste avant le verbe} (ou l'adjectif) de la seconde proposition, c'est-à-dire \textbf{après le sujet}.
\item \zh{那么} (nàme) est une \textbf{conjonction} : il peut se placer \textbf{en tête de la proposition de conséquence} (suivi du sujet) \textbf{ou} devant le verbe (souvent associé à \zh{就} : \zh{那么……就……}).
\item La condition précède \textbf{toujours} la conséquence : l'ordre des propositions est fixe en chinois.
\item Les deux propositions sont séparées par une virgule chinoise \zh{，}.
\end{itemize}
\end{propriete}

\begin{popfigure}
\begin{tikzpicture}[node distance=0.6cm and 1.6cm]
\node[draw=popblue, fill=popblueL, rounded corners, thick, align=center, minimum width=4.2cm, minimum height=1.2cm] (cond) {如果 + condition\\ \emph{si\dots}\\ \zh{如果明天下雨}};
\node[draw=poporange, fill=poporangeL, rounded corners, thick, align=center, minimum width=4.2cm, minimum height=1.2cm, right=of cond] (cons) {(那么 / 就) + conséquence\\ \emph{alors\dots}\\ \zh{我们就不去 / 那么我们就不去}};
\draw[-{Stealth[length=3mm]}, very thick, popdark] (cond) -- (cons) node[midway, above, align=center, font=\small] {condition $\rightarrow$ conséquence};
\node[below=0.35cm of cond, align=center, font=\small, text=popmuted] {proposition 1 : l'hypothèse};
\node[below=0.35cm of cons, align=center, font=\small, text=popmuted] {proposition 2 : le résultat};
\end{tikzpicture}
\legende{Schéma de la structure 如果……，那么/就…… : la condition ouvre toujours la phrase, la conséquence suit.}
\end{popfigure}

\courssub{就 ou 那么 : lequel choisir ?}

Les deux mots sont corrects, mais ils ne s'emploient pas exactement de la même façon.

\begin{definition}[Différence entre 就 et 那么]
\begin{itemize}
\item \zh{就} (jiù) : \textbf{adverbe}, langue \textbf{parlée} courante. Marque une conséquence naturelle, immédiate, automatique : « si A, alors (aussitôt) B ». Place : \textbf{Sujet + 就 + Verbe}.
\item \zh{那么} (nàme) : \textbf{conjonction}, plus \textbf{soutenu/écrit}. Marque la \textbf{conséquence logique}, le raisonnement (« dans ce cas, donc »). Place : \textbf{那么 + Sujet + Verbe} \emph{ou} \textbf{Sujet + 那么 + 就 + Verbe}.
\item Combinaison fréquente : \zh{那么……就……} renforce la logique de l'ensemble (écrit formel).
\end{itemize}
\end{definition}

\begin{exemplebox}[Exemples commentés]
\zh{如果你有时间，就来我家。}\\
\emph{Rúguǒ nǐ yǒu shíjiān, jiù lái wǒ jiā.}\\
« Si tu as le temps, viens chez moi. » (parlé, conséquence immédiate, sujet omis à l'impératif)

\zh{如果他不来，那么我们就开始。}\\
\emph{Rúguǒ tā bù lái, nàme wǒmen jiù kāishǐ.}\\
« S'il ne vient pas, alors nous commencerons. » (\zh{那么} en tête, \zh{就} devant le verbe)

\zh{如果太贵，我就不买。}\\
\emph{Rúguǒ tài guì, wǒ jiù bù mǎi.}\\
« Si c'est trop cher, je n'achète pas. » (au marché de Mokolo, phrase très utile !)

\zh{如果你学汉语，那么你可以去中国。}\\
\emph{Rúguǒ nǐ xué Hànyǔ, nàme nǐ kěyǐ qù Zhōngguó.}\\
« Si tu apprends le chinois, tu peux aller en Chine. » (\zh{那么} introduit la conclusion logique)
\end{exemplebox}

\courssub{La négation dans la condition}

Pour exprimer « si\dots ne\dots pas », on place simplement la négation \zh{不} (ou \zh{没} pour le passé / \zh{没有}) \textbf{devant le verbe de la condition}, à l'intérieur de la proposition introduite par \zh{如果}.

\begin{propriete}[Négation de la condition]
Schéma : \textbf{如果 + sujet + 不 / 没(有) + verbe ，（那么 / 就）+ conséquence}

\zh{如果明天不下雨，我们就去杜阿拉。}\\
\emph{Rúguǒ míngtiān bù xià yǔ, wǒmen jiù qù Dù'ālā.}\\
« S'il ne pleut pas demain, nous irons à Douala. »

\zh{如果你没有钱，就别买这个东西。}\\
\emph{Rúguǒ nǐ méiyǒu qián, jiù bié mǎi zhège dōngxi.}\\
« Si tu n'as pas d'argent, n'achète pas cette chose. » (\zh{别} = impératif négatif « ne ... pas »)
\end{propriete}

\courssub{Variante : l'omission de 如果}

\begin{aretenir}[如果 peut être omis à l'oral]
Dans la \textbf{conversation courante}, quand le contexte est clair, on peut omettre \zh{如果} ; la conséquence avec \zh{就} suffit à signaler l'hypothèse :

\zh{明天下雨，我们就不去。}\\
\emph{Míngtiān xià yǔ, wǒmen jiù bù qù.}\\
« (Si) demain il pleut, nous n'irons pas. »

\textbf{Attention pour l'examen :} dans une production écrite ou une traduction, il est fortement conseillé d'écrire \zh{如果} en entier. L'omission est une facilité orale, pas une règle de grammaire à exploiter au baccalauréat.
\end{aretenir}

\courssub{Comparaison avec le français}

\begin{center}
\begin{tabularx}{\linewidth}{|X|X|X|}
\hline
\rowcolor{popblueL} Point comparé & Français & Chinois \\
\hline
Mot de condition & « si » & \zh{如果} (rúguǒ) \\
\hline
Mot de conséquence & « alors » (souvent omis) & \zh{就} / \zh{那么} (très fréquent, presque obligatoire à l'oral) \\
\hline
Conjugaison & « si + imparfait \dots conditionnel » : \emph{s'il venait, je partirais} & aucune conjugaison : le verbe ne change jamais de forme \\
\hline
Marque de l'hypothèse & les temps verbaux (imparfait, conditionnel) & uniquement les mots-outils \zh{如果} et \zh{就/那么} \\
\hline
Ordre des propositions & libre : « je partirais s'il venait » & fixe : la condition \textbf{toujours avant} la conséquence \\
\hline
\end{tabularx}
\end{center}

\begin{aretenir}[L'idée essentielle]
Le chinois n'a \textbf{pas de conjugaison}. Le sens hypothétique ne repose sur aucun temps verbal : il repose entièrement sur le mot \zh{如果}. Que l'hypothèse porte sur le présent, le futur ou une situation irréelle, la structure reste identique : seuls les mots de temps (\zh{明天}, \zh{以后}\dots) précisent le moment.
\end{aretenir}

\courssub{Erreurs fréquentes des francophones}

\begin{attention}[Ne placez jamais 就 en tête de phrase]
L'erreur n° 1 des élèves francophones est de calquer le français « alors, je\dots » en mettant \zh{就} au début de la conséquence :

\zh{*如果下雨，就我不去。} \quad FAUX

\zh{就} est un adverbe : il suit le sujet et précède le verbe. La phrase correcte est :

\zh{如果下雨，我就不去。}\\
\emph{Rúguǒ xià yǔ, wǒ jiù bù qù.}\\
« S'il pleut, je n'irai pas. »

Autres pièges :
\begin{itemize}
\item Inverser l'ordre : \zh{*我不去，如果下雨} est impossible ; la condition vient d'abord.
\item Oublier la virgule \zh{，} entre les deux propositions à l'écrit.
\item Traduire « si » par \zh{如果} dans une question indirecte (« je ne sais pas s'il viendra ») : le chinois utilise alors la forme affirmative-négative \zh{他来不来} ou \zh{他是否来}. Ce point sera approfondi en Terminale.
\end{itemize}
\end{attention}

\courssub{Point culturel : les sorties scolaires}

\begin{savaistu}[Sorties scolaires : Yaoundé -- Pékin]
Au Cameroun, la « sortie scolaire » (\zh{春游} / \zh{秋游} \emph{chūnyóu / qiūyóu} : sortie de printemps / d'automne) est un moment fort : location de cars, pique-nique (\zh{野餐} \emph{yěcān}) au bord de la Sanaga ou visite du musée national. En Chine, les \zh{春游} / \zh{秋游} sont institutionnalisés : toute l'école part souvent le même jour vers un parc (\zh{公园} \emph{gōngyuán}) ou une montagne célèbre (\zh{名山} \emph{míngshān}), avec des activités encadrées (chants, jeux, \zh{拔河} \emph{báhó} : tir à la corde). La météo (\zh{天气} \emph{tiānqì}) dicte tout : \zh{如果下雨，那么取消} (« s'il pleut, c'est annulé ») est la phrase la plus entendue la veille du départ !
\end{savaistu}

\courssub{Mise en pratique orale}

\begin{experience}[Jeu de rôle : « On organise le week-end »]
\textbf{Consigne :} Par binômes. L'élève A propose une activité (match de foot, cinéma, marché, révision). L'élève B pose une condition météo ou matérielle avec \zh{如果……就/那么……}. Échangez les rôles.
\textbf{Exemple :}
A : \zh{这个周末我们去踢球，你去吗？}
B : \zh{如果不下雨，我就去。如果下雨，那么我在家看书。}
\textbf{Objectif :} Automatiser l'ordre \zh{如果} (condition) \textrightarrow \zh{就/那么} (conséquence) et la place de \zh{就} après le sujet.
\end{experience}

\courssub{Exercice résolu et entraînement}

\begin{exoresolu}[Traduction et construction]
\textbf{Énoncé.} Traduisez en chinois :
1) « Si tu es fatigué, alors repose-toi. »
2) « Si le professeur ne vient pas, nous rentrerons à la maison. »
3) « Si c'est trop loin, je n'y vais pas. »

\tcblower
\textbf{Solution détaillée.}

1) Condition : « tu es fatigué » = \zh{你累} ; conséquence (impératif) : « repose-toi » = \zh{休息}. Avec \zh{juste} devant le verbe (sujet omis) :

\zh{如果你累，就休息。}\\
\emph{Rúguǒ nǐ lèi, jiù xiūxi.}

2) Négation avec \zh{不} devant \zh{来} ; conséquence avec \zh{就} après le sujet \zh{我们} :

\zh{如果老师不来，我们就回家。}\\
\emph{Rúguǒ lǎoshī bù lái, wǒmen jiù huí jiā.}

3) « trop loin » = \zh{太远} ; sujet \zh{我} + \zh{就} + \zh{不去} :

\zh{如果太远，我就不去。}\\
\emph{Rúguǒ tài yuǎn, wǒ jiù bù qù.}

\textbf{Méthode :} repérer la condition et la conséquence en français, placer la condition d'abord, ajouter \zh{如果}, puis insérer \zh{就} (parlé) ou \zh{那么} (écrit/logique) juste avant le verbe de la conséquence (après le sujet).
\end{exoresolu}

\begin{exercice}[À vous de jouer]
Traduisez en chinois avec \zh{如果……，就/那么……} :
\begin{enumerate}
\item Si tu étudies bien, tu réussiras à l'examen. (\zh{通过考试} = réussir l'examen / \zh{考好} = bien composer)
\item S'il pleut demain, le match de football n'aura pas lieu. (\zh{比赛} = match ; \zh{取消} = annuler / \zh{不举行} = ne pas avoir lieu)
\item Si ce ndolé est trop épicé, je ne le mange pas. (\zh{辣} = épicé ; ajouter un classificateur \zh{份} ou \zh{碗})
\item Si nous avons de l'argent, nous irons en Chine.
\end{enumerate}
\end{exercice}

\begin{corrige}[Exercice « À vous de jouer »]
1) \zh{如果你好好学习，你就能通过考试。} \emph{Rúguǒ nǐ hǎohǎo xuéxí, nǐ jiù néng tōngguò kǎoshì.} — \zh{就} suit le sujet \zh{你}, \zh{能} exprime la possibilité future.

2) \zh{如果明天下雨，足球比赛就取消了。} \emph{Rúguǒ míngtiān xià yǔ, zúqiú bǐsài jiù qǔxiāo le.} — \zh{取消了} (annulé / \zh{了} marque le changement d'état) est plus précis que \zh{没有了}.

3) \zh{如果这份 ndolé 太辣，我就不吃。} \emph{Rúguǒ zhè fèn ndolé tài là, wǒ jiù bù chī.} — Classificateur \zh{份} (portion) pour un plat.

4) \zh{如果我们有钱，那么我们就去中国。} \emph{Rúguǒ wǒmen yǒu qián, nàme wǒmen jiù qù Zhōngguó.} — \zh{那么} introduit la conclusion logique, \zh{就} devant le verbe \zh{去}.
\end{corrige}

Maintenant que l'hypothèse est bien maîtrisée, passons à un outil indispensable de la langue quotidienne : les nombres, qui nous permettront de compter, de donner des prix et des quantités.

\coursec{Leçon 6 — Grammaire : 如果……，那么……；les nombres ；有的……有的……；是……的}

\courssub{Observation d'exemples en contexte}

\begin{textebox}[Dialogue : Préparatifs du voyage — 旅行准备]
\zh{小马：明天如果不下雨，那么我们就去爬摄影山。}\\
\emph{Xiǎomǎ : Míngtiān rúguǒ bù xiàyǔ, nàme wǒmen jiù qù pá Shèyǐngshān.}\\
« Petit Ma : Demain, s'il ne pleut pas, nous irons gravir le Mont Cameroun. »\\
\zh{阿华：好啊。有的同学想坐车，有的想走路。}\\
\emph{Āhuá : Hǎo a. Yǒude tóngxué xiǎng zuò chē, yǒude xiǎng zǒulù.}\\
« Ah Hua : D'accord. Certains camarades veulent prendre la voiture, d'autres veulent marcher. »\\
\zh{小马：我们是坐大巴去的，不是坐火车。}\\
\emph{Xiǎomǎ : Wǒmen shì zuò dàbā qù de, bùshì zuò huǒchē.}\\
« Petit Ma : C'est en bus que nous y allons, pas en train. »\\
\zh{阿华：对了，入场费多少钱？}\\
\emph{Āhuá : Duì le, rùchǎngfèi duōshao qián ?}\\
« Ah Hua : Au fait, combien coûte l'entrée ? »\\
\zh{小马：成人票一万五千法郎，学生票八千法郎。}\\
\emph{Xiǎomǎ : Chéngrén piào yī wàn wǔqiān fǎláng, xuésheng piào bāqiān fǎláng.}\\
« Petit Ma : Billet adulte 15\,000 francs, billet étudiant 8\,000 francs. »
\end{textebox}

\begin{aretenir}[Repérage des quatre structures de la leçon]
Dans ce dialogue authentique (contexte camerounais : Mont Cameroun, francs CFA), observez :
\begin{enumerate}
\item \textbf{Hypothèse} : \zh{如果不下雨，那么我们就去……} (Si condition, alors conséquence).
\item \textbf{Nombres} : \zh{一万五千} (15\,000), \zh{八千} (8\,000) — lecture, prix, classifieurs implicites.
\item \textbf{Répartition} : \zh{有的……有的……} (Certains…, d'autres…).
\item \textbf{Mise en relief} : \zh{是……的} (C'est \textbf{en bus} que nous y allons — insistance sur le moyen de transport).
\end{enumerate}
\end{aretenir}

\coursec{L'hypothèse : 如果……，那么/就……}

\courssub{Structure de base et ordre des mots}

\begin{propriete}[Structure conditionnelle réelle ou future]
\zh{如果} (rúguǒ, « si ») introduit la condition ; \zh{那么} (nàme, « alors ») ou \zh{就} (jiù, « alors / dès lors ») introduit la conséquence.
\medskip
\textbf{Schéma :} \zh{如果} + \textbf{Condition (Sujet + Verbe/Adj.)} ， \zh{那么/就} + \textbf{Conséquence (Sujet + Verbe)}.
\medskip
Le sujet de la conséquence peut être omis s'il est identique à celui de la condition.
\end{propriete}

\begin{exemplebox}[如果……，那么/就……]
\zh{如果你累了，就早点休息。}\\
\emph{Rúguǒ nǐ lèi le, jiù zǎodiǎn xiūxi.}\\
« Si tu es fatigué, repose-toi tôt. » (Sujet \zh{你} répété)\\
\zh{如果明天下雨，那么比赛取消。}\\
\emph{Rúguǒ míngtiān xiàyǔ, nàme bǐsài qǔxiāo.}\\
« S'il pleut demain, le match est annulé. » (Sujet différent : \zh{比赛})\\
\zh{如果想去中国，就要办签证。}\\
\emph{Rúguǒ xiǎng qù Zhōngguó, jiù yào bàn qiānzhèng.}\\
« Si (on) veut aller en Chine, (il) faut faire un visa. » (Sujet omis, générique)
\end{exemplebox}

\courssub{\zh{那么} vs \zh{就} : nuances}

\begin{center}
\begin{tabularx}{\linewidth}{|X|X|X|}
\hline
\rowcolor{popblueL} Marqueur & Nuance & Exemple \\
\hline
\zh{那么} (nàme) & Conséquence logique, narration, ton plus posé, souvent à l'écrit ou discours formel. & \zh{如果他不来，那么我们开始吧。} \\
\hline
\zh{就} (jiù) & Conséquence immédiate, inévitable, ton oral, insistance sur la rapidité ou la certitude. & \zh{如果你去，我就去。} \\
\hline
\end{tabularx}
\end{center}

\begin{attention}[Erreurs fréquentes des francophones]
\begin{itemize}
\item \textbf{Ordre des mots :} Ne mettez \textbf{jamais} \zh{就} ou \zh{那么} \emph{avant} le sujet de la conséquence. *\zh{如果下雨，我就不去} (correct) vs *\zh{如果下雨，就我不去} (faux).
\item \textbf{\zh{如果} vs \zh{要是}} : \zh{要是} (yàoshi) est plus oral, familier. Au bac, privilégiez \zh{如果}.
\item \textbf{Négation :} La négation se place \textbf{devant} le verbe dans la proposition qui convient. \zh{如果不下雨，我们就去。} / \zh{如果下雨，我们就不去。}
\item \textbf{Pas de « alors » redondant :} En français « Si..., alors... ». En chinois \zh{如果} suffit à marquer l'hypothèse ; \zh{那么/就} marque la conséquence. Ne rajoutez pas de mot supplémentaire.
\end{itemize}
\end{attention}

\courssub{Variantes : \zh{要是}…, \zh{除非}…}

\begin{definition}[Variantes conditionnelles]
\begin{itemize}
\item \zh{要是} (yàoshi) = \zh{如果} (oral) : \zh{要是你有空，来我家玩。}
\item \zh{除非} (chúfēi) = « à moins que » (condition nécessaire) : \zh{除非下大雨，否则我们照常出发。} (\zh{否则} = sinon).
\end{itemize}
\end{definition}

\coursec{Les nombres}

\courssub{Observation : au marché de Mokolo}

\begin{textebox}[Au marché — Zài shìchǎng]
\zh{顾客：老板，这个多少钱？}\\
\emph{Gùkè : Lǎobǎn, zhège duōshao qián ?}\\
« Cliente : Patron, combien ça coûte ? »\\
\zh{老板：三百五十法郎。}\\
\emph{Lǎobǎn : Sānbǎi wǔshí fǎláng.}\\
« Patron : 350 francs. »\\
\zh{顾客：太贵了！两百块，怎么样？}\\
\emph{Gùkè : Tài guì le ! Liǎng bǎi kuài, zěnmeyàng ?}\\
« Cliente : Trop cher ! 200 francs, ça va ? »\\
\zh{老板：好吧。你要几个？}\\
\emph{Lǎobǎn : Hǎo ba. Nǐ yào jǐ ge ?}\\
« Patron : D'accord. Tu en veux combien ? »\\
\zh{顾客：我要三个。}\\
\emph{Gùkè : Wǒ yào sān ge.}\\
« Cliente : J'en veux trois. »
\end{textebox}

Observons : \zh{三百五十} (350), \zh{两百} (200), \zh{三个} (trois + classifieur). Remarquez deux choses : les nombres se construisent par simple juxtaposition de chiffres, et devant un nom il faut un \cle{classifieur} (\zh{个}). Remarquez aussi que « deux cents » se dit \zh{两百} avec \zh{两} et non \zh{二}. Nous allons expliquer tout cela.

\courssub{Révision : les chiffres de 1 à 10 et unités de base}

\begin{center}
\begin{tabularx}{\linewidth}{|X|X|X|X|}
\hline
\rowcolor{popblueL} Caractère & Pinyin & Nature & Traduction \\
\hline
一 & yī & nombre & un \\
\hline
二 & èr & nombre & deux \\
\hline
三 & sān & nombre & trois \\
\hline
四 & sì & nombre & quatre \\
\hline
五 & wǔ & nombre & cinq \\
\hline
六 & liù & nombre & six \\
\hline
七 & qī & nombre & sept \\
\hline
八 & bā & nombre & huit \\
\hline
九 & jiǔ & nombre & neuf \\
\hline
十 & shí & nombre & dix \\
\hline
零 & líng & nombre & zéro \\
\hline
百 & bǎi & nombre & cent \\
\hline
千 & qiān & nombre & mille \\
\hline
万 & wàn & nombre & dix mille \\
\hline
\end{tabularx}
\end{center}

\begin{attention}[Le ton de \zh{一} (yī) — Règle de sandhi obligatoire]
\zh{一} seul se prononce \emph{yī} (1\ier{} ton).
\begin{itemize}
\item Devant un \textbf{4\ieme{} ton} : devient \emph{yí} (2\ieme{} ton). Ex: \zh{一个} \rightarrow \emph{yí ge}, \zh{一定} \rightarrow \emph{yí dìng}.
\item Devant un \textbf{1\ier{}, 2\ieme{} ou 3\ieme{} ton} : devient \emph{yì} (4\ieme{} ton). Ex: \zh{一千} \rightarrow \emph{yì qiān}, \zh{一百} \rightarrow \emph{yì bǎi}, \zh{一起} \rightarrow \emph{yì qǐ}.
\item Dans les nombres à plusieurs chiffres (dates, téléphones, codes) : on garde souvent \emph{yī} (ton original) pour la clarté, mais la règle de sandhi s'applique dans le flux de la parole naturelle.
\end{itemize}
À l'écrit, on garde toujours le caractère \zh{一}.
\end{attention}

\courssub{De 11 à 99 : la juxtaposition pure}

\begin{propriete}[Formation des nombres de 11 à 99]
Le chinois forme les nombres par \cle{juxtaposition stricte}, sans mot de liaison, sans trait d'union, sans accord.
\begin{itemize}
\item \textbf{11 à 19} : \zh{十} + unité → \zh{十一} (11, « dix-un »), \zh{十二} (12), \zh{十五} (15), \zh{十九} (19).
\item \textbf{20, 30, … 90} : unité + \zh{十} → \zh{二十} (20, « deux-dix »), \zh{三十} (30), \zh{九十} (90).
\item \textbf{Nombres composés (21–99)} : dizaine + unité → \zh{二十一} (21), \zh{三十五} (35), \zh{四十八} (48), \zh{九十九} (99).
\end{itemize}
Structure générale : \textbf{[chiffre]十[chiffre]}.
\end{propriete}

\begin{exemplebox}[De 11 à 99 — Contexte camerounais]
\zh{我们班有三十二个学生。}\\
\emph{Wǒmen bān yǒu sānshí'èr ge xuésheng.}\\
« Notre classe a trente-deux élèves. »\\
\zh{我妹妹今年十四岁。}\\
\emph{Wǒ mèimei jīnnián shísì suì.}\\
« Ma petite sœur a quatorze ans cette année. »\\
\zh{这件球衣八十九块。}\\
\emph{Zhè jiàn qiúyī bāshíjiǔ kuài.}\\
« Ce maillot coûte quatre-vingt-neuf yuans. »
\end{exemplebox}

\begin{savaistu}[Pourquoi le chinois est plus logique que le français]
En français, 72 = « soixante-douze » (60+12), 93 = « quatre-vingt-treize » (4×20+13). En chinois : \zh{七十二} (7-10-2), \zh{九十三} (9-10-3). Le système chinois est \textbf{additif et positionnel} (base 10 pure), exactement comme l'écriture chiffrée arabe. C'est un atout majeur pour le calcul mental !
\end{savaistu}

\courssub{Centaines, milliers, dizaines de milliers : 百, 千, 万}

\begin{propriete}[Juxtaposition multiplicative puis additive]
Le chinois construit les grands nombres par \cle{juxtaposition multiplicative} (chiffre × unité) puis \cle{additive} (somme des blocs), exactement comme l'écriture positionnelle arabe :
\medskip
\zh{三千二百五十六} = 3×1000 + 2×100 + 5×10 + 6 = 3256
\medskip
Unités clés : \zh{十} (10), \zh{百} (100), \zh{千} (1000), \zh{万} (10\,000).
\end{propriete}

\begin{attention}[Règle cruciale : \zh{一} obligatoire devant 百, 千, 万 ; interdit devant 十 (sauf nombres composés)]
\begin{itemize}
\item 10 = \zh{十} (jamais *\zh{一十} seul, sauf dans 11-19 : \zh{十一}).
\item 100 = \zh{一百} (obligatoire).
\item 1000 = \zh{一千} (obligatoire).
\item 10\,000 = \zh{一万} (obligatoire).
\item 110 = \zh{一百一十} (le \zh{一} devant \zh{十} réapparaît car il y a une centaine).
\end{itemize}
\end{attention}

\begin{exemplebox}[Grand nombres : 100 à 99\,999]
\zh{一百} (100) \quad \zh{两百} / \zh{二百} (200) \quad \zh{九百九十九} (999)\\
\zh{一千} (1000) \quad \zh{两千} / \zh{二千} (2000) \quad \zh{五千三百} (5300)\\
\zh{一万} (10\,000) \quad \zh{两万} (20\,000) \quad \zh{一万五千} (15\,000 — prix billet Mont Cameroun)\\
\zh{九万九千九百九十九} (99\,999)
\end{exemplebox}

\courssub{Le zéro interne : \zh{零} (líng)}

\begin{propriete}[Règle du zéro interne]
Quand une position intermédiaire est vide (zéro à l'intérieur du nombre), le chinois l'indique \textbf{obligatoirement} par un seul \zh{零}, quelle que soit la longueur de la série de zéros.
\end{propriete}

\begin{exemplebox}[Le zéro interne — Exemples progressifs]
\zh{一百零五} (\emph{yì bǎi líng wǔ}) = 105 (centaine + zéro + unité)\\
\zh{三百零二} (\emph{sān bǎi líng èr}) = 302\\
\zh{一千零二十} (\emph{yì qiān líng èrshí}) = 1020 (millier + zéro + dizaine)\\
\zh{一千零五} (\emph{yì qiān líng wǔ}) = 1005 (millier + zéro + unité)\\
\zh{两千零五} (\emph{liǎng qiān líng wǔ}) = 2005\\
\zh{一万零五} (\emph{yī wàn líng wǔ}) = 10\,005\\
\zh{一万零五十} (\emph{yī wàn líng wǔshí}) = 10\,050
\end{exemplebox}

\begin{attention}[Pièges mortels au marché de Douala / Yaoundé]
\begin{itemize}
\item \zh{一百零五} = 105 \quad vs \quad \zh{一百五} = 150 (zéro final non prononcé, cent-cinquante).
\item \zh{一千零五} = 1005 \quad vs \quad \zh{一千五} = 1500.
\item Un seul \zh{零} pour plusieurs zéros : 2005 = \zh{两千零五} (jamais *\zh{两千零零五}).
\item Zéro final : ne se dit \textbf{jamais}. 350 = \zh{三百五十} (pas *\zh{三百五十零}).
\end{itemize}
\end{attention}

\begin{popfigure}
\begin{tikzpicture}[font=\small, node distance=1.2cm and 1.5cm]
\node[draw=popblue, fill=popblueL, rounded corners, align=center, minimum width=2.2cm] (shi) {\zh{十} shí\\ 10};
\node[draw=poporange, fill=poporangeL, rounded corners, align=center, minimum width=2.2cm, right=of shi] (bai) {\zh{百} bǎi\\ 100};
\node[draw=poppurple, fill=poppurpleL, rounded corners, align=center, minimum width=2.2cm, right=of bai] (qian) {\zh{千} qiān\\ 1\,000};
\node[draw=popgreen, fill=popgreenL, rounded corners, align=center, minimum width=2.2cm, right=of qian] (wan) {\zh{万} wàn\\ 10\,000};

\node[draw=poppink, fill=poppinkL, rounded corners, align=center, below=of shi] (n35) {35 = \zh{三十五}\\ 3×10 + 5};
\node[draw=poppink, fill=poppinkL, rounded corners, align=center, below=of bai] (n350) {350 = \zh{三百五十}\\ 3×100 + 5×10};
\node[draw=poppink, fill=poppinkL, rounded corners, align=center, below=of n350] (n105) {105 = \zh{一百零五}\\ 1×100 + 0×10 + 5};
\node[draw=poppink, fill=poppinkL, rounded corners, align=center, below=of qian] (n1020) {1020 = \zh{一千零二十}\\ 1×1000 + 0×100 + 2×10};
\node[draw=poppink, fill=poppinkL, rounded corners, align=center, below=of wan] (n15000) {15\,000 = \zh{一万五千}\\ 1×10000 + 5×1000};

\draw[-{Stealth}, thick, popblue] (shi) -- (n35);
\draw[-{Stealth}, thick, poporange] (bai) -- (n350);
\draw[-{Stealth}, thick, poporange] (bai) -- (n105);
\draw[-{Stealth}, thick, poppurple] (qian) -- (n1020);
\draw[-{Stealth}, thick, popgreen] (wan) -- (n15000);
\end{tikzpicture}
\legende{Architecture du système numérique : bases 10, 100, 1000, 10000. \zh{零} marque toute position vide interne.}
\end{popfigure}

\courssub{Deux : \zh{二} (èr) ou \zh{两} (liǎng) ?}

\begin{propriete}[Distribution de « deux »]
\begin{itemize}
\item \zh{二} (\emph{èr}, 4\ieme{} ton) : \textbf{nombres purs, comptage, ordinaux, codes}.
      \zh{二十二} (22), \zh{一百零二} (102), \zh{第二} (2\ieme{}), \zh{二零二五年} (2025), téléphone : \zh{八二六} (8-2-6).
\item \zh{两} (\emph{liǎng}, 3\ieme{} ton) : \textbf{devant un classifieur / spécificatif} (quantité concrète).
      \zh{两个人} (2 personnes), \zh{两本书} (2 livres), \zh{两杯茶} (2 tasses).
\item \textbf{Zone grise acceptée} : devant \zh{百}, \zh{千}, \zh{万} → \zh{两百} / \zh{二百}, \zh{两千} / \zh{二千}, \zh{两万} / \zh{二万}. \zh{两百} est plus courant à l'oral pour les quantités (prix, foule).
\end{itemize}
\end{propriete}

\begin{center}
\begin{tabularx}{\linewidth}{|X|X|X|}
\hline
\rowcolor{popblueL} Situation & Forme correcte & Traduction \\
\hline
Compter / Nombre pur & 一、二、三 (yī, èr, sān) & un, deux, trois \\
\hline
Dans un nombre composé & 三十二 (sānshí'èr) & 32 \\
\hline
Devant classifieur & 两个学生 (liǎng ge xuésheng) & deux élèves \\
\hline
Devant 百 / 千 / 万 & 两百 / 二百 (liǎng bǎi / èr bǎi) & 200 (les deux possibles) \\
\hline
Ordinal & 第二课 (dì-èr kè) & la 2\ieme{} leçon \\
\hline
Date / Téléphone & 二零二五年 / 一三八... & 2025 / 138... \\
\hline
\end{tabularx}
\end{center}

\courssub{Nombre + classifieur + nom : structure obligatoire}

\begin{propriete}[Règle d'or du syntagme nominal quantifié]
En chinois, \textbf{jamais} Nombre + Nom directement. Structure immuable :
\medskip
\textbf{Nombre + Classifieur (Spec.) + Nom}
\medskip
\zh{个} (gè) est le classifieur générique (personnes, objets divers). \zh{本} (běn) pour livres, \zh{件} (jiàn) pour vêtements, \zh{杯} (bēi) pour tasses, \zh{张} (zhāng) pour objets plats (billets, tables, papier).
\end{propriete}

\begin{exemplebox}[Nombres + Classifieurs + Noms]
\zh{三个学生} (\emph{sān ge xuésheng}) « trois élèves »\\
\zh{五本书} (\emph{wǔ běn shū}) « cinq livres »\\
\zh{两杯茶} (\emph{liǎng bēi chá}) « deux tasses de thé »\\
\zh{一件衬衫} (\emph{yí jiàn chènshān}) « une chemise » (sandhi \zh{一} $\rightarrow$ \emph{yí})\\
\zh{我们学校有一万五千个学生。}\\
\emph{Wǒmen xuéxiào yǒu yī wàn wǔqiān ge xuésheng.}\\
« Notre école a quinze mille élèves. »\\
\zh{这件衬衫两百块。}\\
\emph{Zhè jiàn chènshān liǎng bǎi kuài.}\\
« Cette chemise coûte deux cents yuans. » (\zh{块} kuài = classifieur monétaire oral pour \zh{元} yuán / franc CFA).
\end{exemplebox}

\courssub{Applications concrètes : prix, âge, dates, téléphone}

\begin{itemize}
\item \textbf{Prix} : \zh{多少钱？} (\emph{Duōshao qián ?}) — \zh{五百法郎。} (\emph{Wǔbǎi fǎláng.}) « 500 francs. » (\zh{块} kuài souvent utilisé pour francs CFA à l'oral).
\item \textbf{Âge} : Nombre + \zh{岁} (suì). \zh{我十六岁。} (\emph{Wǒ shíliù suì.}) « J'ai 16 ans. » \textbf{Pas de verbe « avoir » !}
\item \textbf{Date (Année)} : Chiffres un par un + \zh{年}. \zh{二零二五年} (\emph{èr líng èr wǔ nián}) = 2025.
\item \textbf{Date (Mois/Jour)} : Nombre + \zh{月} / \zh{号} (oral) ou \zh{日} (écrit). \zh{十月一日} (1\ier{} oct), \zh{三月八号} (8 mars).
\item \textbf{Téléphone} : Chiffres un par un. \zh{1} se prononce \emph{yāo} (避免 confusion avec \zh{7} \emph{qī}). \zh{我的电话是六九零一二三四五六。} (\emph{Wǒ de diànhuà shì liù jiǔ líng yāo èr sān sì wǔ liù.})
\end{itemize}

\begin{methode}[Lire / Écrire un grand nombre en chinois (Méthode pas à pas)]
\begin{enumerate}
\item \textbf{Grouper par 4 chiffres} à partir de la droite (unité \zh{万}).
\item \textbf{Traiter chaque groupe} : dire le chiffre + son unité (\zh{千}, \zh{百}, \zh{十}), insérer \zh{零} pour trou interne.
\item \textbf{Relier les groupes} par \zh{万}.
\item \textbf{Vérifier} : 12\,050 $\rightarrow$ 1万 / 2千 / 0百 / 5十 $\rightarrow$ \zh{一万二千零五十}.
\end{enumerate}
\end{methode}

\begin{exoresolu}[Écrire les nombres en caractères chinois]
Énoncé : Écrivez en caractères : 17 ; 63 ; 99 ; 210 ; 408 ; 1111 ; 15\,000 ; 10\,005. Traduisez : \zh{这个芒果两百法郎。} / \zh{我有三个姐姐。} Corrigez : *\zh{一百五} (pour 105) ; *\zh{二个老师}.
\tcblower
Solution détaillée :
\begin{itemize}
\item 17 = \zh{十七} ; 63 = \zh{六十三} ; 99 = \zh{九十九}.
\item 210 = \zh{二百一十} (ou \zh{两百一十}) ; 408 = \zh{四百零八} (zéro interne) ; 1111 = \zh{一千一百一十一}.
\item 15\,000 = \zh{一万五千} ; 10\,005 = \zh{一万零五} (un seul \zh{零}).
\item \zh{这个芒果两百法郎。} $\rightarrow$ « Cette mangue coûte 200 francs. »
\item \zh{我有三个姐姐。} $\rightarrow$ « J'ai trois grandes sœurs. » (\zh{姐姐} = grande sœur).
\item 105 = \zh{一百零五} (obligation de \zh{零}) ; « deux professeurs » = \zh{两个老师} (\zh{两} devant classifieur).
\end{itemize}
\end{exoresolu}

\begin{exercice}[Entraînement nombres — À vous]
\begin{enumerate}
\item Écrivez en caractères : 28 ; 50 ; 101 ; 306 ; 1200 ; 20\,080 ; 99\,999.
\item Lisez à haute voix (pinyin + tons) : \zh{三十五} ; \zh{六百零六} ; \zh{两千零八} ; \zh{一万两千} ; \zh{八万零八十}.
\item Traduisez en chinois : « 350 francs » ; « 22 élèves » ; « J'ai 18 ans » ; « Le 8 mars 2025 ».
\item Corrigez les fautes : *\zh{二十十} (30) ; *\zh{一千五} (1005) ; *\zh{三个书} ; *\zh{二个人}.
\item Mini-dialogue : Au marché, demandez le prix d'un sac (\zh{包} bāo), le vendeur répond 12\,500 francs. Vous négociez 10\,000 francs.
\end{enumerate}
\end{exercice}

\begin{corrige}[Exercice « Entraînement nombres »]
\begin{enumerate}
\item 28=\zh{二十八} ; 50=\zh{五十} ; 101=\zh{一百零一} ; 306=\zh{三百零六} ; 1200=\zh{一千二百} (ou \zh{十二百} oral) ; 20\,080=\zh{两万零八十} ; 99\,999=\zh{九万九千九百九十九}.
\item \zh{三十五} (sānshíwǔ) ; \zh{六百零六} (liùbǎi líng liù) ; \zh{两千零八} (liǎng qiān líng bā) ; \zh{一万两千} (yī wàn liǎng qiān) ; \zh{八万零八十} (bā wàn líng bāshí).
\item \zh{三百五十法郎} ; \zh{二十二个学生} ; \zh{我十八岁} ; \zh{二零二五年三月八号}.
\item 30=\zh{三十} ; 1005=\zh{一千零五} ; \zh{三本书} (\zh{本} pour livres) ; \zh{两个人}.
\item \zh{A: 这个包多少钱？ B: 一万二千五百法郎。 A: 太贵了，一万块行吗？ B: 好吧，成交。}
\end{enumerate}
\end{corrige}

\coursec{La répartition : 有的……有的……}

\courssub{Structure et sens}

\begin{propriete}[\zh{有的……有的……} (yǒude… yǒude…)]
Exprime une \textbf{répartition} au sein d'un groupe : « Certains…, d'autres… », « Les uns…, les autres… ».
\medskip
\textbf{Schéma :} \zh{有的} + \textbf{Sujet/Sujet omis} + \textbf{Prédicat 1} ， \zh{有的} + \textbf{Sujet/Sujet omis} + \textbf{Prédicat 2}。
\medskip
Le sujet après \zh{有的} peut être explicité (nom) ou omis (reprise du thème).
\end{propriete}

\begin{exemplebox}[有的……有的…… — Exemples variés]
\zh{我们班有的学生喜欢足球，有的喜欢篮球。}\\
\emph{Wǒmen bān yǒude xuésheng xǐhuan zúqiú, yǒude xǐhuan lánqiú.}\\
« Dans notre classe, certains élèves aiment le foot, d'autres le basket. »\\
\zh{今天天气真奇怪，有的地方下雨，有的地方出太阳。}\\
\emph{Jīntiān tiānqì zhēn qíguài, yǒude dìfang xiàyǔ, yǒude dìfang chū tàiyáng.}\\
« La météo d'aujourd'hui est bizarre : par endroits il pleut, par endroits il fait soleil. »\\
\zh{这些水果有的是芒果，有的是菠萝。}\\
\emph{Zhèxiē shuǐguǒ yǒude shì mángguǒ, yǒude shì bōluó.}\\
« Parmi ces fruits, certains sont des mangues, d'autres des ananas. » (Jugement \zh{是})
\end{exemplebox}

\begin{attention}[Pièges francophones]
\begin{itemize}
\item \textbf{Pas de « les uns / les autres » redondant :} \zh{有的} porte déjà la notion de « certains/uns ». Ne dites pas *\zh{有的学生……其他的学生……} (utilisez \zh{其他的} pour « les autres » si besoin, mais \zh{有的……有的} suffit).
\item \textbf{Sujet unique vs Sujets différents :} Si le sujet est le même groupe, on l'omet la 2\ieme{} fois. \zh{有的同学……，有的……}. Si groupes différents, on le répète : \zh{有的男生……，有的女生……}.
\item \textbf{Ponctuation :} Virgule chinoise \zh{，} entre les deux membres.
\end{itemize}
\end{attention}

\courssub{Extension : \zh{有的……有的……还有的……}}

\begin{exemplebox}[Trois groupes ou plus]
\zh{周末有的同学回家，有的去图书馆，还有的去踢球。}\\
\emph{Zhōumò yǒude tóngxué huí jiā, yǒude qù túshūguǎn, hái yǒude qù tī qiú.}\\
« Le week-end, certains rentrent chez eux, d'autres vont à la bibliothèque, et d'autres encore vont jouer au foot. »
\end{exemplebox}

\coursec{La mise en relief : 是……的}

\courssub{Concept clé : focus sur un détail d'action \textbf{passée} (ou future programmée)}

\begin{propriete}[Structure \zh{是……的} (shì… de)]
Utilisée pour \textbf{mettre en relief} (focus) un élément de circonstance (temps, lieu, manière, but, agent) d'une action \textbf{accomplie} (ou décidée). Le verbe principal se place \textbf{avant} \zh{的} (ou à la fin).
\medskip
\textbf{Schéma :} \textbf{Sujet + 是 + \cle{Élément mis en relief} + Verbe + 的} \quad (ou \quad \textbf{Sujet + 是 + \cle{Élément} + 的 + Verbe} — moins courant pour verbes d'action).
\medskip
\zh{的} marque la fin de la focalisation. \zh{是} peut être omis à l'affirmatif (surtout à l'oral), mais \textbf{obligatoire à la négation et à l'interrogation}.
\end{propriete}

\begin{exemplebox}[Focus sur Temps, Lieu, Manière, Moyen, But]
\textbf{Temps :} \zh{我是昨天到的雅温得。} (\emph{Wǒ shì zuótiān dào de Yāwēndé.}) « C'est \textbf{hier} que j'ai atteint Yaoundé. »\\
\textbf{Lieu :} \zh{我们是在市场买的水果。} (\emph{Wǒmen shì zài shìchǎng mǎi de shuǐguǒ.}) « C'est \textbf{au marché} qu'on a acheté les fruits. »\\
\textbf{Moyen/Transport :} \zh{他不是坐飞机去的北京，是坐火车去的。} (\emph{Tā bùshì zuò fēijī qù de Běijīng, shì zuò huǒchē qù de.}) « Ce n'est \textbf{pas en avion} qu'il est allé à Pékin, c'est \textbf{en train}. »\\
\textbf{But :} \zh{我是来学中文的。} (\emph{Wǒ shì lái xué Zhōngwén de.}) « Je suis venu \textbf{pour apprendre le chinois}. »\\
\textbf{Agent (passif/actif) :} \zh{这本书是王老师教的。} (\emph{Zhè běn shū shì Wáng lǎoshī jiāo de.}) « Ce livre a été enseigné \textbf{par M. Wang}. »
\end{exemplebox}

\courssub{Négation et Interrogation : \zh{是} obligatoire}

\begin{propriete}[Négation : \zh{不是……的} ; Question : \zh{是……的吗？}]
\begin{itemize}
\item \textbf{Négation} : \zh{不是} + Élément + \zh{的}. \zh{我不是坐大巴去的，是坐火车去的。}
\item \textbf{Question focus} : \zh{是} + Élément + \zh{的吗？} \zh{你是几点来的？} / \zh{你是坐车来的吗？}
\item \textbf{Question ouverte} : \zh{怎么} / \zh{几点} / \zh{谁} + \zh{的} ? \zh{你怎么来的？} (Comment es-tu venu ? — focus sur manière).
\end{itemize}
\end{propriete}

\begin{attention}[Différence cruciale \zh{了} vs \zh{是……的}]
\begin{center}
\begin{tabularx}{\linewidth}{|X|X|X|}
\hline
\rowcolor{popblueL} Structure & Fonction & Exemple \\
\hline
\zh{Verb + 了} & \textbf{Accomplissement / Changement d'état} (nouvelle info). & \zh{我去了北京。} (Je suis allé à Pékin — fait nouveau). \\
\hline
\zh{是……的} & \textbf{Focus sur détail} d'un fait \textbf{déjà connu/acquis}. & \zh{我是去年去的北京。} (C'est \textbf{l'an dernier} que j'y suis allé — on sait que tu y es allé). \\
\hline
\end{tabularx}
\end{center}
\begin{itemize}
\item \zh{我吃了饭。} (J'ai mangé — action terminée).
\item \zh{我是在食堂吃的饭。} (C'est \textbf{à la cantine} que j'ai mangé — lieu mis en relief).
\item \textbf{Erreur fréquente :} *\zh{我是在食堂吃了饭。} (Mélange incorrect).
\end{itemize}
\end{attention}

\courssub{Position du verbe : avant \zh{的} (standard) ou après \zh{的} (verbes d'état/mentaux)}

\begin{definition}[Position du verbe]
\begin{itemize}
\item \textbf{Verbes d'action (aller, venir, acheter, manger, faire) :} Verbe \textbf{avant} \zh{的}. \zh{我是昨天买的书。}
\item \textbf{Verbes d'état/mentaux / \zh{在} + lieu :} Verbe \textbf{après} \zh{的} possible (structure \zh{是……的} = « C'est … que… »).
      \zh{这本书是我的。} (\emph{Zhè běn shū shì wǒ de.}) « Ce livre est \textbf{à moi}. » (Pas de verbe d'action).
      \zh{会议是在会议室开的。} (\emph{Huìyì shì zài huìyìshì kāi de.}) « La réunion s'est tenue \textbf{dans la salle de réunion}. » (\zh{在} phrase avant verbe).
\end{itemize}
\end{definition}

\coursec{Synthèse, comparaison et entraînement global}

\courssub{Tableau récapitulatif des quatre structures}

\begin{center}
\begin{tabularx}{\linewidth}{|X|X|X|X|}
\hline
\rowcolor{popblueL} Structure & Fonction communicative & Schéma clé & Exemple type \\
\hline
\zh{如果……，那么/就……} & Hypothèse conditionnelle (réel/futur) & 如果 Cond, 那么/就 Conséquence & \zh{如果下雨，我就不去。} \\
\hline
Nombres + Classifieurs & Quantifier, prix, âge, date, tel & Nombre + Classifieur + Nom & \zh{两个朋友}, \zh{一万五千法郎} \\
\hline
\zh{有的……有的……} & Répartition / Énumération alternative & 有的 A, 有的 B & \zh{有的喝茶，有的喝咖啡。} \\
\hline
\zh{是……的} & Focus sur détail d'action passée & 是 [Temps/Lieu/Moyen] Verb 的 & \zh{我是坐飞机去的。} \\
\hline
\end{tabularx}
\end{center}

\courssub{Exercices de synthèse : Transformation et Traduction}

\begin{exercice}[Synthèse — Niveau Bac 1ère A]
\begin{enumerate}
\item \textbf{Transformation (Hypothèse) :} Réécrivez avec \zh{如果……，就……} : « Demain j'ai cours. Je ne vais pas au marché. » $\rightarrow$ \zh{如果明天有课，我就不去市场。} Faites de même : « S'il fait beau, nous grimpons le Mont Cameroun. » / « Si tu n'as pas d'argent, tu ne peux pas acheter ce téléphone. »
\item \textbf{Transformation (Focus \zh{是……的}) :} Mettez l'élément souligné en relief.
      \begin{itemize}
      \item Il est allé à Douala \underline{en train}. (\zh{坐火车})
      \item J'ai acheté ce dictionnaire \underline{hier}. (\zh{昨天})
      \item Nous avons appris le chinois \underline{au Cameroun}. (\zh{在喀麦隆})
      \end{itemize}
\item \textbf{Traduction (Français $\rightarrow$ Chinois) :}
      \begin{enumerate}
      \item « Certains élèves de notre école viennent de Yaoundé, d'autres viennent de Douala. » (\zh{有的……有的……})
      \item « Ce n'est pas 105 francs, c'est 150 francs. » (Nombres + \zh{是……的} possible pour correction).
      \item « Si tu viens demain, appelle-moi. » (\zh{打电话} / \zh{微信}).
      \item « Mon père est né en 1970. » (Date : \zh{一九七零年}).
      \end{enumerate}
\item \textbf{Production écrite (60-80 caractères) :} Écrivez un court message WeChat à un ami chinois pour lui expliquer : la météo de demain (hypothèse), le prix du billet de bus (nombre), ce que font les camarades (répartition), et comment vous irez à l'école (focus \zh{是……的}).
\end{enumerate}
\end{exercice}

\begin{corrige}[Exercice « Synthèse »]
\begin{enumerate}
\item \zh{如果明天天气好，我们就去爬摄影山。} / \zh{如果你没钱，就买不了这个手机。} (\zh{买不了} = ne pas pouvoir acheter).
\item \zh{他是坐火车去的杜阿拉。} / \zh{我是昨天买的这本词典。} / \zh{我们是在喀麦隆学的中文。}
\item \begin{enumerate}
      \item \zh{我们学校有的学生是雅温得来的，有的是杜阿拉来的。} (ou \zh{来自}).
      \item \zh{不是一百零五法郎，是一百五十法郎。} (Correction de prix).
      \item \zh{如果你明天来，给我打电话 / 发微信。}
      \item \zh{我爸爸是一九七零年生的。} (Focus sur année naissance).
      \end{enumerate}
\item \textbf{Modèle de message :}
\zh{小李：你好！如果明天不下雨，我们就去踢球。球票两百法郎。有的同学踢前锋，有的踢后卫。我是骑摩托车去的。你来吗？}
(\emph{Xiǎolǐ: Nǐ hǎo! Rúguǒ míngtiān bù xiàyǔ, wǒmen jiù qù tī qiú. Qiúpiào liǎng bǎi fǎláng. Yǒude tóngxué tī qiánfēng, yǒude tī hòuwèi. Wǒ shì qí mótuōchē qù de. Nǐ lái ma?})
\end{enumerate}
\end{corrige}

\courssub{Expression orale guidée : Jeu de rôle « Au marché de Mokolo »}

\begin{experience}[Jeu de rôle à 3 : Client / Vendeur / Observateur]
\textbf{Contexte :} Marché de Mokolo (Yaoundé). Achat de maillots de foot (\zh{球衣} qiúyī).
\textbf{Rôles :}
\begin{itemize}
\item \textbf{Client :} Demande prix (\zh{多少钱}), négocie (nombres), choisit quantité (classifieur \zh{件}), hypothèse : \zh{如果便宜，我就多买两件}.
\item \textbf{Vendeur :} Donne prix (grand nombre : 1500, 2000…), propose réduction si quantité (\zh{如果买三件，就便宜点}), précise origine (\zh{是广州产的} — focus \zh{是……的}).
\item \textbf{Observateur :} Note les structures cibles utilisées (✓/✗), corrige tons (\zh{一}, \zh{不}) et classifieurs.
\end{itemize}
\textbf{Durée :} 5 min préparation, 3 min jeu, 2 min feedback.
\end{experience}

\begin{aretenir}[L'essentiel de la Leçon 6 — Check-list révision]
\begin{itemize}
\item \textbf{如果……，那么/就……} : Ordre Sujet-\zh{就/那么}-Verbe. \zh{就} = oral/immédiat ; \zh{那么} = écrit/logique. Négation avant verbe concerné.
\item \textbf{Nombres} : Juxtaposition stricte. \zh{零} pour zéro interne (un seul). \zh{万} = 10\,000. \zh{一} obligatoire devant \zh{百/千/万}. \zh{两} devant classifieur, \zh{二} dans nombres/ordinaux. \textbf{Toujours} Nombre + Classifieur + Nom.
\item \textbf{有的……有的……} : Répartition. Sujet omis si identique. Virgule entre les deux.
\item \textbf{是……的} : Focus sur détail (temps/lieu/moyen/but) d'action \textbf{passée/connue}. \zh{是} omis à l'affirmatif, \textbf{obligatoire} en négation (\zh{不是……的}) et question (\zh{是……的吗？}). Verbe d'action avant \zh{的}. Différencier de \zh{了} (accomplissement simple).
\end{itemize}
\end{aretenir}

\coursec{La répartition : 有的……有的……}

Dans la section précédente, nous avons appris à compter et à quantifier grâce aux nombres chinois. Mais dans la vie quotidienne, on ne compte pas toujours : souvent, on veut simplement décrire ce que font les différents membres d'un groupe. Au marché de Mokolo, dans la cour du lycée de Yaoundé, au stade Ahmadou Ahidjo, chacun fait quelque chose de différent. Comment dire en chinois « certains font ceci, d'autres font cela » ? C'est exactement le rôle de la structure \zh{有的……有的……}, que nous allons maintenant étudier.

\courssub{Observation : un samedi au marché}

\begin{textebox}[Au marché de Yaoundé]
星期六，市场上很热闹。有的人买水果，有的人卖鱼，还有的人卖衣服。市场旁边有一个小公园：公园里，有的人跑步，有的人唱歌，还有的人跳舞。下午，我们去看足球比赛。体育场里，有的人大声喊，有的人拍照，还有的人吃花生。喀麦隆的星期六真有意思！\\
\emph{Xīngqīliù, shìchǎng shang hěn rènao. Yǒude rén mǎi shuǐguǒ, yǒude rén mài yú, hái yǒude rén mài yīfu. Shìchǎng pángbiān yǒu yí ge xiǎo gōngyuán : gōngyuán lǐ, yǒude rén pǎobù, yǒude rén chànggē, hái yǒude rén tiàowǔ. Xiàwǔ, wǒmen qù kàn zúqiú bǐsài. Tǐyùchǎng lǐ, yǒude rén dàshēng hǎn, yǒude rén pāizhào, hái yǒude rén chī huāshēng. Kāmàilóng de xīngqīliù zhēn yǒu yìsi !}\\
« Samedi, le marché est très animé. Certains achètent des fruits, d'autres vendent du poisson, d'autres encore vendent des vêtements. À côté du marché, il y a un petit parc : dans le parc, certains courent, d'autres chantent, d'autres encore dansent. L'après-midi, nous allons voir un match de football. Dans le stade, certains crient très fort, d'autres prennent des photos, d'autres encore mangent des cacahuètes. Les samedis au Cameroun sont vraiment intéressants ! »
\end{textebox}

\textbf{Questions d'observation}
\begin{enumerate}
\item Combien d'activités différentes sont décrites au marché ? Au stade ?
\item Quel mot revient devant chaque activité ? Où se place-t-il dans la phrase ?
\item Où se placent les mots de lieu \zh{市场上}, \zh{公园里}, \zh{体育场里} ?
\item Quel mot introduit la troisième activité de chaque série ?
\end{enumerate}

\textbf{Glossaire}

\begin{center}
\begin{tabularx}{\linewidth}{|X|X|X|X|}
\hline
\rowcolor{popblueL} Caractères & Pinyin & Nature & Traduction \\
\hline
有的 & yǒude & pronom indéfini & certains, certaines \\
\hline
还有的 & hái yǒude & pronom indéfini & d'autres encore \\
\hline
热闹 & rènao & adjectif & animé, bruyant (de joie) \\
\hline
卖 & mài & verbe & vendre \\
\hline
跑步 & pǎobù & verbe & courir (sport) \\
\hline
唱歌 & chànggē & verbe & chanter \\
\hline
跳舞 & tiàowǔ & verbe & danser \\
\hline
比赛 & bǐsài & nom / verbe & match ; disputer un match \\
\hline
拍照 & pāizhào & verbe & prendre des photos \\
\hline
花生 & huāshēng & nom & cacahuète, arachide \\
\hline
\end{tabularx}
\end{center}

\courssub{La règle : forme et emploi}

\begin{propriete}[Structure de la répartition 有的……有的……]
La structure \cle{有的……有的……} décrit une \textbf{répartition d'activités à l'intérieur d'un même groupe ou d'un même lieu}, annoncé en tête de phrase.

\textbf{Schéma :}

(Cadre : lieu / groupe) + 有的 (+ nom) + verbe，有的 (+ nom) + verbe（，还有的 + verbe）。

Exemple canonique : \zh{操场上，有的同学跑步，有的同学打球。} \emph{(Cāochǎng shang, yǒude tóngxué pǎobù, yǒude tóngxué dǎqiú.)} « Sur le terrain de sport, certains élèves courent, d'autres jouent au ballon. »

\textbf{Points essentiels :}
\begin{itemize}
\item \zh{有的} signifie « certains » ; il peut être suivi d'un nom (\zh{有的人}, \zh{有的同学}) ou utilisé seul (\zh{有的跑步，有的打球}).
\item Chaque proposition introduite par \zh{有的} contient \textbf{obligatoirement un verbe} : on ne peut pas dire « certains au foot » sans verbe.
\item On peut chaîner trois propositions ou plus ; la dernière est souvent introduite par \zh{还有的} « d'autres encore » : \zh{有的……有的……还有的……}.
\item Le cadre (lieu ou groupe) se place \textbf{en tête de phrase}, avant le premier \zh{有的}.
\item \zh{有的} peut aussi fonctionner comme déterminant devant un nom sujet unique : \zh{有的同学喜欢汉语。} « Certains élèves aiment le chinois. » Dans ce cas, il n'y a pas de répartition, mais une simple restriction à une partie du groupe.
\end{itemize}
\end{propriete}

\begin{popfigure}
\begin{tikzpicture}[
  node distance=0.5cm,
  cadre/.style={rectangle, rounded corners, draw=popdark, fill=popgoldL, thick, align=center, font=\small, minimum height=0.9cm},
  bloc/.style={rectangle, rounded corners, draw=popblue, fill=popblueL, thick, align=center, font=\small, minimum height=0.9cm},
  blocb/.style={rectangle, rounded corners, draw=popgreen, fill=popgreenL, thick, align=center, font=\small, minimum height=0.9cm},
  blocc/.style={rectangle, rounded corners, draw=poppurple, fill=poppurpleL, thick, align=center, font=\small, minimum height=0.9cm},
  fleche/.style={-{Stealth[length=2.5mm]}, thick, popmuted}
]
\node[cadre, align=center] (cadre){Cadre en tête\\ \zh{公园里} \emph{gōngyuán lǐ}\\ « dans le parc »};
\node[bloc, below left=1.1cm and -0.2cm of cadre, align=center] (v1){有的 + V1\\ \zh{有的跑步}\\ « certains courent »};
\node[blocb, below=1.1cm of cadre, align=center] (v2){有的 + V2\\ \zh{有的唱歌}\\ « d'autres chantent »};
\node[blocc, below right=1.1cm and -0.2cm of cadre, align=center] (v3){还有的 + V3\\ \zh{还有的跳舞}\\ « d'autres encore dansent »};
\draw[fleche] (cadre.south) -- (v1.north);
\draw[fleche] (cadre.south) -- (v2.north);
\draw[fleche] (cadre.south) -- (v3.north);
\end{tikzpicture}
\legende{Schéma de la structure : le cadre est annoncé d'abord, puis chaque sous-groupe est introduit par 有的 suivi d'un verbe.}
\end{popfigure}

\courssub{Exemples contextualisés}

\begin{exemplebox}[Au marché, à l'école, au stade]
\begin{enumerate}
\item \zh{市场上，有的人买水果，有的人卖鱼。} \emph{(Shìchǎng shang, yǒude rén mǎi shuǐguǒ, yǒude rén mài yú.)} « Au marché, certains achètent des fruits, d'autres vendent du poisson. »
\item \zh{下课后，有的同学做作业，有的踢足球，还有的聊天。} \emph{(Xiàkè hòu, yǒude tóngxué zuò zuòyè, yǒude tī zúqiú, hái yǒude liáotiān.)} « Après les cours, certains font leurs devoirs, d'autres jouent au foot, d'autres encore bavardent. »
\item \zh{食堂里，有的学生吃米饭，有的学生吃恩多雷。} \emph{(Shítáng lǐ, yǒude xuésheng chī mǐfàn, yǒude xuésheng chī Ēnduōléi.)} « À la cantine, certains élèves mangent du riz, d'autres mangent du ndolé. » (\zh{恩多雷} est la transcription chinoise du mot « ndolé », plat camerounais.)
\item \zh{星期天，我们家的人有的看电视，有的洗衣服，还有的休息。} \emph{(Xīngqītiān, wǒmen jiā de rén yǒude kàn diànshì, yǒude xǐ yīfu, hái yǒude xiūxi.)} « Dimanche, dans notre famille, certains regardent la télé, d'autres font la lessive, d'autres encore se reposent. »
\item \zh{公共汽车上，有的人看手机，有的人睡觉。} \emph{(Gōnggòng qìchē shang, yǒude rén kàn shǒujī, yǒude rén shuìjiào.)} « Dans le bus, certains regardent leur téléphone, d'autres dorment. »
\end{enumerate}
\end{exemplebox}

\courssub{Comparaison avec le français}

\begin{center}
\begin{tabularx}{\linewidth}{|X|X|X|}
\hline
\rowcolor{popblueL} Structure & Exemple en caractères + pinyin & Traduction \\
\hline
Cadre + 有的 + nom + V，有的 + nom + V & \zh{教室里，有的学生读书，有的学生写字。} \emph{Jiàoshì lǐ, yǒude xuésheng dúshū, yǒude xuésheng xiězì.} & En classe, les uns lisent, les autres écrivent. \\
\hline
有的 seul (nom sous-entendu) & \zh{有的喝茶，有的喝咖啡。} \emph{Yǒude hē chá, yǒude hē kāfēi.} & Certains boivent du thé, d'autres du café. \\
\hline
Chaîne à trois termes & \zh{有的唱歌，有的跳舞，还有的聊天。} \emph{Yǒude chànggē, yǒude tiàowǔ, hái yǒude liáotiān.} & Les uns chantent, les autres dansent, d'autres encore bavardent. \\
\hline
有的 + nom (sans répartition) & \zh{有的学生喜欢足球。} \emph{Yǒude xuésheng xǐhuan zúqiú.} & Certains élèves aiment le football. \\
\hline
\end{tabularx}
\end{center}

Le français dit « les uns… les autres… » ou « certains… d'autres… » avec un seul verbe parfois sous-entendu (« les uns au foot, les autres au basket »). Le chinois, lui, \textbf{répète \zh{有的} devant chaque proposition et exige un verbe dans chaque partie}. Autre différence : le français place souvent le lieu à la fin (« certains courent \emph{dans le parc} »), alors que le chinois place le cadre \textbf{en tête}.

\begin{attention}[Pièges fréquents des francophones]
\begin{itemize}
\item \textbf{Ne pas confondre \zh{有的} (certains) et \zh{有} (avoir / il y a).} \zh{有的} est inséparable : on ne peut pas utiliser \zh{有的人} pour dire « il y a des gens ». Comparez : \zh{公园里有人。} « Il y a des gens dans le parc » (existence) et \zh{公园里，有的人跑步。} « Dans le parc, certains courent » (répartition).
\item La phrase \zh{公园有人跑步，有人唱歌} est grammaticalement possible, mais elle est moins élégante et moins claire que \zh{公园里，有的人跑步，有的人唱歌}. Préférez toujours la structure complète avec le cadre en tête.
\item N'oubliez pas le verbe : *\zh{有的足球，有的篮球} est incorrect. Dites \zh{有的踢足球，有的打篮球}.
\item Ne traduisez pas « les uns » par \zh{一些} : \zh{一些} est un quantificateur (« quelques ») qui se place devant un nom (\zh{一些学生} « quelques élèves »), pas un pronom de répartition. On ne peut pas dire *\zh{一些跑步，一些唱歌}.
\item Attention à la ponctuation : les propositions introduites par \zh{有的} sont séparées par la virgule chinoise « ， », jamais par le signe de coordination « 、 », réservé aux énumérations de noms.
\end{itemize}
\end{attention}

\begin{methode}[Construire une phrase avec 有的……有的……]
\begin{enumerate}
\item Identifiez le \textbf{cadre} : le lieu (\zh{市场上}) ou le groupe (\zh{我们班的同学}).
\item Placez ce cadre en tête de phrase, suivi d'une virgule chinoise « ， ».
\item Listez les activités : pour chacune, écrivez \zh{有的} (+ nom si besoin) + verbe (+ complément).
\item S'il y a trois activités ou plus, introduisez la dernière par \zh{还有的}.
\item Vérifiez : chaque partie a-t-elle un verbe ? Le cadre est-il bien en tête ?
\end{enumerate}
\end{methode}

\begin{exoresolu}[Transformation guidée]
Énoncé : traduisez en chinois : « Dans la cour de récréation, certains élèves jouent au football, d'autres jouent au basket, d'autres encore bavardent. »
\tcblower
Solution détaillée :
\begin{enumerate}
\item Cadre : « dans la cour de récréation » → \zh{操场上} (placé en tête).
\item Activité 1 : jouer au football → \zh{有的同学踢足球}.
\item Activité 2 : jouer au basket → \zh{有的同学打篮球}.
\item Activité 3 (dernière) : bavarder → \zh{还有的聊天}.
\end{enumerate}
Phrase complète : \zh{操场上，有的同学踢足球，有的同学打篮球，还有的聊天。} \emph{(Cāochǎng shang, yǒude tóngxué tī zúqiú, yǒude tóngxué dǎ lánqiú, hái yǒude liáotiān.)} Vérification : cadre en tête, trois verbes présents, \zh{还有的} pour la dernière activité.
\end{exoresolu}

\begin{exercice}[À vous de jouer]
\begin{enumerate}
\item Traduisez : « Au stade, certains crient, d'autres chantent, d'autres encore prennent des photos. »
\item Traduisez : « Le matin, dans ma famille, certains mangent du pain, d'autres mangent des beignets. »
\item Corrigez la phrase fautive : *\zh{有的人公园跑步，有的人唱歌。}
\item Complétez avec \zh{有} ou \zh{有的} : \zh{市场上\_\_\_很多人，\_\_\_人买水果，\_\_\_人卖鱼。}
\end{enumerate}
\end{exercice}

\begin{corrige}[Exercice « À vous de jouer »]
\begin{enumerate}
\item \zh{体育场里，有的大声喊，有的唱歌，还有的拍照。} \emph{(Tǐyùchǎng lǐ, yǒude dàshēng hǎn, yǒude chànggē, hái yǒude pāizhào.)}
\item \zh{早上，我们家的人有的吃面包，有的吃油炸糕。} \emph{(Zǎoshang, wǒmen jiā de rén yǒude chī miànbāo, yǒude chī yóuzhágāo.)}
\item Le cadre doit être en tête : \zh{公园里，有的人跑步，有的人唱歌。} \emph{(Gōngyuán lǐ, yǒude rén pǎobù, yǒude rén chànggē.)}
\item \zh{市场上\textbf{有}很多人，\textbf{有的}人买水果，\textbf{有的}人卖鱼。} (premier blanc : existence « il y a » ; les deux suivants : répartition « certains »).
\end{enumerate}
\end{corrige}

\begin{aretenir}[L'essentiel de la section]
\begin{itemize}
\item \zh{有的……有的……} = « certains… d'autres… » : répartition d'activités dans un même groupe ou lieu.
\item Le cadre (lieu ou groupe) se place \textbf{en tête de phrase}.
\item Chaque \zh{有的} est suivi d'un \textbf{verbe obligatoire} ; la dernière activité peut être introduite par \zh{还有的}.
\item \zh{有的} (certains) est inséparable et ne se confond pas avec \zh{有} (il y a) ni avec \zh{一些} (quelques).
\end{itemize}
\end{aretenir}

\begin{experience}[Description orale : la cour de récréation]
Par groupes de trois : observez (ou imaginez) la cour de votre lycée pendant la récréation. Chaque élève produit deux phrases avec \zh{有的……有的……} pour décrire la scène, puis le groupe présente sa description complète à la classe en ajoutant une troisième activité avec \zh{还有的}. Les autres groupes vérifient : cadre en tête ? verbe dans chaque partie ?
\end{experience}

\coursec{La mise en relief : 是……的}

Dans la section précédente, nous avons appris à répartir des personnes ou des choses en groupes grâce à \zh{有的……有的……}. Nous allons maintenant découvrir une structure extrêmement fréquente en chinois parlé : la construction \zh{是……的}, qui permet d'\cle{insister sur les circonstances} d'une action déjà accomplie. C'est l'équivalent chinois de la phrase clivée française « c'est… que… ».

\courssub{Observation : découvrir la structure en contexte}

Lisons d'abord ce petit dialogue entre deux camarades de classe à Yaoundé.

\begin{textebox}[Dialogue : À la sortie du lycée]
A：你是昨天到的吗？\\
\emph{Nǐ shì zuótiān dào de ma ?}\\
« C'est hier que tu es arrivé ? »\\[4pt]
B：不是，我是前天到的。\\
\emph{Bú shì, wǒ shì qiántiān dào de.}\\
« Non, c'est avant-hier que je suis arrivé. »\\[4pt]
A：你是怎么来的？\\
\emph{Nǐ shì zěnme lái de ?}\\
« Comment es-tu venu ? »\\[4pt]
B：我是坐出租车来的。我哥哥是开车来的。\\
\emph{Wǒ shì zuò chūzūchē lái de. Wǒ gēge shì kāichē lái de.}\\
« C'est en taxi que je suis venu. Mon grand frère, c'est en voiture qu'il est venu. »
\end{textebox}

Observons attentivement les phrases en gras de structure. Que remarque-t-on ?

\begin{itemize}
\item Chaque phrase contient \zh{是} au début et \zh{的} à la fin : les deux mots \cle{encadrent} la phrase.
\item L'action décrite (\zh{到} « arriver », \zh{来} « venir ») est \cle{déjà accomplie} : on parle du passé.
\item L'élément placé juste après \zh{是} est précisément celui sur lequel on insiste : le moment (\zh{昨天} « hier », \zh{前天} « avant-hier ») ou la manière (\zh{坐出租车} « en taxi », \zh{开车} « en voiture »).
\item Remarque importante : il n'y a \textbf{pas de} \zh{了} dans ces phrases, alors que l'action est passée ! La structure \zh{是……的} suffit à elle seule à situer l'action dans le passé.
\end{itemize}

\courssub{La règle : forme et emploi}

\begin{propriete}[La construction 是……的]
La construction \zh{是……的} sert à \cle{mettre en relief les circonstances} (le moment, le lieu, la manière, la personne avec qui) d'une action \textbf{déjà accomplie}. L'élément souligné se place \textbf{entre} \zh{是} et le verbe ; \zh{的} ferme la phrase.

\textbf{Schéma général :}

Sujet + 是 + (moment / lieu / manière) + verbe + 的

\textbf{Emplois :}
\begin{itemize}
\item Insister sur le \textbf{moment} : \zh{我是去年学的中文。} \emph{Wǒ shì qùnián xué de Zhōngwén.} « C'est l'année dernière que j'ai appris le chinois. »
\item Insister sur le \textbf{lieu} : \zh{我们是在学校学的汉语。} \emph{Wǒmen shì zài xuéxiào xué de Hànyǔ.} « C'est à l'école que nous avons appris le chinois. »
\item Insister sur la \textbf{manière} : \zh{他是坐出租车来的。} \emph{Tā shì zuò chūzūchē lái de.} « C'est en taxi qu'il est venu. »
\item Insister sur la \textbf{personne accompagnante} : \zh{她是跟妈妈去的杜阿拉。} \emph{Tā shì gēn māma qù de Dù'ālā.} « C'est avec sa maman qu'elle est allée à Douala. »
\end{itemize}

\textbf{Condition essentielle :} l'action doit être \cle{passée ou réalisée}. On n'utilise jamais \zh{是……的} pour le futur.
\end{propriete}

\begin{popfigure}
\begin{tikzpicture}[node distance=0.35cm]
\node[draw=popblue, fill=popblueL, rounded corners, minimum height=0.9cm, align=center] (suj) {\textbf{Sujet}\\ 我};
\node[draw=popdark, fill=popcream, rounded corners, minimum height=0.9cm, align=center, right=of suj] (shi) {\textbf{是}};
\node[draw=poporange, fill=poporangeL, rounded corners, minimum height=0.9cm, align=center, right=of shi, line width=1.5pt] (circ) {\textbf{Circonstance}\\ 昨天 / 在学校 / 坐出租车};
\node[draw=popgreen, fill=popgreenL, rounded corners, minimum height=0.9cm, align=center, right=of circ] (verb) {\textbf{Verbe}\\ 来 / 学 / 买};
\node[draw=popdark, fill=popcream, rounded corners, minimum height=0.9cm, align=center, right=of verb] (de) {\textbf{的}};
\node[draw=poppurple, dashed, thick, rounded corners, fit=(circ), inner sep=6pt, label={[poppurple]above:\textbf{élément mis en relief}}] {};
\draw[-{Stealth}, popmuted] (suj) -- (shi);
\draw[-{Stealth}, popmuted] (shi) -- (circ);
\draw[-{Stealth}, popmuted] (circ) -- (verb);
\draw[-{Stealth}, popmuted] (verb) -- (de);
\end{tikzpicture}
\legende{Schéma de la structure 是……的 : la circonstance mise en relief (encadrée) se place entre 是 et le verbe.}
\end{popfigure}

\courssub{Comparaison avec le français}

\begin{aretenir}[是……的 = « c'est… que… »]
Le français utilise la \textbf{phrase clivée} : « C'est \emph{hier} que je suis arrivé. » Le chinois, lui, \textbf{encadre} la phrase avec \zh{是}…\zh{的}, \textbf{sans inversion} et sans changer l'ordre des mots :

\zh{我 是 昨天 到 的。} = mot à mot : « moi — être — hier — arriver — (marque) ».
\end{aretenir}

\begin{center}
\begin{tabularx}{\linewidth}{|X|X|X|}\hline
\rowcolor{popblueL} \textbf{Français (clivée)} & \textbf{Chinois (caractères + pinyin)} & \textbf{Élément mis en relief} \\ \hline
C'est hier que je suis arrivé. & 我是昨天到的。 \emph{Wǒ shì zuótiān dào de.} & le moment \\ \hline
C'est à l'école que nous avons appris le chinois. & 我们是在学校学的汉语。 \emph{Wǒmen shì zài xuéxiào xué de Hànyǔ.} & le lieu \\ \hline
C'est en taxi qu'il est venu. & 他是坐出租车来的。 \emph{Tā shì zuò chūzūchē lái de.} & la manière \\ \hline
C'est avec son frère qu'elle est partie. & 她是跟哥哥走的。 \emph{Tā shì gēn gēge zǒu de.} & l'accompagnant \\ \hline
\end{tabularx}
\end{center}

\courssub{Exemples commentés}

\begin{exemplebox}[是……的 en contexte camerounais]
\zh{我是在雅温得出生的。}\\
\emph{Wǒ shì zài Yǎwēndé chūshēng de.}\\
« C'est à Yaoundé que je suis né. » (mise en relief du \textbf{lieu})\\[4pt]
\zh{她是去年去中国的。}\\
\emph{Tā shì qùnián qù Zhōngguó de.}\\
« C'est l'année dernière qu'elle est allée en Chine. » (mise en relief du \textbf{moment})\\[4pt]
\zh{我们是在杜阿拉机场告别的。}\\
\emph{Wǒmen shì zài Dù'ālā jīchǎng gàobié de.}\\
« C'est à l'aéroport de Douala que nous nous sommes dit au revoir. »\\[4pt]
\zh{这本书是他在市场上买的。}\\
\emph{Zhè běn shū shì tā zài shìchǎng shàng mǎi de.}\\
« C'est au marché qu'il a acheté ce livre. »
\end{exemplebox}

\courssub{La position de 的 quand il y a un complément d'objet}

Quand le verbe est suivi d'un complément d'objet, deux positions sont possibles pour \zh{的} :

\begin{propriete}[Position de 的 avec un objet]
\begin{itemize}
\item \textbf{En fin de phrase} (le plus courant) : \zh{我是昨天买的书。} \emph{Wǒ shì zuótiān mǎi de shū.} « C'est hier que j'ai acheté ce livre. »
\item \textbf{Juste après le verbe}, surtout si l'objet est long ou précédé d'un démonstratif : \zh{我是昨天买的这本书。} \emph{Wǒ shì zuótiān mǎi de zhè běn shū.}
\end{itemize}
Les deux phrases sont correctes ; à l'examen, la fin de phrase est la solution la plus sûre.
\end{propriete}

\courssub{Négation, question et omission de 是}

\begin{propriete}[Négation et interrogation]
\textbf{Négation :} on remplace \zh{是} par \zh{不是} ; le \zh{的} reste.

\zh{我不是昨天来的。} \emph{Wǒ bú shì zuótiān lái de.} « Ce n'est pas hier que je suis arrivé. »

\textbf{Question :} on insère le mot interrogatif à la place de la circonstance.

\zh{你是哪天到的？} \emph{Nǐ shì nǎ tiān dào de ?} « Quel jour es-tu arrivé ? »

\zh{她是怎么去中国的？} \emph{Tā shì zěnme qù Zhōngguó de ?} « Comment est-elle allée en Chine ? »

\textbf{Omission de 是 :} dans la langue parlée, à la forme affirmative, \zh{是} peut être omis : \zh{我昨天买的。} \emph{Wǒ zuótiān mǎi de.} « (C'est) hier (que) je l'ai acheté. » Le \zh{的}, lui, ne s'omet jamais.
\end{propriete}

\begin{center}
\begin{tabularx}{\linewidth}{|X|X|X|}\hline
\rowcolor{popblueL} \textbf{Forme} & \textbf{Exemple (caractères + pinyin)} & \textbf{Traduction} \\ \hline
Affirmative & 我是前天到的。 \emph{Wǒ shì qiántiān dào de.} & C'est avant-hier que je suis arrivé. \\ \hline
Négative & 我不是前天到的。 \emph{Wǒ bú shì qiántiān dào de.} & Ce n'est pas avant-hier que je suis arrivé. \\ \hline
Interrogative & 你是哪天到的？ \emph{Nǐ shì nǎ tiān dào de ?} & Quel jour es-tu arrivé ? \\ \hline
Orale (是 omis) & 我前天到的。 \emph{Wǒ qiántiān dào de.} & (C'est) avant-hier que je suis arrivé. \\ \hline
\end{tabularx}
\end{center}

\courssub{Exercice résolu et pièges à éviter}

\begin{exoresolu}[Transformer une phrase simple en mise en relief]
\textbf{Énoncé.} Transformez la phrase \zh{他去年去了中国。} \emph{Tā qùnián qù le Zhōngguó.} (« Il est allé en Chine l'année dernière. ») pour mettre en relief le moment « l'année dernière ».
\tcblower
\textbf{Solution détaillée.}
\begin{enumerate}
\item Identifier l'élément à mettre en relief : \zh{去年} (le moment).
\item Placer \zh{是} avant cet élément et \zh{的} en fin de phrase.
\item \textbf{Supprimer} \zh{了} : la structure \zh{是……的} marque déjà l'accompli, les deux ne se combinent pas.
\end{enumerate}
Résultat : \zh{他是去年去中国的。} \emph{Tā shì qùnián qù Zhōngguó de.} « C'est l'année dernière qu'il est allé en Chine. »
\end{exoresolu}

\begin{attention}[Ne jamais oublier le 的 final]
L'erreur la plus fréquente des francophones est d'écrire * \zh{我是昨天来} en oubliant le \zh{的} final. En chinois, \zh{是} et \zh{的} fonctionnent \textbf{en couple} : ils encadrent la phrase comme des parenthèses. Sans \zh{的}, la phrase est incorrecte ou change de sens. Retenez : \textbf{有是必有的} — « s'il y a 是 (de mise en relief), il doit y avoir 的 ».
\end{attention}

\begin{attention}[Interdit pour le futur]
\zh{是……的} ne s'emploie \textbf{que pour une action passée}. La phrase * \zh{我是明天去的} est incorrecte. Pour le futur, dites simplement : \zh{我明天去。} \emph{Wǒ míngtiān qù.} « J'y vais demain. » De même, ne combinez jamais \zh{是……的} avec \zh{了} : * \zh{我是昨天来了的} est fautif.
\end{attention}

\courssub{Vocabulaire clé de la section}

\begin{center}
\begin{tabularx}{\linewidth}{|X|X|X|X|}\hline
\rowcolor{popblueL} \textbf{Caractères} & \textbf{Pinyin} & \textbf{Nature} & \textbf{Traduction} \\ \hline
出生 & chūshēng & verbe & naître \\ \hline
告别 & gàobié & verbe & dire au revoir, quitter \\ \hline
市场 & shìchǎng & nom & marché \\ \hline
雅温得 & Yǎwēndé & nom propre & Yaoundé \\ \hline
杜阿拉 & Dù'ālā & nom propre & Douala \\ \hline
布埃亚 & Bù'āiyà & nom propre & Buea \\ \hline
公共汽车 & gōnggòng qìchē & nom & bus \\ \hline
出租车 & chūzūchē & nom & taxi \\ \hline
\end{tabularx}
\end{center}

\begin{exercice}[À vous de jouer]
\begin{enumerate}
\item Traduisez en chinois avec \zh{是……的} : « C'est en bus que nous sommes allés à Buea. »
\item Mettez à la forme négative : \zh{她是在学校学的中文。}
\item Posez la question sur le lieu : \zh{我是在雅温得买的这本书。}
\item Corrigez la phrase fautive : * \zh{他是明天来中国的。}
\end{enumerate}
\end{exercice}

\begin{corrige}[Exercice « À vous de jouer »]
\begin{enumerate}
\item \zh{我们是坐公共汽车去布埃亚的。} \emph{Wǒmen shì zuò gōnggòng qìchē qù Bù'āiyà de.}
\item \zh{她不是在学校学的中文。} \emph{Tā bú shì zài xuéxiào xué de Zhōngwén.} « Ce n'est pas à l'école qu'elle a appris le chinois. »
\item \zh{你是在哪儿买的这本书？} \emph{Nǐ shì zài nǎr mǎi de zhè běn shū ?} « Où as-tu acheté ce livre ? »
\item \zh{是……的} est impossible pour le futur. Correction : \zh{他明天来中国。} \emph{Tā míngtiān lái Zhōngguó.} « Il vient en Chine demain. »
\end{enumerate}
\end{corrige}

\begin{aretenir}[L'essentiel sur 是……的]
\begin{itemize}
\item Schéma : Sujet + 是 + circonstance + verbe + 的.
\item Emploi : insister sur le moment, le lieu, la manière ou l'accompagnant d'une action \textbf{passée}.
\item Négation : \zh{不是……的} ; question : mot interrogatif à la place de la circonstance.
\item \zh{是} peut être omis à l'oral, jamais \zh{的}.
\item Jamais de \zh{了} dans la structure, jamais pour le futur.
\item Équivalent français : la phrase clivée « c'est… que… », mais sans inversion en chinois.
\end{itemize}
\end{aretenir}

\coursec{Synthèse, comparaison et entraînement}

Dans les sections précédentes, nous avons étudié séparément l'hypothèse (\zh{如果……，那么/就……}), les nombres, la répartition (\zh{有的……有的……}) et la mise en relief (\zh{是……的}). Nous les réunissons ici pour en comparer le fonctionnement, éviter les interférences du français et nous entraîner à les mobiliser ensemble, comme le requiert l'épreuve du baccalauréat.

\courssub{Observation guidée : les quatre structures en contexte}

Avant d'analyser, lisez le texte ci-dessous. Repérez visuellement (surlignez si possible) : 1. la phrase d'hypothèse, 2. le nombre composé, 3. la série \zh{有的……}, 4. la phrase \zh{是……的}.

\begin{textebox}[Dimanche au parc — 星期天在公园]
如果星期天天气好，我们就去公园。\\
\emph{Rúguǒ xīngqītiān tiānqì hǎo, wǒmen jiù qù gōngyuán.}\\
公园里有一百多个人：有的跳舞，有的跑步，还有的唱歌。\\
\emph{Gōngyuán lǐ yǒu yì bǎi duō gè rén : yǒude tiàowǔ, yǒude pǎobù, hái yǒude chànggē.}\\
我们是早上九点到的。\\
\emph{Wǒmen shì zǎoshang jiǔ diǎn dào de.}
\end{textebox}

Traduction : « S'il fait beau dimanche, nous irons au parc. Il y a plus de cent personnes au parc : certains dansent, d'autres courent, d'autres encore chantent. C'est à neuf heures du matin que nous sommes arrivés. »

\begin{exemplebox}[Repérage commenté]
\begin{itemize}
\item \zh{如果星期天天气好，我们就去公园。} $\rightarrow$ Hypothèse : \zh{如果} + condition (sujet + verbe), \zh{就} + conséquence (sujet + verbe). \zh{那么} est une variante plus formelle de \zh{就}.
\item \zh{一百多个人} $\rightarrow$ Nombre \zh{一百} (100) + \zh{多} (plus de) + classificateur \zh{个} + nom. Structure : [Nombre + 多 + Classif. + Nom] = « plus de [Nombre]… ».
\item \zh{有的跳舞，有的跑步，还有的唱歌。} $\rightarrow$ Répartition en trois groupes. \zh{有的} (certains) se répète devant chaque activité. \zh{还有的} marque l'ajout d'un dernier groupe (« d'autres encore »).
\item \zh{我们是早上九点到的。} $\rightarrow$ Mise en relief de la circonstance de temps (\zh{早上九点}) d'une action passée (\zh{到}). Structure : Sujet + \zh{是} + Circonstance + Verbe + \zh{的}.
\end{itemize}
\end{exemplebox}

\courssub{Lexique de synthèse : mots-outils et nombres clés}

\begin{center}
\begin{tabularx}{\linewidth}{|X|X|X|X|}
\hline
\rowcolor{popblueL} Caractères & Pinyin & Nature & Traduction \\
\hline
如果 & rúguǒ & conjonction & si (introduit l'hypothèse) \\
\hline
那么 & nàme & adverbe & alors, dans ce cas (formel) \\
\hline
就 & jiù & adverbe & alors, donc (conséquence, courant) \\
\hline
有的 & yǒude & pronom indéfini & certains, les uns \\
\hline
还有的 & hái yǒude & pronom indéfini & d'autres encore \\
\hline
是 & shì & verbe / outil & c'est… (mise en relief) \\
\hline
的 & de & particule & marque la fin de la mise en relief (\zh{是……的}) \\
\hline
零 & líng & nombre & zéro (zéro interne obligatoire) \\
\hline
两 & liǎng & nombre & deux (devant classificateur) \\
\hline
多 & duō & adjectif / suffixe & plus de (après un nombre rond) \\
\hline
\end{tabularx}
\end{center}

\courssub{Tableau récapitulatif des quatre structures}

\begin{center}
\begin{tabularx}{\linewidth}{|X|X|X|X|}
\hline
\rowcolor{popblueL} Structure & Schéma & Fonction & Exemple (caractères + pinyin + traduction) \\
\hline
\zh{如果……，就/那么……} & \zh{如果} + condition，(\zh{那么})\zh{就} + conséquence & Exprimer une hypothèse et sa conséquence & \zh{如果明天下雨，我就不去市场。} \emph{Rúguǒ míngtiān xiàyǔ, wǒ jiù bú qù shìchǎng.} « S'il pleut demain, je n'irai pas au marché. » \\
\hline
Les nombres & nombre + classif. + nom ; \zh{零} pour 0 interne ; \zh{两} devant classif. & Compter, quantifier, dater, donner l'heure, le prix & \zh{我们学校有两千零五十个学生。} \emph{Wǒmen xuéxiào yǒu liǎng qiān líng wǔshí gè xuésheng.} « Notre école a deux mille cinquante élèves. » \\
\hline
\zh{有的……有的……} & \zh{有的} + V/Adj，\zh{有的} + V/Adj & Décrire la répartition d'un groupe (« les uns… les autres… ») & \zh{同学们有的看书，有的写字。} \emph{Tóngxuémen yǒude kànshū, yǒude xiězì.} « Les élèves, les uns lisent, les autres écrivent. » \\
\hline
\zh{是……的} & Sujet + \zh{是} + circonstance + verbe + \zh{的} & Mettre en relief quand/où/comment/avec qui une action \textbf{passée} a eu lieu & \zh{他是去年去中国的。} \emph{Tā shì qùnián qù Zhōngguó de.} « C'est l'année dernière qu'il est allé en Chine. » \\
\hline
\end{tabularx}
\end{center}

\begin{popfigure}
\begin{tikzpicture}[
box/.style={draw=popblue, fill=popblueL, rounded corners=2pt, align=center, font=\small, inner sep=4pt},
sch/.style={draw=poporange, fill=poporangeL, rounded corners=2pt, align=center, font=\small, inner sep=4pt},
fct/.style={draw=popgreen, fill=popgreenL, rounded corners=2pt, align=center, font=\small, inner sep=4pt},
ex/.style={draw=poppurple, fill=poppurpleL, rounded corners=2pt, align=center, font=\small, inner sep=4pt},
head/.style={fill=popblue, text=white, rounded corners=2pt, align=center, font=\small\bfseries, inner sep=4pt}]
\node[head, minimum width=2.6cm] (h1) at (0,0) {Structure};
\node[head, minimum width=3.4cm] (h2) at (3.3,0) {Schéma};
\node[head, minimum width=2.8cm] (h3) at (6.6,0) {Fonction};
\node[head, minimum width=4.6cm] (h4) at (10.5,0) {Exemple traduit};
\node[box, minimum width=2.6cm, align=center] (a1) at (0,-1.5){如果……\\就/那么……};
\node[sch, minimum width=3.4cm, align=center] (a2) at (3.3,-1.5){如果 + condition\\就/那么 + conséquence};
\node[fct, minimum width=2.8cm] (a3) at (6.6,-1.5) {Hypothèse};
\node[ex, minimum width=4.6cm, align=center] (a4) at (10.5,-1.5){如果天气好，我们就去公园。\\« S'il fait beau, nous irons au parc. »};
\node[box, minimum width=2.6cm] (b1) at (0,-3.1) {Les nombres};
\node[sch, minimum width=3.4cm, align=center] (b2) at (3.3,-3.1){Nombre + classif. + nom\\零 = 0 interne ; 两 + classif.};
\node[fct, minimum width=2.8cm, align=center] (b3) at (6.6,-3.1){Quantité, prix,\\heure, date};
\node[ex, minimum width=4.6cm, align=center] (b4) at (10.5,-3.1){两百零三个学生\\« deux cent trois élèves »};
\node[box, minimum width=2.6cm, align=center] (c1) at (0,-4.7){有的……\\有的……};
\node[sch, minimum width=3.4cm, align=center] (c2) at (3.3,-4.7){有的 + V/Adj，\\有的 + V/Adj};
\node[fct, minimum width=2.8cm, align=center] (c3) at (6.6,-4.7){Répartition\\d'un groupe};
\node[ex, minimum width=4.6cm, align=center] (c4) at (10.5,-4.7){有的跳舞，有的跑步。\\« Les uns dansent, les autres courent. »};
\node[box, minimum width=2.6cm] (d1) at (0,-6.3) {是……的};
\node[sch, minimum width=3.4cm, align=center] (d2) at (3.3,-6.3){是 + circonstance\\+ verbe + 的};
\node[fct, minimum width=2.8cm, align=center] (d3) at (6.6,-6.3){Relief action\\passée};
\node[ex, minimum width=4.6cm, align=center] (d4) at (10.5,-6.3){我们是早上九点到的。\\« C'est à 9 h que nous sommes arrivés. »};
\end{tikzpicture}
\legende{Les quatre structures de la leçon : schéma, fonction et exemple}
\end{popfigure}

\courssub{Comparaison globale avec le français}

Le français possède des équivalents pour chacune de ces structures, mais le fonctionnement chinois est plus régulier (pas de conjugaison, pas d'accord). Comparons systématiquement :

\begin{center}
\begin{tabularx}{\linewidth}{|X|X|X|}
\hline
\rowcolor{popblueL} Notion & Français & Chinois \\
\hline
Hypothèse & « Si + indicatif imparfait/présent $\rightarrow$ conditionnel/présent/futur » (accords de temps complexes) & \zh{如果} + condition，\zh{就} + conséquence : \textbf{aucun temps verbal}, aucune conjugaison, seul l'ordre des mots et les mots-outils comptent. \\
\hline
Nombres & « Deux cents », « mille un » ; zéro interne souvent muet (« deux cent trois ») & Chaque rang énoncé explicitement (\zh{两百零三}) ; \textbf{zéro interne obligatoire} (\zh{零}) ; \zh{两} (et non \zh{二}) devant classificateur. \\
\hline
Répartition & « Les uns… les autres… », « Certains… d'autres… » (accords, articles) & \zh{有的……有的……} répété devant chaque prédicat, \textbf{sans article, sans accord}, neutre. \\
\hline
Mise en relief & Phrase clivée « C'est… que/qui… » (ordre figé) & \zh{是……的} : \zh{是} avant l'élément souligné, \zh{的} \textbf{après le verbe} (action accomplie). \\
\hline
\end{tabularx}
\end{center}

\begin{attention}[Les trois pièges majeurs — récapitulatif]
\begin{enumerate}
\item \cle{就 mal placé} : \zh{就} se place \textbf{immédiatement devant le verbe} (ou l'auxiliaire/négation) de la conséquence, jamais en début de proposition, jamais devant le sujet. \checkmark \zh{如果下雨，我就不去。} \quad \textcolor{red}{\ding{55}} \zh{如果下雨，就我不去。}
\item \cle{的 final oublié dans \zh{是……的}} : Le \zh{的} final est \textbf{obligatoire} pour clore la construction de mise en relief d'une action passée. \checkmark \zh{我是昨天来的。} \quad \textcolor{red}{\ding{55}} \zh{我是昨天来。} (phrase incomplète/incorrecte).
\item \cle{二 (\zh{èr}) et 两 (\zh{liǎng}) confondus} : Devant un classificateur (et pour « deux cents/mille » ronds), on emploie \zh{两} : \zh{两个人}, \zh{两本书}, \zh{两百人}. \zh{二} s'emploie dans le comptage pur, les numéros, dates, étages, et nombres composés \textbf{non suivis de classificateur} : \zh{二月} (février), \zh{二十二} (22), \zh{二楼} (2e étage), \zh{一百二} (120 — oral).
\end{enumerate}
\end{attention}

\begin{methode}[Choisir la bonne structure en production]
Face à une phrase à traduire (thème) ou à transformer, posez-vous ces quatre questions \textbf{dans l'ordre} :
\begin{enumerate}
\item La phrase exprime-t-elle une \textbf{condition / hypothèse} (« Si… ») ? $\rightarrow$ Construire : \zh{如果} + condition，\zh{就} + conséquence.
\item Décrit-elle un \textbf{groupe divisé en sous-ensembles} agissant différemment (« Les uns… les autres… ») ? $\rightarrow$ Construire : \zh{有的} + action 1，\zh{有的} + action 2 (+\zh{还有的} + action 3).
\item Souligne-t-elle une \textbf{circonstance} (temps, lieu, moyen, accompagnateur) d'une \textbf{action déjà accomplie} ? $\rightarrow$ Construire : Sujet + \zh{是} + circonstance + verbe + \zh{的}.
\item Donne-t-elle une \textbf{donnée chiffrée} (quantité, prix, date, heure, âge) ? $\rightarrow$ Écrire le nombre en toutes lettres (chiffres acceptés en écriture rapide) avec \zh{零} pour zéro interne et \zh{两} devant le classificateur.
\end{enumerate}
\end{methode}

\begin{savaistu}[Ces structures au baccalauréat LV2]
L'épreuve écrite de chinois (Première A) teste très fréquemment la \textbf{transformation de phrases} :
\begin{itemize}
\item Transformer une phrase simple en phrase \zh{是……的} (relief temps/lieu/moyen).
\item Fusionner deux phrases parallèles en une phrase \zh{有的……有的……}.
\item Compléter une phrase hypothétique avec \zh{就} ou \zh{那么}.
\item Écrire un nombre complexe en lettres (avec \zh{零} et \zh{两}).
\end{itemize}
Maîtriser la mécanique de ces quatre structures, c'est sécuriser 15 à 20\% des points de grammaire. Entraînez-vous à les identifier \textbf{au premier coup d'œil}.

\end{savaistu}

\courssub{Exercices d'application : complétion}

\begin{exoresolu}[Compléter avec le bon mot-outil]
Énoncé : Complétez les phrases avec \zh{就}, \zh{那么}, \zh{有的}, \zh{是}, \zh{的}, \zh{零} ou \zh{两}.
\begin{enumerate}
\item 如果明天有空，我\textcolor{white}{（　）}去看你。
\item 我们班有三百\textcolor{white}{（　）}五个学生。
\item 操场上，\textcolor{white}{（　）}同学踢球，\textcolor{white}{（　）}同学跑步。
\item 我\textcolor{white}{（　）}在杜阿拉出生\textcolor{white}{（　）}。
\item 我有\textcolor{white}{（　）}个哥哥。
\end{enumerate}
\tcblower
Solution détaillée :
\begin{enumerate}
\item \zh{如果明天有空，我就去看你。} — \zh{就} placé devant le verbe \zh{看} marque la conséquence.
\item \zh{我们班有三百零五个学生。} — Zéro interne (pas de dizaines) : \zh{零} obligatoire.
\item \zh{操场上，有的同学踢球，有的同学跑步。} — Répartition : \zh{有的} répété devant chaque groupe-sujet.
\item \zh{我是在杜阿拉出生的。} — Mise en relief du lieu : \zh{是} + lieu + verbe + \zh{的}.
\item \zh{我有两个哥哥。} — Devant le classificateur \zh{个}, on utilise \zh{两} (et non \zh{二}).
\end{enumerate}
\end{exoresolu}

\begin{exercice}[À toi de compléter]
Complétez avec \zh{就}, \zh{那么}, \zh{有的}, \zh{是}, \zh{的}, \zh{零} ou \zh{两} :
\begin{enumerate}
\item 如果太贵，我\textcolor{white}{（　）}不买。
\item 雅温得有一千\textcolor{white}{（　）}八个学生参加比赛。
\item 市场里，\textcolor{white}{（　）}人买鱼，\textcolor{white}{（　）}人买水果。
\item 他\textcolor{white}{（　）}坐公共汽车来学校\textcolor{white}{（　）}。
\item 我家有\textcolor{white}{（　）}只狗。
\end{enumerate}
\end{exercice}

\begin{corrige}[Exercice « À toi de compléter »]
\begin{enumerate}
\item \zh{如果太贵，我就不买。} (\zh{juste devant le verbe \zh{不买}}.)
\item \zh{雅温得有一千零八个学生参加比赛。} (Zéro interne centaines $\rightarrow$ unités : \zh{零}.)
\item \zh{市场里，有的人买鱼，有的人买水果。} (Répartition : \zh{有的……有的……} ; \zh{人} ajouté pour la fluidité).
\item \zh{他是坐公共汽车来学校的。} (Mise en relief du moyen : \zh{是} + moyen + verbe + \zh{的}.)
\item \zh{我家有两只狗。} (\zh{两} devant le classificateur \zh{只}.)
\end{enumerate}
\end{corrige}

\courssub{Exercices de transformation}

\begin{exoresolu}[Transformer les phrases selon le modèle]
Énoncé :
\begin{enumerate}
\item Transforme en phrase avec \zh{是……的} (mets en relief l'élément souligné — ici l'heure) : \zh{我们早上九点到了。}
\item Réunis en une phrase avec \zh{有的……有的……} : \zh{一些同学唱歌。一些同学跳舞。}
\end{enumerate}
\tcblower
Solution détaillée :
\begin{enumerate}
\item On place \zh{是} avant l'élément à souligner (\zh{早上九点}) et \zh{的} après le verbe : \zh{我们是早上九点到的。} \emph{Wǒmen shì zǎoshang jiǔ diǎn dào de.} « C'est à neuf heures du matin que nous sommes arrivés. »
\item On remplace « 一些 » par \zh{有的} devant chaque groupe, sujet commun \zh{同学们} en tête : \zh{同学们有的唱歌，有的跳舞。} \emph{Tóngxuémen yǒude chànggē, yǒude tiàowǔ.} « Les élèves, les uns chantent, les autres dansent. »
\end{enumerate}
\end{exoresolu}

\begin{exercice}[Transforme à ton tour]
\begin{enumerate}
\item Mets en relief le \textbf{lieu} avec \zh{是……的} : \zh{我在北京学习中文。}
\item Mets en relief l'\textbf{accompagnateur} avec \zh{是……的} : \zh{我跟妈妈去的市场。} (Déjà souligné $\rightarrow$ reconstruis la phrase source simple puis transforme).
\item Réunis : \zh{一些学生喜欢足球。一些学生喜欢篮球。}
\item Réunis : \zh{一些人喝茶。一些人喝咖啡。}
\end{enumerate}
\end{exercice}

\begin{corrige}[Exercice « Transforme à ton tour »]
\begin{enumerate}
\item \zh{我是在北京学习中文的。} \emph{Wǒ shì zài Běijīng xuéxí Zhōngwén de.} « C'est à Pékin que j'apprends le chinois. »
\item Phrase source : \zh{我跟妈妈去市场。} $\rightarrow$ Relief : \zh{我是跟妈妈去市场的。} \emph{Wǒ shì gēn māma qù shìchǎng de.} « C'est avec maman que je vais/ suis allé au marché. »
\item \zh{学生们有的喜欢足球，有的喜欢篮球。} \emph{Xuéshengmen yǒude xǐhuan zúqiú, yǒude xǐhuan lánqiú.}
\item \zh{有的人喝茶，有的人喝咖啡。} \emph{Yǒude rén hē chá, yǒude rén hē kāfēi.}
\end{enumerate}
\end{corrige}

\courssub{Traduction : thème et version}

\begin{exercice}[Thème : français $\rightarrow$ chinois]
Traduis en chinois (caractères + pinyin) :
\begin{enumerate}
\item S'il fait beau demain, nous irons au stade.
\item Il y a deux cent cinq élèves dans notre école.
\item Au marché, les uns vendent des légumes, les autres vendent des fruits.
\item C'est en 2023 que je suis arrivé à Yaoundé.
\end{enumerate}
\end{exercice}

\begin{corrige}[Thème]
\begin{enumerate}
\item \zh{如果明天天气好，我们就去体育场。} \emph{Rúguǒ míngtiān tiānqì hǎo, wǒmen jiù qù tǐyùchǎng.}
\item \zh{我们学校有两百零五个学生。} \emph{Wǒmen xuéxiào yǒu liǎng bǎi líng wǔ gè xuésheng.}
\item \zh{市场里，有的卖蔬菜，有的卖水果。} \emph{Shìchǎng lǐ, yǒude mài shūcài, yǒude mài shuǐguǒ.} (Sujet \zh{人} omis, récupéré du contexte).
\item \zh{我是二零二三年到雅温得的。} \emph{Wǒ shì èr líng èr sān nián dào Yǎwēndé de.} (Année lue chiffre par chiffre).
\end{enumerate}
\end{corrige}

\begin{exercice}[Version : chinois $\rightarrow$ français]
Traduis en français :
\begin{enumerate}
\item 如果明天下雨，我们就不去杜阿拉。
\item 这个学校有三千零二十个学生。
\item 公园里，有的老人下棋，有的老人跳舞。
\item 她是坐飞机去中国的。
\end{enumerate}
\end{exercice}

\begin{corrige}[Version]
\begin{enumerate}
\item « S'il pleut demain, nous n'irons pas à Douala. »
\item « Cette école compte trois mille vingt élèves. »
\item « Au parc, certains personnes âgées jouent aux échecs, d'autres dansent. »
\item « C'est en avion qu'elle est allée en Chine. »
\end{enumerate}
\end{corrige}

\courssub{Production guidée : « Dimanche au quartier »}

\begin{experience}[Rédige ton paragraphe — Expression écrite]
Rédige un court paragraphe (5--7 phrases) intitulé \zh{星期天在我们的街区} (« Dimanche dans notre quartier ») en utilisant \textbf{au moins trois des quatre structures} de la leçon. Respecte la ponctuation chinoise (， 。)。

\textbf{Plan suggéré :}
\begin{enumerate}
\item Une hypothèse météo : \zh{如果星期天天气好，……就……}
\item Une quantité (habitants, enfants, maisons) : pense à \zh{零} et \zh{两}.
\item Une répartition d'activités : \zh{有的……有的……还有的……}
\item Une mise en relief (heure de sortie, moyen de transport, compagnie) : \zh{是……的}
\end{enumerate}

\textbf{Modèle d'amorce (ne le recopie pas, inspire-toi en) :}
\zh{如果星期天天气好，我们一家人就出去。街区里有两百多个人：有的聊天，有的打球，还有的买东西。我们是下午三点出门的。}
\end{experience}

\courssub{Auto-évaluation}

\begin{aretenir}[Grille d'auto-évaluation — Leçon 6]
Relis ta production et coche chaque point validé :
\begin{itemize}
\item \square J'ai placé \zh{就} \textbf{juste devant le verbe} de la conséquence (pas devant le sujet).
\item \square Dans \zh{是……的}, je n'ai pas oublié le \zh{的} \textbf{final}, après le verbe.
\item \square Dans mes nombres, j'ai écrit \zh{零} pour chaque \textbf{zéro interne} (ex. \zh{三百零五}, \zh{一千零八}).
\item \square Devant un classificateur, j'ai utilisé \zh{两} et non \zh{二} (ex. \zh{两个人}, \zh{两只狗}).
\item \square Dans \zh{有的……有的……}, j'ai bien \textbf{répété} \zh{有的} devant chaque groupe.
\item \square Ma ponctuation est chinoise : ， 。 ： ； et non , . : ;.
\end{itemize}
\textbf{Bilan :} 6/6 $\rightarrow$ Maîtrise acquise. $<$ 6 $\rightarrow$ Revois la fiche de la structure concernée et refais les exercices de transformation correspondants.
\end{aretenir}

\newpage

\coursec{Activité d'intégration}

\coursec{Leçon 6 : Intégration des structures — 如果……就…… \textperiodcentered\ Nombres \textperiodcentered\ 有的……有的…… \textperiodcentered\ 是……的}

\courssub{Objectifs de la leçon}
\begin{itemize}
    \item Maîtriser l'emploi combiné des quatre structures cibles : hypothèse conditionnelle (\zh{如果……，就/那么……}), numération (cardinaux, ordinaux, prix en FCFA), répartition (\zh{有的……，有的……}) et mise en relief du passé (\zh{是……的}).
    \item Produire un texte narratif court (6--8 phrases) cohérent, ancré dans le contexte camerounais (Yaoundé : marché Mokolo, stade Ahidjo, transport, météo).
    \item Passer de l'écrit à l'oral : lecture expressive avec tons corrects, écoute sélective pour identification grammaticale par les pairs.
    \item Développer l'auto-évaluation et l'esprit critique via un vote argumenté sur la richesse lexicale (HSK 2--3) et la correction syntaxique.
\end{itemize}

\courssub{Observation et mise en situation}
\begin{textebox}[Situation de communication : Un dimanche à Yaoundé]
C'est dimanche matin, 9 h 30. Le ciel est chargé au-dessus de \zh{雅温得} (Yǎwēndé, Yaoundé). Vous êtes en groupe de trois élèves. Votre mission : rédiger ensemble le récit de cette journée type, depuis le lever jusqu'au retour à la maison, en imaginant les aléas de la météo, l'ambiance au \zh{莫科洛市场} (Mòkēluò Shìchǎng, Marché Mokolo) ou au \zh{阿马杜·阿希乔体育场} (Āmǎdù·Āxījū Tǐyùchǎng, Stade Omnisports), le coût de la vie et le moyen de rentrer.

\textbf{Contraintes obligatoires (une par structure étudiée) :}
\begin{enumerate}
    \item \textbf{Hypothèse météo :} Une phrase avec \zh{如果……，就/那么……} (ex. : s'il pleut, on prend le taxi ; s'il fait beau, on marche).
    \item \textbf{Nombres \& Prix :} Au moins trois nombres différents (personnes, heure, prix en FCFA, quantité d'achats).
    \item \textbf{Répartition :} Une phrase (ou deux) avec \zh{有的……，有的……} décrivant une foule hétérogène (acheteurs, supporters, vendeurs, piétons).
    \item \textbf{Mise en relief :} Une phrase avec \zh{是……的} pour insister sur \emph{le moment}, \emph{le lieu}, \emph{le moyen} ou \emph{la manière} d'une action passée (ex. : c'est \textbf{en taxi} que nous sommes rentrés / c'est \textbf{à 14 h} que nous sommes arrivés).
\end{enumerate}
Le texte final fera 6 à 8 phrases. Il sera lu à haute voix devant la classe.
\end{textebox}

\courssub{Rappels grammaticaux et lexicaux}

\begin{propriete}[Structure 1 : 如果……，就/那么…… (Hypothèse conditionnelle réaliste)]
\textbf{Schéma :} \zh{如果} + Condition (Sujet + Verbe/Adj) + \zh{，} + \zh{就 / 那么} + Conséquence (Sujet + Verbe). \\
\textbf{Emploi :} Exprime une condition réaliste et sa conséquence logique. \zh{就} (jiù) est courant à l'oral ; \zh{那么} (nàme) plus formel/écrit. Le sujet de la conséquence s'omet s'il est identique à celui de la condition. \\
\textbf{Exemples :} \\
\zh{如果下雨，我们就打车去市场。} (Rúguǒ xiàyǔ, wǒmen jiù dǎchē qù shìchǎng.) — S'il pleut, nous prendrons un taxi pour aller au marché. \\
\zh{如果不下雨，那么我们走路去。} (Rúguǒ bù xiàyǔ, nàme wǒmen zǒulù qù.) — S'il ne pleut pas, alors nous irons à pied. \\
\zh{如果你累了，就休息一下吧。} (Rúguǒ nǐ lèi le, jiù xiūxi yíxià ba.) — Si tu es fatigué, repose-toi un peu.
\end{propriete}

\begin{propriete}[Structure 2 : Les nombres (Cardinaux, Classificateurs, Prix en FCFA)]
\textbf{Unités clés :} \zh{十} (shí, 10), \zh{百} (bǎi, 100), \zh{千} (qiān, 1\,000), \zh{万} (wàn, 10\,000). \textbf{Attention :} 10\,000 = \zh{一万} (yī wàn), jamais \zh{十千}. \\
\textbf{Règle du \zh{两} (liǎng) :} On utilise \zh{两} pour 200 (\zh{两百}), 2\,000 (\zh{两千}), 20\,000 (\zh{两万}) et devant un classificateur (\zh{两个人}, \zh{两斤}). \zh{二} (èr) s'emploie pour le chiffre 2 pur, les numéros (tél., étage) et les décimales. \\
\textbf{Classificateurs obligatoires (sélection HSK 1--3) :} \\
\begin{center}
\begin{tabularx}{\linewidth}{|X|X|X|}
\hline
\rowcolor{popblueL} \textbf{Classif.} & \textbf{Pinyin} & \textbf{Usage principal} \\
\hline
\zh{个} & gè & Général (personnes, objets) \\
\zh{位} & wèi & Personnes (registre poli) \\
\zh{斤} & jīn & Poids (500 g), fruits/légumes \\
\zh{张} & zhāng & Objets plats (billets, tables, lits) \\
\zh{辆} & liàng & Véhicules (voiture, moto, vélo) \\
\hline
\end{tabularx}
\end{center}
\textbf{Prix au Cameroun (FCFA) :} On dit le nombre + \zh{法郎} (Fǎláng). Pas de \zh{元} (yuán) ni \zh{块} (kuài) qui sont pour le RMB. \\
\textbf{Exemples :} \\
\zh{十五个人} (shíwǔ gè rén) — 15 personnes. \\
\zh{两万五千法郎} (liǎng wàn wǔ qiān Fǎláng) — 25\,000 FCFA. \\
\zh{三斤芒果} (sān jīn mángguǒ) — 3 catties (1,5 kg) de mangues. \\
\zh{第八辆摩托车} (dì bā liàng mótuōchē) — La 8\textsuperscript{e} moto-taxi (ordinaux : \zh{第} + nombre + classif.).
\end{propriete}

\begin{propriete}[Structure 3 : 有的……，有的…… (Répartition : « Certains…, d'autres… »)]
\textbf{Schéma :} \zh{有的} + Sujet/Ø + Prédicat (V/Adj) + \zh{，} + \zh{有的} + Sujet/Ø + Prédicat (V/Adj). \\
\textbf{Emploi :} Décrit la diversité au sein d'un groupe. \zh{有的} = « Certains (d'entre eux) ». Le sujet de la 2\textsuperscript{e} partie s'omet s'il est du même type. On peut ajouter \zh{还有的} pour un 3\textsuperscript{e} groupe. \\
\textbf{Exemples :} \\
\zh{市场里人很多，有的买菜，有的卖水果。} (Shìchǎng lǐ rén hěn duō, yǒude mǎi cài, yǒude mài shuǐguǒ.) — Au marché, certains achètent des légumes, d'autres vendent des fruits. \\
\zh{看球的人有的穿绿衣服，有的穿黄衣服，还有的打旗子。} (Kàn qiú de rén yǒude chuān lǜ yīfu, yǒude chuān huáng yīfu, hái yǒude dǎ qízi.) — Spectateurs : certains en vert, d'autres en jaune, d'autres agitent des drapeaux.
\end{propriete}

\begin{propriete}[Structure 4 : 是……的 (Mise en relief : Temps / Lieu / Moyen / Manière d'action PASSÉE)]
\textbf{Schéma :} \zh{是} + \textbf{Élément mis en relief (Temps / Lieu / Moyen / Manière)} + \zh{的}. Le verbe principal reste \textbf{avant} \zh{的} (souvent répété pour la clarté) ou s'omet si connu. \textbf{Négation :} \zh{不是……的}. \\
\textbf{Règle d'or :} L'action est \textbf{déjà accomplie}. \textbf{Interdiction formelle de \zh{了} après le verbe} dans cette structure. \\
\textbf{Ordre des compléments (rappel) :} Sujet + \zh{是} + [Temps / Lieu / Moyen / Manière] + Verbe + \zh{的}. \\
\textbf{Exemples :} \\
\zh{我是坐摩托车回来的。} (Wǒ shì zuò mótuōchē huílái de.) — C'est \textbf{en moto-taxi} que je suis revenu (Moyen). \\
\zh{我们是十四点到的体育场。} (Wǒmen shì shísì diǎn dào de tǐyùchǎng.) — C'est \textbf{à 14 h} que nous sommes arrivés au stade (Temps). \\
\zh{他是在莫科洛市场买的芒果。} (Tā shì zài Mòkēluò Shìchǎng mǎi de mángguǒ.) — C'est \textbf{au marché Mokolo} qu'il a acheté les mangues (Lieu). \\
\zh{她是走着回去的。} (Tā shì zǒuzhe huíqù de.) — C'est \textbf{à pied} qu'elle est rentrée (Manière).
\end{propriete}

\begin{definition}[Lexique utile : « Dimanche à Yaoundé » (HSK 2--3)]
\begin{center}
\begin{tabularx}{\linewidth}{|X|X|X|X|}
\hline
\rowcolor{popblueL} \textbf{Caractères} & \textbf{Pinyin} & \textbf{Nature} & \textbf{Traduction} \\
\hline
雅温得 & Yǎwēndé & n. & Yaoundé \\
莫科洛市场 & Mòkēluò Shìchǎng & n. & Marché Mokolo \\
阿马杜·阿希乔体育场 & Āmǎdù·Āxījū Tǐyùchǎng & n. & Stade Omnisports (Ahidjo) \\
摩托车 / 奥凯 & mótuōchē / àokǎi & n. & Moto-taxi (terme courant : « okada » / 奥凯) \\
出租车 & chūzūchē & n. & Taxi (voiture) \\
下雨 & xiàyǔ & v. & Pleuvoir \\
天气 / 天气预报 & tiānqì / tiānqì yùbào & n. & Météo / Bulletin météo \\
看球 & kàn qiú & v.o. & Regarder un match (football) \\
买菜 & mǎi cài & v.o. & Faire les courses (légumes/viande) \\
芒果 / 香蕉 & mángguǒ / xiāngjiāo & n. & Mangue / Banane \\
炸鱼 & zhá yú & n. & Poisson frit (spécialité rue) \\
法郎 & Fǎláng & n. & Franc CFA \\
便宜 / 贵 & piányi / guì & adj. & Pas cher / Cher \\
拥挤 / 堵车 & yǒngjǐ / dǔchē & adj./v. & Bondé / Embouteillage \\
附近 & fùjìn & n./loc. & Voisinage, près de \\
小饭馆 & xiǎo fànguǎn & n. & Petit restaurant populaire \\
\hline
\end{tabularx}
\end{center}
\end{definition}

\begin{attention}[Erreurs fréquentes des francophones — Pièges à éviter]
\begin{itemize}
    \item \textbf{\zh{如果……会……} :} \textcolor{red}{*如果下雨，我们会打车。} $\rightarrow$ \textbf{Correct :} \zh{如果下雨，我们就打车。} (\zh{就} marque la conséquence ; \zh{me} exprime la probabilité future, pas la conséquence conditionnelle immédiate).
    \item \textbf{\zh{是……了} :} \textcolor{red}{*我是坐摩托车回来了的。} $\rightarrow$ \textbf{Correct :} \zh{我是坐摩托车回来的。} (\zh{是……的} verrouille le passé, \zh{了} est redondant et incorrect ici).
    \item \textbf{Place de \zh{就} :} \textcolor{red}{*我们打就车。} $\rightarrow$ \textbf{Correct :} \zh{我们就打车。} (\zh{juste avant le verbe}).
    \item \textbf{\zh{二} vs \zh{两} :} \textcolor{red}{*二千人, *两个人 (correct), *二个人 (incorrect pour quantité).} $\rightarrow$ \textbf{Règle :} \zh{两} pour quantités (200, 2000 + classif.), \zh{二} pour comptage pur, numéros, décimales.
    \item \textbf{Classificateur manquant :} \textcolor{red}{*三芒果, *五法郎。} $\rightarrow$ \textbf{Correct :} \zh{三个芒果 / 三斤芒果}, \zh{五千法郎} (\zh{法郎} peut agir comme classif. monétaire, mais \zh{个} interdit pour fruits).
    \item \textbf{\zh{有的……也有的……} :} \textcolor{red}{*有的买菜，也有的卖菜。} $\rightarrow$ \textbf{Correct :} \zh{有的买菜，有的卖菜。} (\zh{也} est pléonastique, la structure implique déjà l'alternance).
\end{itemize}
\end{attention}

\begin{methode}[Stratégie : Rédiger un récit intégré en 4 étapes]
\begin{enumerate}
    \item \textbf{Scénario \& Mots-clés (3 min) :} Choisir le lieu (Mokolo / Stade), la météo, les achats, le retour. Lister 5--7 mots-clés chinois (ex. \zh{下雨}, \zh{打车}, \zh{三斤}, \zh{坐摩托车}).
    \item \textbf{Squelette grammatical (5 min) :} Écrire les 4 phrases « contraintes » (une par structure) en caractères + pinyin. Vérifier l'ordre : \zh{如果……，就……} ; \zh{有的……，有的……} ; \zh{是[Time/Lieu/Moyen]……的} ; Nombres + Classif.
    \item \textbf{Chair \& Liaison (7 min) :} Ajouter phrases de contexte (début : \zh{今天是星期天……}, transition : \zh{中午……}, fin : \zh{最后……}). Utiliser connecteurs : \zh{因为……所以……}, \zh{虽然……但是……}, \zh{然后}.
    \item \textbf{Relecture croisée (5 min) :} Le \emph{vérificateur} contrôle : tons (pinyin), \zh{的/得/地}, \zh{了/过} absence dans \zh{是……的}, \zh{两} vs \zh{二}, classif., ponctuation chinoise (\zh{，} \zh{。} \zh{；}).
\end{enumerate}
\end{methode}

\begin{savaistu}[Culture : Transports et monnaie — Cameroun vs Chine]
\begin{itemize}
    \item \textbf{Moto-taxi :} Au Cameroun, \zh{奥凯} (àokǎi, de l'anglais « okada ») ou \zh{摩托车} (mótuōchē). En Chine, \zh{摩托车} existe mais les \zh{网约车} (wǎngyuēchē, VTC type Didi) et le métro (\zh{地铁}) dominent en ville. Le casque est obligatoire dans les deux pays.
    \item \textbf{Monnaie :} Cameroun : \zh{非洲法郎} (Fēizhōu Fǎláng) / \zh{中非法郎} (Zhōngfēi Fǎláng), code XAF, taux fixe 1 EUR = 655,957 FCFA. Chine : \zh{人民币} (Rénmínbì), unité \zh{元} (yuán) / oral \zh{块} (kuài), 1 \zh{元} = 10 \zh{角} (jiǎo) = 100 \zh{分} (fēn). \textbf{Ne jamais dire \zh{块} pour les FCFA.}
    \item \textbf{Marché :} \zh{莫科洛市场} (Mokolo) = grand marché de gros et détail, très dense. Chine : \zh{农贸市场} (nóngmào shìchǎng, marché agricole) ou \zh{批发市场} (pīfā shìchǎng, marché de gros). Négociation (\zh{还价} huánjià) courante dans les deux cultures.
\end{itemize}
\end{savaistu}

\begin{popfigure}
\begin{center}\tcbox[colback=popink,colframe=popink,coltext=white,arc=9pt,boxsep=4pt,left=6pt,right=6pt]{\bfseries\small Leçon 6 Intégration}\end{center}
\vspace{-2pt}
\begin{tcbraster}[raster columns=2,raster equal height=rows,raster column skip=6pt,raster row skip=6pt,enhanced,blankest]
\begin{tcolorbox}[enhanced,colback=white,colframe=poporange,boxrule=0.9pt,arc=3pt,title={如果……就/那么 Hypothèse},fonttitle=\bfseries\scriptsize,coltitle=white,colbacktitle=poporange,left=4pt,right=4pt,top=2pt,bottom=2pt,boxsep=1pt,toptitle=1.5pt,bottomtitle=1.5pt,halign=left]
\begin{itemize}[nosep,leftmargin=*,itemsep=1pt,font=\scriptsize]
\item Condition réaliste
\item 就 avant verbe
\item Pas de 会
\end{itemize}
\end{tcolorbox}
\begin{tcolorbox}[enhanced,colback=white,colframe=popgreen,boxrule=0.9pt,arc=3pt,title={Nombres Cardinaux / Prix},fonttitle=\bfseries\scriptsize,coltitle=white,colbacktitle=popgreen,left=4pt,right=4pt,top=2pt,bottom=2pt,boxsep=1pt,toptitle=1.5pt,bottomtitle=1.5pt,halign=left]
\begin{itemize}[nosep,leftmargin=*,itemsep=1pt,font=\scriptsize]
\item 万 千 百 十
\item 两 vs 二
\item Classif. 斤/辆/位
\item FCFA = 法郎
\end{itemize}
\end{tcolorbox}
\begin{tcolorbox}[enhanced,colback=white,colframe=poppurple,boxrule=0.9pt,arc=3pt,title={有的……有的 Répartition},fonttitle=\bfseries\scriptsize,coltitle=white,colbacktitle=poppurple,left=4pt,right=4pt,top=2pt,bottom=2pt,boxsep=1pt,toptitle=1.5pt,bottomtitle=1.5pt,halign=left]
\begin{itemize}[nosep,leftmargin=*,itemsep=1pt,font=\scriptsize]
\item Sujet omis (ellipse)
\item Virgule obligatoire
\item 还有的 (3e groupe)
\end{itemize}
\end{tcolorbox}
\begin{tcolorbox}[enhanced,colback=white,colframe=poppink,boxrule=0.9pt,arc=3pt,title={是……的 Mise en relief},fonttitle=\bfseries\scriptsize,coltitle=white,colbacktitle=poppink,left=4pt,right=4pt,top=2pt,bottom=2pt,boxsep=1pt,toptitle=1.5pt,bottomtitle=1.5pt,halign=left]
\begin{itemize}[nosep,leftmargin=*,itemsep=1pt,font=\scriptsize]
\item Action PASSÉE
\item Interdit 了
\item Temps / Lieu / Moyen
\item Négat. 不是……的
\end{itemize}
\end{tcolorbox}
\end{tcbraster}

\legende{Carte mentale des 4 structures cibles de la Leçon 6}
\end{popfigure}

\courssub{Production intégrée guidée}

\begin{experience}[Activité de groupe : « Récit d'un dimanche à Yaoundé » (45 min)]
\textbf{Organisation :} Groupes de 3 élèves. Rôles tournants : \emph{Secrétaire} (écrit caractères + pinyin), \emph{Lecteur} (entraîne la lecture), \emph{Vérificateur} (valide les 4 structures + tons + classif.).

\textbf{Déroulement :}
\begin{enumerate}
    \item \textbf{Préparation (10 min) :} Brainstorming scénario (météo, lieu, achats, retour). Remplir une fiche brouillon : \textit{Mots-clés / 4 phrases squelettes / Connecteurs}.
    \item \textbf{Rédaction collaborative (15 min) :} Écriture du texte final (6--8 phrases) sur feuille A4 / tablette. Contraintes : caractères simplifiés, pinyin au-dessus ou en marge, ponctuation chinoise. Respect strict de l'ordre S + T + L + V + C.
    \item \textbf{Répétition orale (5 min) :} Lecteur s'entraîne (tons, liaisons, débit). Secrétaire et Vérificateur corrigent : \zh{就} (4\textsuperscript{e} ton), \zh{有} (3\textsuperscript{e} ton \textrightarrow\ 2\textsuperscript{e} ton devant 4\textsuperscript{e} ton), \zh{的} (ton neutre), \zh{法郎} (2-2).
    \item \textbf{Présentation plénière (10 min) :} Chaque groupe passe. Lecteur lit \textbf{deux fois} : 1. débit naturel, 2. débit lent. Classe écoute sans écrire.
    \item \textbf{Identification grammaticale (5 min) :} Après chaque groupe, 30 s pour noter : \textit{Phrase n°X = 如果结构 / Phrase n°Y = 有的结构 / etc.} Élève tiré au sort donne réponse orale ; classe valide/corrige.
\end{enumerate}
\textbf{Critères d'évaluation (grille partagée) :} (1) 4 structures présentes et correctes (8 pts), (2) Vocabulaire HSK 2--3 varié (4 pts), (3) Cohérence narrative (3 pts), (4) Prononciation / Tons (3 pts), (5) Ponctuation / Caractères lisibles (2 pts). Total 20 pts.
\end{experience}

\begin{exoresolu}[Production modèle commentée : « Un dimanche à Yaoundé »]
\textbf{Texte intégral (Caractères + Pinyin + Traduction) :}

\zh{今天是星期天，天气预报说会下雨。} \\
(Jīntiān shì xīngqítiān, tiānqì yùbào shuō huì xiàyǔ.) \\
Aujourd'hui c'est dimanche, la météo annonce qu'il va pleuvoir.

\zh{如果下雨，我们就打出租车去莫科洛市场；如果不下雨，那么我们就走路去。} \\
(Rúguǒ xiàyǔ, wǒmen jiù dǎ chūzūchē qù Mòkēluò Shìchǎng; rúguǒ bù xiàyǔ, nàme wǒmen jiù zǒulù qù.) \\
S'il pleut, nous prendrons un taxi pour aller au marché Mokolo ; s'il ne pleut pas, alors nous irons à pied.

\zh{我们是九点半出发的，十点到的市场。} \\
(Wǒmen shì jiǔ diǎn bàn chūfā de, shí diǎn dào de shìchǎng.) \\
C'est à 9 h 30 que nous sommes partis, c'est à 10 h que nous sommes arrivés au marché.

\zh{市场里人非常多，大概有两千多人。} \\
(Shìchǎng lǐ rén fēicháng duō, dàgài yǒu liǎng qiān duō rén.) \\
Il y a énormément de monde au marché, environ plus de deux mille personnes.

\zh{有的买芒果和香蕉，有的卖衣服和鞋子，还有的在吃炸鱼。} \\
(Yǒude mǎi mángguǒ hé xiāngjiāo, yǒude mài yīfu hé xiézi, hái yǒude zài chī zhá yú.) \\
Certains achètent des mangues et des bananes, d'autres vendent des vêtements et des chaussures, et d'autres encore sont en train de manger du poisson frit.

\zh{我买了三斤芒果，花了一万两千法郎。} \\
(Wǒ mǎi le sān jīn mángguǒ, huā le yī wàn liǎng qiān Fǎláng.) \\
J'ai acheté 3 catties (1,5 kg) de mangues, j'ai dépensé 12\,000 FCFA.

\zh{中午我们是在市场附近的小饭馆吃的饭。} \\
(Zhōngwǔ wǒmen shì zài shìchǎng fùjìn de xiǎo fànguǎn chī de fàn.) \\
Midi, c'est dans un petit restaurant près du marché que nous avons mangé.

\zh{下午三点，因为堵车，我们是坐摩托车回家的。} \\
(Xiàwǔ sān diǎn, yīnwèi dǔchē, wǒmen shì zuò mótuōchē huíjiā de.) \\
À 15 h, à cause des embouteillages, c'est en moto-taxi que nous sommes rentrés.

\tcblower
\textbf{Commentaire détaillé des choix de langue (Corrigé commenté) :}

\begin{enumerate}
    \item \textbf{Phrase 1 (Contexte) :} \zh{今天是星期天} (Temps + Copule). \zh{天气预报} (voc. HSK 3). \zh{会下雨} (Modal \zh{会} + V : probabilité future).
    \item \textbf{Phrase 2 (Structure 1 : 假设) :} \textbf{Deux branches} (pluie / pas pluie) pour montrer l'alternance. \zh{就} placé \textbf{avant} le verbe (\zh{打} / \zh{走}). \textbf{Piège évité :} Pas de \zh{会} après \zh{就} (*\zh{我们就会打车} : \zh{就} suffit pour la conséquence). Ponctuation : point-virgule \zh{；} pour séparer deux propositions longues coordonnées.
    \item \textbf{Phrase 3 (Structure 4 : 是……的 — Temps) :} Double mise en relief : \zh{是九点半出发的} / \zh{十点到的}. \textbf{Règle respectée :} Pas de \zh{了} après \zh{出发} ni \zh{到}. L'élément mis en relief (l'heure) intercalé entre \zh{是} et \zh{的}.
    \item \textbf{Phrase 4 (Nombres) :} \zh{大概} (environ). \zh{两千多人} (plus de 2000). \textbf{Note :} \zh{两} (liǎng) pour 2000. Classificateur \zh{人} implicite après \zh{多}.
    \item \textbf{Phrase 5 (Structure 3 : 有的……有的……) :} Trois groupes (\zh{有的…}，\zh{有的…}，\zh{还有的…}). Sujets omis (récupérés du contexte « gens du marché »). Prédicats variés : \zh{买}, \zh{卖}, \zh{在吃} (action en cours). \zh{炸鱼} (plat de rue typique).
    \item \textbf{Phrase 6 (Nombres + Classif. + Prix + Aspect) :} \zh{三斤} (3 catties = 1,5 kg, unité standard fruits). \zh{花了} (dépenser) + \zh{了} (aspect accompli, \textbf{pas} structure \zh{是……的}). \zh{一万两千} (12\,000) : structure \zh{万} + \zh{千} correcte. \zh{法郎} (FCFA).
    \item \textbf{Phrase 7 (Structure 4 : 是……的 — Lieu) :} \zh{是在……吃的}. Mise en relief du \textbf{lieu} (\zh{在小饭馆}). Verbe \zh{吃} répété avant \zh{的} (recommandé niveau Première pour clarté). \zh{附近} (près de).
    \item \textbf{Phrase 8 (Structure 4 : 是……的 — Moyen + Cause) :} \zh{因为堵车} (cause) placée \textbf{avant} la structure \zh{是……的} (ordre standard : Cause $\rightarrow$ Focus). \zh{坐摩托车} (moyen). \zh{堵车} (embouteillage), \zh{快} (rapide) implicite dans la justification.
\end{enumerate}

\textbf{Check-list d'auto-correction (à distribuer aux élèves) :}
\begin{itemize}
    \item [\checkmark] \zh{如果} \dots \zh{，就/那么} \dots (virgule, \zh{juste avant verbe}, pas de \zh{会}).
    \item [\checkmark] Nombres : \zh{两} pour 200/2000/20000 + classif. ; \zh{万}/\zh{千} corrects ; \zh{法郎} sans \zh{元}/\zh{块}.
    \item [\checkmark] \zh{有的} \dots \zh{，有的} \dots (virgule, ellipse sujet, \zh{还有的} possible).
    \item [\checkmark] \zh{是} [Temps/Lieu/Moyen] \zh{的} (pas de \zh{了} sur verbe, action passée, ordre S + 是 + Complément + V + 的).
\end{itemize}
\end{exoresolu}

\courssub{Entraînement individuel}

\begin{exercice}[Exercices d'application — À faire en classe ou en devoir]
\textbf{A. Transformation : Hypothèse (Si $\rightarrow$ 如果……就……)}
Traduire en chinois (caractères + pinyin).
\begin{enumerate}
    \item S'il fait chaud, je bois de l'eau froide. (\zh{热} rè, \zh{喝} hē, \zh{冰水} bīngshuǐ)
    \item S'ils ne vont pas au stade, ils regardent le match à la télé. (\zh{去} qù, \zh{看电视} kàn diànshì)
    \item Si ce n'est pas cher, nous en achetons deux. (\zh{便宜} piányi, \zh{买} mǎi, \zh{两个} liǎng gè)
\end{enumerate}

\textbf{B. Nombres et Classificateurs : Compléter}
Écrire en caractères et pinyin.
\begin{enumerate}
    \item 35\,000 FCFA \hfill \zh{\_\_\_\_\_\_\_\_\_\_\_\_}
    \item 8 moto-taxis \hfill \zh{\_\_\_\_\_\_\_\_\_\_\_\_}
    \item 1,5 kg de bananes \hfill \zh{\_\_\_\_\_\_\_\_\_\_\_\_}
    \item La 3\textsuperscript{e} personne \hfill \zh{\_\_\_\_\_\_\_\_\_\_\_\_}
\end{enumerate}

\textbf{C. Répartition : Relier et écrire une phrase \zh{有的……有的……}}
Contexte : \textit{Au stade, il y a beaucoup de supporters.}
\begin{itemize}
    \item Groupe A : chanter / Groupe B : taper dans les mains / Groupe C : manger des arachides.
\end{itemize}
$\rightarrow$ \zh{\_\_\_\_\_\_\_\_\_\_\_\_}

\textbf{D. Mise en relief \zh{是……的} : Reformuler la phrase en insistant sur l'élément souligné.}
\begin{enumerate}
    \item Nous \underline{sommes arrivés à 10 h}. (\zh{我们十点到了。})
    \item Il \underline{a acheté les mangues au marché Mokolo}. (\zh{他在莫科洛市场买了芒果。})
    \item Elle \underline{est rentrée en taxi}. (\zh{她坐出租车回去了。})
    \item Nous \underline{avons mangé au restaurant}. (\zh{我们在饭馆吃了饭。})
\end{enumerate}

\textbf{E. Traduction thème (Français $\rightarrow$ Chinois) :}
\begin{enumerate}
    \item Dimanche dernier, s'il pleuvait, nous prenions le taxi ; s'il ne pleuvait pas, nous marchions. (Imparfait d'habitude $\rightarrow$ \zh{如果……就……} au passé).
    \item Il y avait plus de 5\,000 personnes. Certains regardaient le match, d'autres vendaient de l'eau.
    \item C'est à 16 h que le match a fini. C'est à moto que nous sommes rentrés.
\end{enumerate}
\end{exercice}

\begin{corrige}[Exercices d'application]
\textbf{A.}
\begin{enumerate}
    \item \zh{如果天热，我就喝冰水。} (Rúguǒ tiān rè, wǒ jiù hē bīngshuǐ.)
    \item \zh{如果他们不去体育场，就看电视。} (Rúguǒ tāmen bù qù tǐyùchǎng, jiù kàn diànshì.)
    \item \zh{如果不贵，我们就买两个。} (Rúguǒ bù guì, wǒmen jiù mǎi liǎng gè.)
\end{enumerate}
\textbf{B.}
\begin{enumerate}
    \item \zh{三万五千法郎} (sān wàn wǔ qiān Fǎláng)
    \item \zh{八辆摩托车} (bā liàng mótuōchē)
    \item \zh{三斤香蕉} (sān jīn xiāngjiāo)
    \item \zh{第三个人 / 第三位} (dì sān gè rén / dì sān wèi)
\end{enumerate}
\textbf{C.}
\zh{体育场里人很多，有的唱歌，有的鼓掌，还有的吃花生。} (Tǐyùchǎng lǐ rén hěn duō, yǒude chànggē, yǒude gǔzhǎng, hái yǒude chī huāshēng.)
\textbf{D.}
\begin{enumerate}
    \item \zh{我们是十点到的。} (Focus temps)
    \item \zh{他是在莫科洛市场买的芒果。} (Focus lieu)
    \item \zh{她是坐出租车回去的。} (Focus moyen)
    \item \zh{我们是在饭馆吃的饭。} (Focus lieu)
\end{enumerate}
\textbf{E.}
\begin{enumerate}
    \item \zh{上个星期天，如果下雨，我们就打车；如果不下雨，我们就走路。} (Contexte passé habituel : \zh{如果……就……} s'utilise aussi pour le passé récurrent).
    \item \zh{有五千多人，有的看球，有的卖水。}
    \item \zh{比赛是十六点结束的。我们是坐摩托车回家的。}
\end{enumerate}
\end{corrige}

\courssub{Bilan et prolongements}

\begin{aretenir}[Bilan grammatical et communicatif — Leçon 6]
Cette leçon d'intégration a permis de vérifier l'acquisition opérationnelle des quatre structures cibles dans une situation de communication authentique (\emph{Vie familiale et intégration sociale : Dimanche à Yaoundé}).

\begin{itemize}
    \item \textbf{Hypothèse (\zh{如果……就……}) :} L'élève place \zh{就} avant le verbe, utilise la virgule, distingue \zh{juste} (oral) / \zh{那么} (formel). \textbf{Erreur résiduelle :} ajout de \zh{会} ou inversion \zh{就}/V.
    \item \textbf{Nombres :} Maîtrise du système \zh{万}/\zh{千}, alternance \zh{二}/\zh{两}, accord Classificateur/Nom (\zh{三斤} $\neq$ *\zh{三个斤}). \textbf{Point dur :} Prix en FCFA (\zh{法郎}) sans classif. monétaire chinois.
    \item \textbf{Répartition (\zh{有的……有的……}) :} Ellipse du sujet maîtrisée, variation des prédicats (V, \zh{在}V, Adj). \textbf{Erreur :} omission de la virgule ou ajout pléonastique de \zh{也}.
    \item \textbf{Mise en relief (\zh{是……的}) :} Compréhension du verrouillage \textbf{passé}. Absence de \zh{了} = marqueur de réussite. Distinction Temps / Lieu / Moyen / Manière opérationnelle.
\end{itemize}

\textbf{Compétences transversales travaillées :}
\begin{itemize}
    \item \textbf{Écriture collaborative :} Négociation de sens, rôles (scribe / relecteur / orateur).
    \item \textbf{Phonologie :} Lecture à voix haute $\rightarrow$ contrôle tons (\zh{就} jiù 4, \zh{有} yǒu 3$\rightarrow$2, \zh{是} shì 4, \zh{法郎} Fǎláng 2-2, ton neutre \zh{de}).
    \item \textbf{Écoute sélective :} Identification structures cibles dans flux oral pairs (analyse linguistique).
    \item \textbf{Culture :} Ancrage camerounais (Mokolo, Ahidjo, okada, FCFA, poisson frit) $\leftrightarrow$ Vocabulaire chinois standard HSK 2--3.
\end{itemize}
\end{aretenir}

\begin{exercice}[Prolongements — Projet / Devoir (au choix)]
\begin{enumerate}
    \item \textbf{Carnet de voyage illustré :} Transformer le texte du groupe en 6--8 vignettes (photos/dessins) avec légendes en chinois (caractères + pinyin). Ajouter 2 dialogues : marchander (\zh{太贵了，便宜一点儿吧!}) et demander le chemin (\zh{请问，去体育场怎么走?}).
    \item \textbf{Saynète filmée (1--2 min) :} Jouer la scène (rôles : narrateur, vendeur, ami, chauffeur moto). Sous-titrer en chinois. Envoyer sur la plateforme OnBuch+.
    \item \textbf{Version « Chine » :} Réécrire le récit en transposant le contexte : \zh{北京} (Běijīng), \zh{潘家园旧货市场} (Marché aux puces Panjiayuan), \zh{工人体育场} (Stade des Ouvriers), \zh{共享单车} (vélos partagés) / \zh{地铁} (métro), \zh{元/块} (RMB). Comparer les différences lexicologiques.
\end{enumerate}
\end{exercice}

\newpage

\coursec{Méthodes et exercices résolus}

\courssub{Méthodes et exercices résolus}

\begin{methode}[Construire une phrase hypothétique avec \cle{如果……，那么/就……}]
\textbf{Objectif :} Exprimer une condition et sa conséquence logique (si… alors…).
\textbf{Étapes :}
\begin{enumerate}
    \item \textbf{Identifier la condition (protase) :} C'est l'événement supposé. La placer \textbf{après \cle{如果} (rúguǒ)} en début de phrase (Sujet + 如果 + Condition).
    \item \textbf{Identifier la conséquence (apodose) :} C'est le résultat si la condition est réalisée. La placer \textbf{après \cle{那么} (nàme) ou \cle{就} (jiù)} en début de la proposition principale (Sujet + 那么/就 + Conséquence).
    \item \textbf{Choisir \cle{那么} ou \cle{就} :} \cle{那么} insiste sur la suite logique (« dans ce cas ») ; \cle{就} insiste sur l'immédiateté, l'inévitabilité ou la décision (« dès lors, tout de suite »).
    \item \textbf{Vérifier l'ordre des mots :} Sujet + 如果 + Condition ， Sujet + 那么/就 + Verbe/Adjectif + Complément. \textbf{Pas de « alors » (alors) en chinois comme en français.}
    \item \textbf{Gérer la négation :} Elle se place \textbf{avant} le verbe dans la proposition concernée (\zh{如果不下雨，我们就去}。).
\end{enumerate}
\textbf{Structures à retenir :}
\begin{itemize}
    \item \zh{如果 + 条件，那么/就 + 结果。}
    \item \zh{如果你有空，那么我们一起去吃饭。} (Rúguǒ nǐ yǒukòng, nàme wǒmen yīqǐ qù chīfàn.)
    \item \zh{如果明天下雨，比赛就取消。} (Rúguǒ míngtiān xiàyǔ, bǐsài jiù qǔxiāo.)
\end{itemize}
\end{methode}

\begin{exoresolu}[Transformation et production : Hypothèses scolaires]
\textbf{Énoncé :}
\begin{enumerate}
    \item Transforme les paires de phrases en une seule phrase avec \cle{如果……，那么……} ou \cle{如果……，就……}.
    \begin{itemize}
        \item (a) \zh{明天不下雨。} \zh{我们去爬山。}
        \item (b) \zh{你努力学习。} \zh{你一定能通过HSK 3。}
        \item (c) \zh{我有钱。} \zh{我买一张去北京的机票。}
    \end{itemize}
    \item Complète les dialogues suivants en utilisant la structure hypothétique (choisis \cle{那么} ou \cle{就} selon le contexte).
    \begin{itemize}
        \item A: \zh{我想买那个中文词典。} B: \zh{如果太贵，\_\_\_\_\_\_ 我们不买。}
        \item A: \zh{王老师病了。} B: \zh{如果他明天不来，\_\_\_\_\_\_ 谁教我们中文？}
    \end{itemize}
    \item \textbf{Production :} Écris 3 phrases sur tes projets pour les vacances (Noël / Grande vacances) en utilisant \cle{如果……，就……}. Intègre du vocabulaire camerounais (ex: \zh{回村里}, \zh{吃恩多莱}, \zh{看足球赛}).
\end{enumerate}

\tcblower

\textbf{Solution détaillée :}

\textbf{1. Transformation (Règle : Sujet + 如果 + Condition， Sujet + 那么/就 + Résultat)}
\begin{itemize}
    \item (a) \zh{如果明天不下雨，我们就去爬山。} (Rúguǒ míngtiān bù xiàyǔ, wǒmen jiù qù páshān.) $\rightarrow$ « Si demain il ne pleut pas, nous irons faire de la randonnée. » \textit{Justification :} Action future décidée, \cle{就} marque la conséquence immédiate/décision.
    \item (b) \zh{如果你努力学习，那么你一定能通过HSK 3。} (Rúguǒ nǐ nǔlì xuéxí, nàme nǐ yídìng néng tōngguò HSK 3.) $\rightarrow$ « Si tu étudies assidûment, alors tu pourras certainement passer le HSK 3. » \textit{Justification :} Inférence logique générale, \cle{那么} convient bien.
    \item (c) \zh{如果我有钱，我就买一张去北京的机票。} (Rúguǒ wǒ yǒu qián, wǒ jiù mǎi yī zhāng qù Běijīng de jīpiào.) $\rightarrow$ « Si j'ai de l'argent, j'achèterai un billet d'avion pour Pékin. » \textit{Justification :} Décision personnelle conditionnelle, \cle{就} est naturel.
\end{itemize}

\textbf{2. Complétion (Choix \cle{那么} vs \cle{就})}
\begin{itemize}
    \item B: \zh{如果太贵，\textbf{那么} 我们不买。} (Inférence logique : « Dans ce cas, on n'achète pas »).
    \item B: \zh{如果他明天不来，\textbf{就} 谁教我们中文？} (Conséquence immédiate/pratique sur l'organisation : « Du coup, qui nous enseigne ? »).
\end{itemize}

\textbf{3. Production (Exemples de réponses attendues) :}
\begin{enumerate}
    \item \zh{如果放假，我就回村里看望奶奶。} (Si je suis en vacances, je retourne au village voir grand-mère.)
    \item \zh{如果周末有比赛，我就和朋友们一起看足球赛。} (S'il y a un match le week-end, je regarde le foot avec mes amis.)
    \item \zh{如果妈妈做恩多莱，我就多吃两碗饭。} (Si maman fait du ndolé, je mange deux bols de riz de plus.)
\end{enumerate}
\textbf{Conclusion :} Maîtriser l'ordre \cle{如果} (Condition) $\rightarrow$ \cle{那么/就} (Conséquence) et le choix du connecteur logique est essentiel pour l'expression de l'hypothèse au Bac.
\end{exoresolu}

\begin{methode}[Maîtriser les nombres cardinaux et ordinaux (HSK 1-3) : lecture, écriture, classificateurs]
\textbf{Objectif :} Compter, donner des numéros, des dates, des prix, des âges sans erreur de tons ni de structure.
\textbf{Étapes :}
\begin{enumerate}
    \item \textbf{Les bases (0-10) :} Mémoriser les tons critiques : \cle{一} (yī $\rightarrow$ yí/yì selon ton suivant), \cle{二} (èr) vs \cle{两} (liǎng), \cle{七} (qī), \cle{八} (bā).
    \item \textbf{Règle du \cle{一} (yī) :}
        \begin{itemize}
            \item Seul / à la fin / ordinal : \textbf{yī} (1er ton).
            \item Devant ton 4 (ex: \zh{一定} yídìng) : \textbf{yí} (2e ton).
            \item Devant tons 1, 2, 3 (ex: \zh{一个} yí gè) : \textbf{yí} (2e ton) \textit{ou} neutre en comptage rapide.
        \end{itemize}
    \item \textbf{Règle du \cle{二} vs \cle{两} :} \cle{二} pour compter (1, 2, 3...), numéros, dates, numéros de téléphone. \cle{两} \textbf{obligatoire} devant classificateur (\zh{两个人} liǎng gè rén, \zh{两本书} liǎng běn shū). Exception : \zh{两点钟} (2 heures).
    \item \textbf{Construction 11-99 :} Dizaine + \cle{十} (shí) + Unité. Ex: \zh{二十一} (èrshíyī), \zh{三十} (sānshí). \textit{Attention :} 10-19 = \cle{十} + unité (\zh{十一} shíyī, pas *yīshíyī).
    \item \textbf{Grands nombres :} \cle{百} (bǎi), \cle{千} (qiān), \cle{万} (wàn = 10 000), \cle{亿} (yì = 100 000 000). \textbf{Zéro (\cle{零} líng)} : un seul \cle{零} pour combler un trou entre unités non nulles (\zh{一百零一} 101, \zh{一万零一百} 10 100).
    \item \textbf{Ordination :} Ajouter \cle{第} (dì) devant le cardinal. \zh{第一} (dì yī), \zh{第二} (dì èr), \zh{第二十一} (dì èrshíyī).
    \item \textbf{Classificateurs (Quantificateurs) :} Nombre + Classificateur + Nom. \cle{个} (gè) universel ; \cle{本} (běn) livres ; \cle{张} (zhāng) objets plats ; \cle{位} (wèi) personnes (politesse) ; \cle{只} (zhī) animaux/paires ; \cle{条} (tiáo) longs/étroits ; \cle{辆} (liàng) véhicules.
\end{enumerate}
\end{methode}

\begin{exoresolu}[Nombres : Dictée, conversion, classificateurs et pièges]
\textbf{Énoncé :}
\begin{enumerate}
    \item \textbf{Dictée / Conversion :} Écris en caractères (simplifiés) et en pinyin (avec tons) :
    \begin{itemize}
        \item 14 ; 20 ; 52 ; 100 ; 101 ; 110 ; 250 ; 1 000 ; 10 001 ; 50 000.
    \end{itemize}
    \item \textbf{Choix \cle{二} / \cle{两} :} Complète avec \zh{二} ou \zh{两} (caractères).
    \begin{itemize}
        \item \_\_\_\_ 个学生 (2 élèves) ; \_\_\_\_ 月 (février) ; \_\_\_\_ 点钟 (2 heures) ; 第 \_\_\_\_ 课 (leçon 2) ; \_\_\_\_ 百元 (200 yuans) ; 2024 年 \_\_\_\_ 月 \_\_\_\_ 日.
    \end{itemize}
    \item \textbf{Classificateurs :} Associe le nombre + classificateur correct au nom.
    \begin{itemize}
        \item 3 livres : \_\_\_\_ \_\_\_\_ 书
        \item 5 professeurs (politesse) : \_\_\_\_ \_\_\_\_ 老师
        \item 1 voiture : \_\_\_\_ \_\_\_\_ 车
        \item 8 billets (train/avion) : \_\_\_\_ \_\_\_\_ 票
        \item 1 chemise : \_\_\_\_ \_\_\_\_ 衬衫
    \end{itemize}
    \item \textbf{Application contextuelle (Cameroun/Chine) :} Traduis en chinois.
    \begin{itemize}
        \item « Yaoundé a environ 4 millions d'habitants. » (万 wàn)
        \item « Le match commence à 15h30. » (点 分)
        \item « C'est mon 3e voyage en Chine. » (第 次)
    \end{itemize}
\end{enumerate}

\tcblower

\textbf{Solution détaillée :}

\textbf{1. Dictée / Conversion (Attention aux zéros et à \cle{万})}
\begin{center}
\begin{tabularx}{\linewidth}{|X|X|X|}
\hline \rowcolor{popblueL} \textbf{Nombre} & \textbf{Caractères} & \textbf{Pinyin} \\ \hline
14 & 十四 & shí sì \\ \hline
20 & 二十 & èr shí \\ \hline
52 & 五十二 & wǔ shí èr \\ \hline
100 & 一百 & yī bǎi \\ \hline
101 & 一百零一 & yī bǎi \textbf{líng} yī \\ \hline
110 & 一百一十 & yī bǎi yī shí \\ \hline
250 & 二百五十 & èr bǎi wǔ shí \\ \hline
1 000 & 一千 & yī qiān \\ \hline
10 001 & 一万零一 & yī wàn \textbf{líng} yī \\ \hline
50 000 & 五万 & wǔ wàn \\ \hline
\end{tabularx}
\end{center}
\textit{Justification :} \cle{零} obligatoire pour 101 (trou aux dizaines) et 10 001 (trou aux milliers/centaines/dizaines). Pas de \cle{零} final. 110 s'écrit « cent dix » (yī bǎi yī shí), pas « cent dix zéro ».

\textbf{2. \cle{二} vs \cle{两} (Règle : Classificateur $\rightarrow$ \cle{两} ; Comptage/Date/Numéro $\rightarrow$ \cle{二})}
\begin{itemize}
    \item \zh{两} 个学生 (Classificateur \cle{个})
    \item \zh{二} 月 (Numéro de mois/Date)
    \item \zh{两} 点钟 (Heure : exception, on dit \cle{两点} et non \cle{二点})
    \item 第 \zh{二} 课 (Ordinal)
    \item \zh{二} 百元 (Comptage pur : 200 = deux cents. Mais \zh{两百元} est aussi accepté à l'oral pour « deux cents yuans » comme quantité. Au Bac, \cle{二百} est la forme standard pour le nombre 200).
    \item 2024 年 \zh{二} 月 \zh{二} 日 (Dates : toujours \cle{二}).
\end{itemize}

\textbf{3. Classificateurs}
\begin{itemize}
    \item 3 livres : \zh{三 本 书} (sān \textbf{běn} shū)
    \item 5 professeurs : \zh{五 位 老师} (wǔ \textbf{wèi} lǎoshī) — \cle{位} marque le respect.
    \item 1 voiture : \zh{一 辆 车} (yī \textbf{liàng} chē)
    \item 8 billets : \zh{八 张 票} (bā \textbf{zhāng} piào) — objets plats.
    \item 1 chemise : \zh{一 件 衬衫} (yī \textbf{jiàn} chènshān) — vêtements (haut).
\end{itemize}

\textbf{4. Traduction contextuelle}
\begin{itemize}
    \item \zh{雅温得大约有四百万人。} (Yǎwēndé dàyuē yǒu sì bǎi wàn rén.) — 4 millions = 400 $\times$ 10 000 = \cle{四百万}.
    \item \zh{比赛下午三点半开始。} (Bǐsài xiàwǔ sān diǎn bàn kāishǐ.) — 15h = 下午三点 ; 30 min = \cle{半} (bàn).
    \item \zh{这是我第三次去中国。} (Zhè shì wǒ dì sān cì qù Zhōngguó.) — Ordinal \cle{第} + cardinal + classificateur d'action \cle{次} (cì).
\end{itemize}
\textbf{Conclusion :} La confusion \cle{二}/\cle{两}, l'oubli du \cle{零} dans les grands nombres et le mauvais classificateur sont les 3 pièges majeurs notés au Bac.
\end{exoresolu}

\begin{methode}[Exprimer la répartition / classification avec \cle{有的……，有的……}]
\textbf{Objectif :} Décrire une diversité au sein d'un groupe (« certains…, d'autres… »).
\textbf{Étapes :}
\begin{enumerate}
    \item \textbf{Identifier le groupe de référence (Sujet global) :} Souvent omis en chinois si le contexte est clair, ou placé en \textbf{Topique} (début de phrase). Ex: \zh{我们班同学，有的……}
    \item \textbf{Structurer les deux membres :} \cle{有的} (yǒude) + Prédicat 1 ， \cle{有的} (yǒude) + Prédicat 2. (Peut y avoir 3+ membres avec \cle{还有的}).
    \item \textbf{Accorder le verbe :} \cle{有的} fonctionne comme un pronom indéfini (« certains »). Le verbe qui suit s'accorde logiquement (verbe d'état, action, ou verbe + complément).
    \item \textbf{Ne pas confondre avec \cle{有} (avoir) + Objet :} \zh{有的学生喜欢足球} (Certains élèves aiment le foot) $\neq$ \zh{学生有足球} (Les élèves ont un ballon).
    \item \textbf{Variante \cle{有的……，有的……，还有的……} :} Pour énumérer plus de deux catégories.
\end{enumerate}
\textbf{Schéma :} [Sujet global ,] \cle{有的} + Prédicat 1 ， \cle{有的} + Prédicat 2 (， \cle{还有的} + Prédicat 3)。
\end{methode}

\begin{exoresolu}[Répartition : Description de classe et vie quotidienne]
\textbf{Énoncé :}
\begin{enumerate}
    \item \textbf{Ordre des mots :} Remets les mots dans l'ordre pour former des phrases correctes avec \cle{有的……有的……}.
    \begin{itemize}
        \item (a) \zh{喜欢 / 同学 / 有的 / 吃 / 恩多莱 / 有的 / 喜欢 / 吃 / 米饭}
        \item (b) \zh{去 / 有的 / 上海 / 旅游 / 有的 / 回 / 村里 / 同学 / 放假}
    \end{itemize}
    \item \textbf{Traduction (Français $\rightarrow$ Chinois) :}
    \begin{itemize}
        \item « Dans notre classe, certains sont Chinois, d'autres sont Camerounais. » (Notre classe = \zh{我们班} wǒmen bān ; Chinois = \zh{中国人} ; Camerounais = \zh{喀麦隆人}).
        \item « Pour le déjeuner, certains mangent au réfectoire, d'autres apportent leur repas. » (Réfectoire = \zh{食堂} shítáng ; Apporter = \zh{自带} zìdài / \zh{带饭} dàifàn).
    \end{itemize}
    \item \textbf{Production orale/écrite :} Décris la situation des élèves de 1ère A après le Bac en utilisant \cle{有的……有的……还有的……} (au moins 3 catégories : études, travail, voyage, mariage...).
\end{enumerate}

\tcblower

\textbf{Solution détaillée :}

\textbf{1. Remise en ordre (Structure : Sujet global + 有的 + Pred1， 有的 + Pred2)}
\begin{itemize}
    \item (a) \zh{同学有的喜欢吃恩多莱，有的喜欢吃米饭。} (Tóngxué yǒude xǐhuan chī ēnduōlái, yǒude xǐhuan chī mǐfàn.) $\rightarrow$ « Les camarades, certains aiment manger du ndolé, d'autres aiment manger du riz. » \textit{Note :} Le sujet global \zh{同学} est en topique.
    \item (b) \zh{放假，同学有的去上海旅游，有的回村里。} (Fàngjià, tóngxué yǒude qù Shànghǎi lǚyóu, yǒude huí cūnlǐ.) $\rightarrow$ « Pendant les vacances, les élèves, certains vont à Shanghai en voyage, d'autres retournent au village. » \textit{Note :} Complément de temps \zh{放假} en topique initial.
\end{itemize}

\textbf{2. Traduction}
\begin{itemize}
    \item \zh{我们班，有的是中国人，有的是喀麦隆人。} (Wǒmen bān, yǒude shì Zhōngguórén, yǒude shì Kāmàilóngrén.) $\rightarrow$ \textit{Structure :} \cle{有的} + \cle{是} + Nationalité.
    \item \zh{午饭，有的在食堂吃，有的自带饭吃。} (Wǔfàn, yǒude zài shítáng chī, yǒude zìdài fàn chī.) $\rightarrow$ \textit{Structure :} Topique \zh{午饭} (Déjeuner), \cle{有的} + Locatif \zh{在食堂} + Verbe, \cle{有的} + Verbe composé \zh{自带饭}.
\end{itemize}

\textbf{3. Production (Exemple de réponse Bac) :}
\zh{高考以后，同学有的考大学，有的找工作，还有的想去中国留学。}
(Gāokǎo yǐhòu, tóngxué yǒude kǎo dàxué, yǒude zhǎo gōngzuò, hái yǒude xiǎng qù Zhōngguó liúxué.)
« Après le Bac, les élèves, certains passent le concours universitaire, d'autres cherchent du travail, et d'autres encore veulent aller étudier en Chine. »
\textbf{Conclusion :} \cle{有的} répété structure la pensée. Au Bac, l'utiliser pour nuancer un avis (\zh{有的人认为……，有的人认为……}) rapporte des points en expression écrite.
\end{exoresolu}

\begin{methode}[Mettre en relief une information passée avec \cle{是……的} (Focus sur temps/lieu/manière/auteur)]
\textbf{Objectif :} Insister sur \textbf{un détail précis} (temps, lieu, manière, auteur, but) d'une action \textbf{déjà accomplie} (passé).
\textbf{Étapes :}
\begin{enumerate}
    \item \textbf{Vérifier le prérequis :} L'action est \textbf{terminée} (aspect accompli). \cle{是……的} ne s'utilise pas pour le futur ou l'habitude présente.
    \item \textbf{Identifier l'élément mis en relief (Focus) :} Temps ? Lieu ? Manière ? Auteur (passif) ? But ? Moyen/Instrument ?
    \item \textbf{Construire la structure :} \cle{是} (souvent omis à l'oral) + \textbf{Élément Focus} + \textbf{Verbe (+ Compléments restants)} + \cle{的}.
    \item \textbf{Place de l'objet (si verbe transitif) :} L'objet reste \textbf{après le verbe} (ou avant \cle{是} si c'est l'objet qui est focusé, structure différente). \textit{Standard :} Sujet + \cle{是} + Focus + Verbe + Objet + \cle{的}.
    \item \textbf{Négation :} \cle{不是……的} (Focus non réalisé) ou \cle{没} + Verbe (Action non réalisée — \cle{是……的} impossible car action non accomplie).
    \item \textbf{Question :} Mettre le mot interrogatif \textbf{à la place du Focus}. Ex: \zh{你是几点来的？} (C'est à quelle heure que tu es venu ?).
\end{enumerate}
\textbf{Tableau récapitulatif des Focus courants :}
\begin{center}
\begin{tabularx}{\linewidth}{|X|X|X|}
\hline \rowcolor{popblueL} \textbf{Focus} & \textbf{Structure} & \textbf{Exemple} \\ \hline
Temps & 是 + 时间 + 动词 + 的 & \zh{我是昨天来的。} (C'est hier que je suis venu.) \\ \hline
Lieu & 是 + 地点 + 动词 + 的 & \zh{我们是在食堂吃的饭。} (C'est au réfectoire qu'on a mangé.) \\ \hline
Manière & 是 + 方式 + 动词 + 的 & \zh{他是坐飞机来的。} (C'est en avion qu'il est venu.) \\ \hline
Auteur (Passif) & 是 + 人 + 动词 + 的 & \zh{这本书是王老师买的。} (Ce livre a été acheté par M. Wang.) \\ \hline
But & 是 + 目的 + 动词 + 的 & \zh{我是来学中文的。} (Je suis venu pour apprendre le chinois.) \\ \hline
\end{tabularx}
\end{center}
\end{methode}

\begin{exoresolu}[Le focus \cle{是……的} : Transformation, correction, questions]
\textbf{Énoncé :}
\begin{enumerate}
    \item \textbf{Transformation (Focus sur l'élément en gras) :} Transforme la phrase neutre (passé accompli avec \cle{了}) en phrase \cle{是……的} pour mettre l'élément en gras en relief.
    \begin{itemize}
        \item (a) \zh{我 \textbf{昨天} 去了双球。} (Hier)
        \item (b) \zh{他在 \textbf{喀麦隆} 学的中文。} (Déjà \cle{的} $\rightarrow$ vérifier structure)
        \item (c) \zh{他们 \textbf{坐火车} 去了北京。} (Moyen de transport)
        \item (d) \zh{这道菜是 \textbf{我妈妈} 做的。} (Auteur — Passif)
    \end{itemize}
    \item \textbf{Correction d'erreurs (Typiques francophones) :} Corrige les phrases suivantes.
    \begin{itemize}
        \item (1) \zh{*我是去年去中国了。} (Focus temps)
        \item (2) \zh{*他是开车来的学校。} (Focus moyen + lieu mal placé)
        \item (3) \zh{*你是几点来了？} (Question focus temps)
    \end{itemize}
    \item \textbf{Questions au Bac :} Pose la question pour obtenir la réponse soulignée.
    \begin{itemize}
        \item Réponse : \zh{这个汉字是李老师教我的。} (Focus : Auteur)
        \item Réponse : \zh{我们是走路回家的。} (Focus : Manière)
    \end{itemize}
\end{enumerate}

\tcblower

\textbf{Solution détaillée :}

\textbf{1. Transformation (Règle : Sujet + 是 + Focus + Verbe + Objet/Compl + 的 ; \cle{了} disparaît)}
\begin{itemize}
    \item (a) \zh{我是昨天去双球的。} (Wǒ shì zuótiān qù Shuāngqiú de.) — \textit{Focus Temps. \cle{了} supprimé.}
    \item (b) \zh{他是在喀麦隆学的中文。} (Tā shì zài Kāmàilóng xué de Zhōngwén.) — \textit{Focus Lieu. \cle{在} conservé devant le lieu. Structure originale correcte mais \cle{是} ajouté pour l'insistance.}
    \item (c) \zh{他们是坐火车去北京的。} (Tāmen shì zuò huǒchē qù Běijīng de.) — \textit{Focus Moyen. \cle{坐火车} placé entre \cle{是} et \cle{去}.}
    \item (d) \zh{这道菜是我妈妈做的。} (Zhè dào cài shì wǒ māma zuò de.) — \textit{Focus Auteur (Passif). Sujet = Objet de l'action.}
\end{itemize}

\textbf{2. Correction d'erreurs}
\begin{itemize}
    \item (1) \zh{*我是去年去中国了。} $\rightarrow$ \zh{我是去年去中国的。} (\textbf{Erreur :} \cle{了} incompatible avec \cle{是……的}. \cle{的} marque l'accompli/le focus, \cle{了} marque l'aspect accompli simple. On ne garde qu'un marqueur).
    \item (2) \zh{*他是开车来的学校。} $\rightarrow$ \zh{他是开车来学校的。} OU \zh{他是开车来的。} (\textbf{Erreur :} L'objet/lieu \zh{学校} ne doit pas couper le lien Verbe-\cle{的}. Soit il reste après le verbe de mouvement \zh{来}, soit il est omis si contexte connu).
    \item (3) \zh{*你是几点来了？} $\rightarrow$ \zh{你是几点来的？} (\textbf{Erreur :} Même chose, pas de \cle{了} dans \cle{是……的}. Le mot interrogatif \cle{几点} remplace le focus temps).
\end{itemize}

\textbf{3. Questions (Règle : Mot interrogatif = Focus)}
\begin{itemize}
    \item Réponse : \zh{这个汉字是李老师教我的。} $\rightarrow$ Question : \zh{这个汉字是\textbf{谁}教的？} (Shéi = Qui). \textit{Note :} \cle{我} (moi) reste après le verbe car c'est l'objet réel de «教».
    \item Réponse : \zh{我们是走路回家的。} $\rightarrow$ Question : \zh{你们是\textbf{怎么}回家的？} (Zěnme = Comment/Par quel moyen). \textit{Note :} \cle{走路} (marcher) = manière.
\end{itemize}
\textbf{Conclusion :} \cle{是……的} = « C'est ... que... » (Cleft sentence). Trois piliers : \textbf{Action passée} + \textbf{Focus unique} + \textbf{的 final}. Oublier \cle{的} ou garder \cle{了} = 0 point à la question de grammaire.
\end{exoresolu}

\begin{methode}[Synthèse : Combiner les 4 structures dans un texte cohérent (Niveau Bac)]
\textbf{Objectif :} Produire un paragraphe (80-100 caractères) intégrant : Hypothèse (\cle{如果}), Nombres (\cle{两}/\cle{二}), Répartition (\cle{有的}), Focus (\cle{是……的}).
\textbf{Démarche en 5 étapes (Méthode « REDAC ») :}
\begin{enumerate}
    \item \textbf{R}éfléchir au \textbf{scénario} (Contexte Cameroun/Chine : rentrée, voyage, fête, résultats).
    \item \textbf{E}baucher le \textbf{plan} : 1 phrase intro (contexte/chiffres) $\rightarrow$ 1 phrase hypothèse (conseil/avenir) $\rightarrow$ 1 phrase répartition (réalités diverses) $\rightarrow$ 1 phrase focus (rappel fait marquant passé).
    \item \textbf{D}éployer le \textbf{vocabulaire} précis (Classificateurs, Tons, Mots-outils).
    \item \textbf{A}ssembler en \textbf{phrases complètes} (Ponctuation chinoise ， 。).
    \item \textbf{C}ontrôler : \cle{了} vs \cle{的} ? \cle{二} vs \cle{两} ? Ordre \cle{如果}... \cle{就} ? \cle{有的} répété ?
\end{enumerate}
\end{methode}

\begin{exoresolu}[Expression écrite intégrée : « Mon projet d'été »]
\textbf{Énoncé :} « Écris un court message (environ 80-100 caractères chinois) à ton ami chinois Liu Wei pour lui parler de ton été prochain. Tu dois \textbf{impérativement} utiliser :
\begin{enumerate}
    \item Une phrase avec \cle{如果……，就……} (ton projet conditionnel).
    \item Un nombre supérieur à 100 (ex: prix, distance, jours) + classificateur correct.
    \item La structure \cle{有的……有的……} (pour décrire ce que font tes camarades).
    \item Une phrase \cle{是……的} (pour raconter un souvenir marquant de l'an dernier).
\end{enumerate}
Vocabulaire aide : \zh{机票} (jīpiào), \zh{元} (yuán), \zh{实习} (shíxí), \zh{去年} (qùnián), \zh{难忘} (nánwàng). »

\tcblower

\textbf{Solution détaillée (Modèle de rédaction Bac — Note 18-20/20) :}

\zh{刘伟，你好！}
\zh{暑假有\textbf{两个月}，我想回村里看望奶奶。} \textit{(Nombres : 两个月 — classificateur 个 ; 两 obligatoire devant classif.)}
\zh{\textbf{如果}机票\textbf{不太贵}，我\textbf{就}买票回去。} \textit{(Hypothèse : 如果...就... ; 不太贵 = pas trop cher).}
\zh{同学\textbf{有的}去实习，\textbf{有的}留在雅温得打工。} \textit{(Répartition : 有的...有的... ; 留在 = rester à).}
\zh{\textbf{去年}我\textbf{是}坐火车\textbf{回去的}，很难忘。} \textit{(Focus : 是...的 sur le moyen 坐火车 ; 去年 = temps passé obligatoire).}
\zh{祝你暑假快乐！}
\zh{你的朋友，}
\zh{张明}

\textbf{Analyse grammaticale du modèle :}
\begin{itemize}
    \item \textbf{Phrase 1 (Nombres) :} \zh{两个月} (liǎng gè yuè) — Correct. \zh{两} + Classif. \zh{个} + Nom temps.
    \item \textbf{Phrase 2 (Hypothèse) :} \zh{如果...就...} — Structure parfaite. Négation \zh{不太贵} avant adjectif.
    \item \textbf{Phrase 3 (Répartition) :} \zh{同学} (Topique) + \zh{有的}... \zh{有的}... — Parallélisme parfait (Verbe + Lieu/Action).
    \item \textbf{Phrase 4 (Focus) :} \zh{去年} (Temps passé global) + \zh{是} + \zh{坐火车} (Focus Moyen) + \zh{回去的} (Verbe + \cle{的}). \textbf{Pas de \cle{了}.} \zh{难忘} (inoubliable) commentaire final.
\end{itemize}
\textbf{Conseil Bac :} Compte les caractères (ici $\approx$ 95). Vérifie chaque mot-outil (\cle{的, 了, 是, 在, 有}). Relis à voix haute pour les tons (\cle{liǎng}, \cle{jiù}, \cle{yǒude}, \cle{zuò}).
\end{exoresolu}

\begin{attention}[Les 5 erreurs qui coûtent des points au Bac (Leçon 6)]
\begin{enumerate}
    \item \textbf{\cle{如果}... \cle{那么} vs \cle{就} / Ordre des mots :} Écrire *\zh{那么如果我去，我去} (Inversion FR « Alors si j'y vais... »). \textbf{Réflexe :} \cle{如果} \textbf{toujours} en premier (ou après sujet), \cle{那么/就} \textbf{toujours} en début de 2e proposition. Ne jamais traduire « alors » par un mot isolé en début de phrase.
    \item \textbf{Nombres : \cle{二} vs \cle{两} et \cle{零} :} Dire *\zh{二个书} (au lieu de \zh{两本書}), *\zh{两月} (au lieu de \zh{二月}), *\zh{一百一} (oral toléré mais écrit \zh{一百一十} ou \zh{一百零一} selon contexte), *\zh{一万一千} pour 11 000 (correct) mais *\zh{一万零一千} (faux, pas de \cle{零} si pas de trou). \textbf{Réflexe :} Classificateur $\rightarrow$ \cle{两} ; Date/Numéro/Comptage pur $\rightarrow$ \cle{二}. Un seul \cle{零} par « trou ».
    \item \textbf{\cle{有的}... \cle{有的}... : Oubli de la répétition ou confusion avec \cle{有}.} Écrire *\zh{有的学生喜欢，其他的不喜欢} (Français « certains..., les autres... »). \textbf{Réflexe :} Structure symétrique obligatoire : \cle{有的} + Prédicat 1， \cle{有的} + Prédicat 2. \cle{其他的} (les autres) existe mais change la structure.
    \item \textbf{\cle{是……的} : Garder \cle{了} ou oublier \cle{的}.} *\zh{我是昨天去了北京的} (Double marqueur) ou *\zh{我是昨天去北京} (Focus non fermé). \textbf{Réflexe :} \cle{是……的} = « C'est... que... » (Passé). \cle{了} = Aspect accompli (Neutre). \textbf{Choisir un seul.} Question $\rightarrow$ Mot interrogatif \textbf{entre \cle{是} et \cle{的}}.
    \item \textbf{Classificateurs : \cle{个} « universel » abusif.} Dire *\zh{三个书}, *\zh{两个车}, *\zh{五个老师} (au lieu de \cle{位}). \textbf{Réflexe :} Apprendre les couples \textbf{Nom + Classif.} par cœur (\zh{书-本}, \zh{车-辆}, \zh{老师-位}, \zh{票-张}, \zh{衣服-件}, \zh{路-条}). Au Bac, 1 erreur de classif = -0,5 à -1 pt.
\end{enumerate}
\end{attention}

\newpage

\coursec{Exercices}

\courssub{Je vérifie mes connaissances}

\begin{exercice}[QCM : Mots-outils et structure]
Choisissez le mot qui convient pour compléter chaque phrase. Une seule réponse est correcte.
\begin{enumerate}
    \item 如果明天下雨，我们 \_\_\_\_\_\_ 不去爬山。
    \begin{itemize}
        \item A. 就 (jiù)
        \item B. 才 (cái)
        \item C. 也 (yě)
    \end{itemize}
    \item 公园里，\_\_\_\_\_\_ 人在跑步，\_\_\_\_\_\_ 人在下棋。
    \begin{itemize}
        \item A. 有的……有的 (yǒude……yǒude)
        \item B. 有些……有些 (yǒuxiē……yǒuxiē)
        \item C. 一个……一个 (yígè……yígè)
    \end{itemize}
    \item 这本书 \_\_\_\_\_\_ 我昨天在书店买 \_\_\_\_\_\_。
    \begin{itemize}
        \item A. 是……了 (shì……le)
        \item B. 是……的 (shì……de)
        \item C. 在……的 (zài……de)
    \end{itemize}
    \item 我们学校有一千 \_\_\_\_\_\_ 五十个学生。(1050)
    \begin{itemize}
        \item A. 零 (líng)
        \item B. 十 (shí)
        \item C. 百 (bǎi)
    \end{itemize}
\end{enumerate}
\end{exercice}

\begin{exercice}[Vrai ou Faux ? Justifiez]
Déterminez si les phrases suivantes sont correctes (V) ou incorrectes (F). Pour les phrases incorrectes, expliquez l'erreur (ordre des mots, classificateur, particule, nombre).
\begin{enumerate}
    \item 就我不去市场。
    \item 我是明天去的北京。
    \item 我有二个朋友。
    \item 一百五 (pour exprimer 105).
    \item 如果你不忙，那么我们去喝茶。
\end{enumerate}
\end{exercice}

\begin{exercice}[Vocabulaire : Les nombres et classificateurs]
Associez chaque nombre en chiffres à son écriture en caractères chinois corrects, puis choisissez le classificateur approprié pour l'objet indiqué.
\begin{center}
\begin{tabularx}{\linewidth}{|X|X|X|X|}
\hline
\rowcolor{popblueL} Chiffres & Caractères chinois & Objet & Classificateur \\
\hline
105 & & 学生 (étudiant) & \\
\hline
2300 & & 书 (livre) & \\
\hline
99 & & 车 (voiture) & \\
\hline
1008 & & 岁 (âge) & \\
\hline
\end{tabularx}
\end{center}
\textbf{Classificateurs à disposition :} 个 (gè), 本 (běn), 辆 (liàng), 岁 (suì).
\end{exercice}

\begin{exercice}[Pinyin $\rightarrow$ Caractères]
Écoutez (ou lisez) les syllabes en pinyin et écrivez les caractères chinois correspondants. Attention aux tons et à l'écriture des nombres.
\begin{enumerate}
    \item rúguǒ, nàme, jiù
    \item yǒude, yǒude
    \item shì, de
    \item yīqiān líng wǔshí
    \item liǎngqiān sānbǎi
\end{enumerate}
\end{exercice}

\courssub{Je m'entraîne}

\begin{exercice}[Traduction : Français $\rightarrow$ Chinois]
Traduisez les phrases suivantes en chinois (caractères + pinyin). Utilisez les structures étudiées.
\begin{enumerate}
    \item S'il fait beau demain, nous irons au marché.
    \item S'il est trop cher, je ne l'achète pas.
    \item Dans le parc, certains courent, d'autres jouent aux échecs.
    \item C'est hier que j'ai acheté ce dictionnaire à la librairie.
    \item Il y a mille cinquante-huit élèves dans notre lycée.
\end{enumerate}
\end{exercice}

\begin{exercice}[Transformation : Mise en relief avec \zh{是……的}]
Transformez les phrases suivantes en utilisant la structure \zh{是……的} pour mettre l'accent sur l'élément souligné (temps, lieu, manière, objet). Conservez le sens original.
\begin{enumerate}
    \item 我 \underline{昨天} 在商店买了一件衬衫。
    \item 他 \underline{去年} 去了中国。
    \item 我们 \underline{坐公交车} 来的。
    \item 她 \underline{用中文} 写了这封信。
    \item 他们 \underline{在喀麦隆} 认识的。
\end{enumerate}
\end{exercice}

\begin{exercice}[Fusion : Répartition avec \zh{有的……有的……}]
Regroupez les deux phrases en une seule phrase en utilisant la structure \zh{有的……有的……}.
\begin{enumerate}
    \item 一些同学唱歌。一些同学跳舞。
    \item 一些人买水果。一些人买鱼。
    \item 一些学生踢足球。一些学生看书。
    \item 一些商店卖衣服。一些商店卖鞋子。
    \item 有些老师教中文。有些老师教法文。
\end{enumerate}
\end{exercice}

\begin{exercice}[Texte à trous : Contexte camerounais]
Complétez le dialogue entre Paul et Li Mei avec les mots manquants : \zh{就}、\zh{那么}、\zh{有的}、\zh{有的}、\zh{是}、\zh{的}、\zh{零}、\zh{两}、\zh{百}、\zh{千}.
\begin{textebox}[Dialogue au marché de Mfoundi]
Paul: 听说这个周末有足球赛，\zh{如果} 不下雨，\zh{\_\_\_\_\_\_} 我们 \zh{\_\_\_\_\_\_} 一起去看吧？
Li Mei: 好主意！听说体育场能容纳 \zh{一\_\_\_\_\_\_\_\_\_\_} \zh{三\_\_\_\_\_\_\_\_\_\_} 人呢。
Paul: 哇，真大！平时 \zh{\_\_\_\_\_\_} 人卖饮料，\zh{\_\_\_\_\_\_} 人卖烤鱼，很热闹。
Li Mei: 对了，这双球鞋真漂亮。\zh{\_\_\_\_\_\_} 你 \zh{\_\_\_\_\_\_} 在这里买 \zh{\_\_\_\_\_\_} 吗？
Paul: 不是，\zh{我 \_\_\_\_\_\_ 上个星期在网上买 \_\_\_\_\_\_}。
\end{textebox}
\end{exercice}

\begin{exercice}[Nombres : Écriture et emploi]
\begin{enumerate}
    \item Écrivez en chiffres arabes, puis en caractères chinois (chiffres complexes 壹, 贰... facultatif, standard acceptés) :
    \begin{itemize}
        \item 三十五
        \item 一百零五
        \item 两千三百
        \item 九百九十九
    \end{itemize}
    \item Utilisez chacun de ces quatre nombres dans une phrase complète avec un nom et son classificateur approprié (ex: 35 étudiants, 105 livres, 2300 yuans, 999 marches).
\end{enumerate}
\end{exercice}

\begin{exercice}[Remise en ordre : Ordre des mots]
Remettez les mots dans l'ordre correct pour former une phrase grammaticale. Ajoutez la ponctuation.
\begin{enumerate}
    \item 去 / 如果 / 明天 / 我 / 有空 / 就 / 图书馆 / 。
    \item 这双 / 是 / 我 / 买 / 的 / 鞋 / 在 / 那个 / 小店 / 里 / 。
    \item 有的 / 学 / 中文 / 有的 / 学 / 英文 / 同学 / 同学 / 。
    \item 一千 / 我们 / 零 / 十 / 个 / 学生 / 班 / 里 / 有 / 。
\end{enumerate}
\end{exercice}

\courssub{Je me prépare au Bac}

\begin{exercice}[Compréhension de l'écrit : Texte original]
Lisez le texte suivant et répondez aux questions en français.
\begin{textebox}[Un week-end à Yaoundé]
周末，小马和他的中国朋友王伟去雅温得市中心逛街。\zh{如果} 天气好，\zh{他们就} 想去博物馆。结果早上下了一点雨，\zh{那么} 他们 \zh{就} 先去了一家餐馆吃\zh{ndolé}。餐馆里人很多，\zh{有的} 吃\zh{ndolé}，\zh{有的} 喝啤酒，\zh{有的} 看手机。饭后，雨停了。他们去了博物馆。王伟 \zh{是} 今年 \zh{一月} 来 \zh{喀麦隆} \zh{的}。他 \zh{说}：“这 \zh{是} 我 \zh{第一次} 吃 \zh{ndolé}，\zh{真} 好吃！”博物馆里大概 \zh{有} \zh{两千} \zh{多} 件展品。他们 \zh{看了} \zh{三个} \zh{小时} 才 出来。
\end{textebox}
\textbf{Questions :}
\begin{enumerate}
    \item Pourquoi n'ont-ils pas été au musée le matin ? (Structure \zh{如果……那么……})
    \item Décrivez l'ambiance au restaurant en utilisant la structure \zh{有的……有的……} (relevez la phrase du texte).
    \item Depuis quand Wang Wei est-il au Cameroun ? Citez la phrase exacte avec la structure \zh{是……的}.
    \item Combien y a-t-il d'objets exposés au musée environ ? (Écrivez le nombre en chiffres).
    \item Traduisez en français la phrase : \zh{“这是我第一次吃ndolé，真好吃！”}
\end{enumerate}
\end{exercice}

\begin{exercice}[Thème : Type Baccalauréat]
Traduisez le paragraphe suivant en chinois (caractères simplifiés + pinyin). Soignez les structures grammaticales (\zh{如果……就/那么……}、\zh{有的……有的……}、\zh{是……的}、 nombres).
\begin{quote}
« Demain, s'il ne pleut pas, nous irons au stade voir le match de football. Le stade est grand, il peut contenir deux mille cinq cents personnes. À l'entrée, certains vendent des billets, d'autres vérifient les sacs. C'est à huit heures précises que nous sommes arrivés au stade. »
\end{quote}
\textit{Indications : stade = 体育场 (tǐyùchǎng) ; match = 比赛 (bǐsài) ; contenir = 容纳 (róngnà) ; billet = 票 (piào) ; vérifier = 检查 (jiǎnchá) ; sac = 包 (bāo) ; précises = 整 (zhěng).}
\end{exercice}

\begin{exercice}[Version : Type Baccalauréat]
Traduisez le texte suivant en français courant et correct.
\begin{textebox}[Texte à traduire]
如果太贵，我就不买。商店里有一百多个客人，有的买衬衫，有的买鞋。我是上个星期买的这双鞋，很舒服。如果你也喜欢，这个周末我们一起去那家店吧？
\end{textebox}
\end{exercice}

\begin{exercice}[Expression écrite : Sujet de type Bac]
\textbf{Sujet :} « Les habitudes de consommation changent : entre marchés traditionnels et achats en ligne. »
Rédigez un court texte en chinois (environ \textbf{80 à 100 caractères chinois}) pour décrire une scène d'achat (marché ou en ligne) en utilisant \textbf{au moins trois} des structures suivantes :
\begin{itemize}
    \item \zh{如果……，就/那么……}
    \item \zh{有的……，有的……}
    \item \zh{是……的} (mise en relief temps/lieu/manière)
    \item Nombres complexes (centaines, milliers, \zh{零})
\end{itemize}
\textit{Pistes : Décrivez le lieu, les gens, ce qu'ils font, quand/ comment vous avez acheté quelque chose.}
\end{exercice}

\newpage

\coursec{Corrigés des exercices}

\begin{corrige}[Exercice 1 — QCM : Mots-outils et structure]
\textbf{Réponses :}
\begin{enumerate}
    \item \textbf{A. 就 (jiù)} — 如果明天下雨，我们就不去爬山。\\\emph{Rúguǒ míngtiān xià yǔ, wǒmen jiù bù qù pá shān.} « S'il pleut demain, nous n'irons pas faire de randonnée. » La structure \zh{如果……就……} exige \zh{就} (ou \zh{那么}) dans la principale ; \zh{才} marque la tardiveté, \zh{也} l'addition : aucun ne convient.
    \item \textbf{A. 有的……有的 (yǒude……yǒude)} — 公园里，有的人在跑步，有的人在下棋。\\\emph{Gōngyuán lǐ, yǒude rén zài pǎobù, yǒude rén zài xià qí.} « Certains courent, d'autres jouent aux échecs. » \zh{有的} + nom est la structure de répartition vue en leçon ; \zh{有些} existe mais n'est pas la forme ciblée ici.
    \item \textbf{B. 是……的 (shì……de)} — 这本书是我昨天在书店买的。\\\emph{Zhè běn shū shì wǒ zuótiān zài shūdiàn mǎi de.} « C'est hier que j'ai acheté ce livre à la librairie. » Mise en relief du temps : \zh{是……的}, sans \zh{了}.
    \item \textbf{A. 零 (líng)} — 我们学校有一千零五十个学生。\\\emph{Wǒmen xuéxiào yǒu yīqiān líng wǔshí gè xuésheng.} 1050 = \zh{一千零五十} : \zh{零} comble la centaine absente.
\end{enumerate}
\end{corrige}

\begin{corrige}[Exercice 2 — Vrai ou Faux ? Justifiez]
\begin{enumerate}
    \item \textbf{F} — 就我不去市场。 est incorrect : \zh{就} est un adverbe, il se place \emph{après} le sujet, jamais en tête. Correction : \zh{我就不去市场。} \emph{Wǒ jiù bù qù shìchǎng.} « Moi, je n'irai pas au marché. »
    \item \textbf{F} — 我是明天去的北京。 est incorrect : la structure \zh{是……的} ne sert que pour des actions \emph{passées} ; \zh{明天} est futur. Correction : \zh{我是昨天去的北京。} ou, pour le futur : \zh{我明天去北京。}
    \item \textbf{F} — 我有二个朋友。 est incorrect : devant un classificateur, « deux » se dit \zh{两} et non \zh{二}. Correction : \zh{我有两个朋友。} \emph{Wǒ yǒu liǎng gè péngyou.}
    \item \textbf{F} — 一百五 exprime 150, non 105. Pour 105, il faut \zh{零} : \zh{一百零五} \emph{yībǎi líng wǔ}.
    \item \textbf{V} — 如果你不忙，那么我们去喝茶。 est correcte : \zh{如果……那么……} bien formé, ordre des mots correct.
\end{enumerate}
\end{corrige}

\begin{corrige}[Exercice 3 — Vocabulaire : Les nombres et classificateurs]
\begin{center}
\begin{tabularx}{\linewidth}{|X|X|X|X|}
\hline
\rowcolor{popblueL} Chiffres & Caractères chinois & Objet & Classificateur \\
\hline
105 & 一百零五 (yībǎi líng wǔ) & 学生 & 个 (gè) → 一百零五个学生 \\
\hline
2300 & 两千三百 (liǎngqiān sānbǎi) & 书 & 本 (běn) → 两千三百本书 \\
\hline
99 & 九十九 (jiǔshí jiǔ) & 车 & 辆 (liàng) → 九十九辆车 \\
\hline
1008 & 一千零八 (yīqiān líng bā) & 岁 & 岁 (suì) → 一千零八岁 \\
\hline
\end{tabularx}
\end{center}
\textbf{Points de vigilance :} \zh{零} est obligatoire dans 105 et 1008 (rang intermédiaire vide) ; « deux mille » = \zh{两千} (\zh{两} devant \zh{千}) ; le classificateur des livres est \zh{本}, celui des véhicules \zh{辆}.
\end{corrige}

\begin{corrige}[Exercice 4 — Pinyin → Caractères]
\begin{enumerate}
    \item rúguǒ → \zh{如果} ; nàme → \zh{那么} ; jiù → \zh{就}
    \item yǒude → \zh{有的} (deux fois)
    \item shì → \zh{是} ; de → \zh{的}
    \item yīqiān líng wǔshí → \zh{一千零五十} (1050)
    \item liǎngqiān sānbǎi → \zh{两千三百} (2300)
\end{enumerate}
\textbf{Vigilance :} ne pas confondre \zh{是} (shì, être) et \zh{十} (shí, dix) ; \zh{的} se prononce \emph{de} (ton neutre) dans la structure \zh{是……的}.
\end{corrige}

\begin{corrige}[Exercice 5 — Traduction : Français → Chinois]
\begin{enumerate}
    \item 如果明天天气好，我们就去市场。\\\emph{Rúguǒ míngtiān tiānqì hǎo, wǒmen jiù qù shìchǎng.}
    \item 如果太贵，我就不买。\\\emph{Rúguǒ tài guì, wǒ jiù bù mǎi.}
    \item 公园里，有的人跑步，有的人下棋。\\\emph{Gōngyuán lǐ, yǒude rén pǎobù, yǒude rén xià qí.}
    \item 这本词典是我昨天在书店买的。\\\emph{Zhè běn cídiǎn shì wǒ zuótiān zài shūdiàn mǎi de.}
    \item 我们中学有一千零五十八个学生。\\\emph{Wǒmen zhōngxué yǒu yīqiān líng wǔshíbā gè xuésheng.}
\end{enumerate}
\textbf{Points de vigilance :} dans la phrase 4, pas de \zh{了} avec \zh{是……的} ; dans la phrase 5, \zh{零} est indispensable (1058 : la centaine est vide).
\end{corrige}

\begin{corrige}[Exercice 6 — Transformation : Mise en relief avec 是……的]
\begin{enumerate}
    \item 我是昨天在商店买了一件衬衫的。\\\emph{Wǒ shì zuótiān zài shāngdiàn mǎi le yí jiàn chènshān de.} (accent sur « hier »)
    \item 他是去年去中国的。\\\emph{Tā shì qùnián qù Zhōngguó de.} (accent sur « l'année dernière »)
    \item 我们是坐公交车来的。\\\emph{Wǒmen shì zuò gōngjiāochē lái de.} (accent sur le moyen de transport)
    \item 她是用中文写这封信的。\\\emph{Tā shì yòng Zhōngwén xiě zhè fēng xìn de.} (accent sur la langue utilisée)
    \item 他们是在喀麦隆认识的。\\\emph{Tāmen shì zài Kāmàilóng rènshi de.} (accent sur le lieu)
\end{enumerate}
\textbf{Règle rappelée :} \zh{是} se place juste \emph{avant} l'élément mis en relief ; \zh{的} se met en fin de phrase (ou après le verbe) ; le verbe garde éventuellement \zh{了} quand l'accent n'est pas sur l'accomplissement.
\end{corrige}

\begin{corrige}[Exercice 7 — Fusion : Répartition avec 有的……有的……]
\begin{enumerate}
    \item 有的同学唱歌，有的同学跳舞。\\\emph{Yǒude tóngxué chàng gē, yǒude tóngxué tiào wǔ.}
    \item 有的人买水果，有的人买鱼。\\\emph{Yǒude rén mǎi shuǐguǒ, yǒude rén mǎi yú.}
    \item 有的学生踢足球，有的学生看书。\\\emph{Yǒude xuésheng tī zúqiú, yǒude xuésheng kàn shū.}
    \item 有的商店卖衣服，有的商店卖鞋子。\\\emph{Yǒude shāngdiàn mài yīfu, yǒude shāngdiàn mài xiézi.}
    \item 有的老师教中文，有的老师教法文。\\\emph{Yǒude lǎoshī jiāo Zhōngwén, yǒude lǎoshī jiāo Fǎwén.}
\end{enumerate}
\textbf{Vigilance :} le nom (\zh{同学}, \zh{人}, \zh{商店}…) se répète après chaque \zh{有的} ; on ne dit pas \zh{有的唱歌，跳舞} sans sujet répété.
\end{corrige}

\begin{corrige}[Exercice 8 — Texte à trous : Contexte camerounais]
Dialogue complété :
\begin{itemize}
    \item Paul : 听说这个周末有足球赛，\zh{如果}不下雨，\zh{那么}我们\zh{就}一起去看吧？
    \item Li Mei : 好主意！听说体育场能容纳\zh{一千}三\zh{百}人呢。
    \item Paul : 哇，真大！平时\zh{有的}人卖饮料，\zh{有的}人卖烤鱼，很热闹。
    \item Li Mei : 对了，这双球鞋真漂亮。\zh{是}你\zh{的}…… → formulation correcte : \zh{是你在这里买的吗？}
    \item Paul : 不是，我\zh{是}上个星期在网上买\zh{的}。
\end{itemize}
\textbf{Ordre des mots placés :} 那么、就、千、百、有的、有的、是、的、是、的。\\\textbf{Remarque :} « mille trois cents » = \zh{一千三百} ; la question avec \zh{是……的……吗} interroge sur les circonstances de l'achat (ici le lieu).
\end{corrige}

\begin{corrige}[Exercice 9 — Nombres : Écriture et emploi]
\textbf{1. Écriture :}
\begin{itemize}
    \item 三十五 → 35
    \item 一百零五 → 105
    \item 两千三百 → 2300
    \item 九百九十九 → 999
\end{itemize}
\textbf{2. Phrases (modèles) :}
\begin{itemize}
    \item 我们班有三十五个学生。\emph{Wǒmen bān yǒu sānshíwǔ gè xuésheng.} « Il y a 35 élèves dans notre classe. »
    \item 图书馆有一百零五本中文书。\emph{Túshūguǎn yǒu yībǎi líng wǔ běn Zhōngwén shū.} « La bibliothèque a 105 livres de chinois. »
    \item 这部手机两千三百块。\emph{Zhè bù shǒujī liǎngqiān sānbǎi kuài.} « Ce téléphone coûte 2300 yuans. »
    \item 这座山有九百九十九级台阶。\emph{Zhè zuò shān yǒu jiǔbǎi jiǔshíjiǔ jí táijiē.} « Cette montagne a 999 marches. »
\end{itemize}
\textbf{Vigilance :} classificateurs adaptés : \zh{个} (personnes), \zh{本} (livres), \zh{块} (argent), \zh{级} (marches).
\end{corrige}

\begin{corrige}[Exercice 10 — Remise en ordre : Ordre des mots]
\begin{enumerate}
    \item 如果我明天有空，就去图书馆。\\\emph{Rúguǒ wǒ míngtiān yǒu kòng, jiù qù túshūguǎn.} « Si j'ai le temps demain, j'irai à la bibliothèque. »
    \item 这双鞋是我在那个小店里买的。\\\emph{Zhè shuāng xié shì wǒ zài nàge xiǎo diàn lǐ mǎi de.} « C'est dans cette petite boutique que j'ai acheté ces chaussures. »
    \item 有的同学学中文，有的同学学英文。\\\emph{Yǒude tóngxué xué Zhōngwén, yǒude tóngxué xué Yīngwén.}
    \item 我们班里有 一千零十个学生。→ phrase correcte : \zh{我们班有一千零十个学生。} \emph{Wǒmen bān yǒu yīqiān líng shí gè xuésheng.} (1010 élèves)
\end{enumerate}
\textbf{Méthode :} repérer d'abord le sujet, puis la subordonnée en \zh{如果}, placer \zh{就} avant le verbe de la principale, et \zh{的} en finale pour \zh{是……的}.
\end{corrige}

\begin{corrige}[Exercice 11 — Compréhension de l'écrit : Un week-end à Yaoundé]
\begin{enumerate}
    \item Ils n'ont pas été au musée le matin parce qu'il a plu un peu : \zh{如果天气好，他们就想去博物馆。结果早上下了一点雨，那么他们就先去了一家餐馆。} La condition (« s'il fait beau ») n'étant pas remplie, ils sont d'abord allés au restaurant.
    \item \zh{餐馆里人很多，有的吃ndolé，有的喝啤酒，有的看手机。} « Au restaurant, il y avait beaucoup de monde : certains mangeaient du ndolé, d'autres buvaient de la bière, d'autres regardaient leur téléphone. »
    \item Wang Wei est au Cameroun depuis janvier de cette année : \zh{王伟是今年一月来喀麦隆的。} La structure \zh{是……的} met en relief la date d'arrivée.
    \item Environ 2000 objets : \zh{大概有两千多件展品} → 2000 (et plus).
    \item \zh{这是我第一次吃ndolé，真好吃！} → « C'est la première fois que je mange du ndolé, c'est vraiment délicieux ! »
\end{enumerate}
\textbf{Barème indicatif (10 pts) :} 2 pts par question (1 pt sens, 1 pt exactitude de la citation ou de la structure).
\end{corrige}

\begin{corrige}[Exercice 12 — Thème : Type Baccalauréat]
\textbf{Proposition de traduction :}

明天如果不下雨，我们就去体育场看足球比赛。体育场很大，能容纳两千五百人。在门口，有的人卖票，有的人检查包。我们是八点整到体育场的。

\emph{Míngtiān rúguǒ bù xià yǔ, wǒmen jiù qù tǐyùchǎng kàn zúqiú bǐsài. Tǐyùchǎng hěn dà, néng róngnà liǎngqiān wǔbǎi rén. Zài ménkǒu, yǒude rén mài piào, yǒude rén jiǎnchá bāo. Wǒmen shì bā diǎn zhěng dào tǐyùchǎng de.}

\textbf{Barème indicatif (10 pts) :} \zh{如果……就……} correct (2 pts) ; nombre 2500 = \zh{两千五百} (2 pts) ; \zh{有的……有的……} (2 pts) ; \zh{是……的} avec \zh{八点整} (2 pts) ; exactitude générale et ponctuation chinoise (2 pts).

\textbf{Points de vigilance :} « deux mille cinq cents » = \zh{两千五百} (\zh{两}, pas \zh{二}) ; « à huit heures précises » = \zh{八点整} placé après \zh{是} ; pas de \zh{了} dans la phrase en \zh{是……的}.
\end{corrige}

\begin{corrige}[Exercice 13 — Version : Type Baccalauréat]
\textbf{Traduction de référence :}

« Si c'est trop cher, je ne l'achète pas. Il y a plus d'une centaine de clients dans le magasin : certains achètent des chemises, d'autres achètent des chaussures. C'est la semaine dernière que j'ai acheté cette paire de chaussures ; elles sont très confortables. Si toi aussi tu les aimes, ce week-end nous irons ensemble dans cette boutique, d'accord ? »

\textbf{Barème indicatif (10 pts) :} hypothèse \zh{如果……就……} (2 pts) ; nombre \zh{一百多个} « plus d'une centaine » (2 pts) ; répartition \zh{有的……有的……} (2 pts) ; \zh{是……的} rendu par une mise en relief française (2 pts) ; fluidité et correction du français (2 pts).

\textbf{Points de vigilance :} \zh{一百多个} = « plus de cent » (\zh{多} = « et plus ») ; \zh{这双鞋} = « cette paire de chaussures » ; \zh{很舒服} qualifie les chaussures, pas l'achat.
\end{corrige}

\begin{corrige}[Exercice 14 — Expression écrite : Sujet de type Bac]
\textbf{Proposition de rédaction modèle (environ 95 caractères) :}

上个星期天，我和妈妈去了莫foundi市场。如果下雨，我们就在网上买东西。市场里人很多，有的卖水果，有的卖鱼，还有的卖衣服。妈妈买了一条鱼，二十块钱。我的鞋是上个星期在网上买的，很方便，也很便宜。

\emph{Shàng ge xīngqītiān, wǒ hé māma qù le shìchǎng. Rúguǒ xià yǔ, wǒmen jiù zài wǎngshàng mǎi dōngxi. Shìchǎng lǐ rén hěn duō, yǒude mài shuǐguǒ, yǒude mài yú, hái yǒude mài yīfu. Māma mǎi le yì tiáo yú, èrshí kuài qián. Wǒ de xié shì shàng ge xīngqī zài wǎngshàng mǎi de, hěn fāngbiàn, yě hěn piányi.}

« Dimanche dernier, je suis allé au marché avec maman. S'il pleut, nous achetons en ligne. Au marché, il y avait beaucoup de monde : certains vendaient des fruits, d'autres du poisson, d'autres encore des vêtements. Maman a acheté un poisson pour vingt yuans. Mes chaussures, c'est la semaine dernière que je les ai achetées en ligne : c'est pratique et bon marché. »

\textbf{Barème indicatif (20 pts) :} respect de la longueur 80–100 caractères (3 pts) ; au moins trois structures ciblées correctement employées (3 × 3 = 9 pts) ; exactitude des caractères et des tons (4 pts) ; cohérence du récit et ponctuation chinoise (4 pts).

\textbf{Points de vigilance :}
\begin{itemize}
    \item \zh{如果……就……} : \zh{就} après le sujet de la principale ;
    \item \zh{有的……有的……} : répéter le nom ou utiliser \zh{有的} seul si le nom est déjà connu ;
    \item \zh{是……的} : uniquement pour un achat \emph{déjà fait}, sans \zh{了} final ;
    \item classificateurs : \zh{条} pour le poisson, \zh{双} pour les chaussures.
\end{itemize}
\end{corrige}

\newpage

\coursec{Fiche bilan}

\begin{popfigure}
\begin{center}\tcbox[colback=popink,colframe=popink,coltext=white,arc=9pt,boxsep=4pt,left=6pt,right=6pt]{\bfseries\small Leçon 6 Grammaire}\end{center}
\vspace{-2pt}
\begin{tcbraster}[raster columns=3,raster equal height=rows,raster column skip=6pt,raster row skip=6pt,enhanced,blankest]
\begin{tcolorbox}[enhanced,colback=popblue,colframe=popblue,coltext=white,boxrule=0.9pt,arc=3pt,left=4pt,right=4pt,top=3pt,bottom=3pt,boxsep=1pt,halign=left]\bfseries\scriptsize Hypothèse \zh{如果\ldots 那么/就}\end{tcolorbox}
\begin{tcolorbox}[enhanced,colback=poporange,colframe=poporange,coltext=white,boxrule=0.9pt,arc=3pt,left=4pt,right=4pt,top=3pt,bottom=3pt,boxsep=1pt,halign=left]\bfseries\scriptsize Les nombres \zh{数字} (\zh{万}, \zh{亿})\end{tcolorbox}
\begin{tcolorbox}[enhanced,colback=popgreen,colframe=popgreen,coltext=white,boxrule=0.9pt,arc=3pt,left=4pt,right=4pt,top=3pt,bottom=3pt,boxsep=1pt,halign=left]\bfseries\scriptsize Répartition \zh{有的\ldots 有的\ldots}\end{tcolorbox}
\begin{tcolorbox}[enhanced,colback=poppurple,colframe=poppurple,coltext=white,boxrule=0.9pt,arc=3pt,left=4pt,right=4pt,top=3pt,bottom=3pt,boxsep=1pt,halign=left]\bfseries\scriptsize Mise en relief \zh{是\ldots 的} (passé)\end{tcolorbox}
\begin{tcolorbox}[enhanced,colback=poppink,colframe=poppink,coltext=white,boxrule=0.9pt,arc=3pt,left=4pt,right=4pt,top=3pt,bottom=3pt,boxsep=1pt,halign=left]\bfseries\scriptsize Vigilance ordre, \zh{的/得/地}\end{tcolorbox}
\end{tcbraster}

\legende{Carte mentale de la Leçon 6 : les 4 structures clés}
\end{popfigure}

\begin{bilan}
\end{bilan}

\courssub{Lexique clé (HSK 2--3)}
\begin{center}
\begin{tabularx}{\linewidth}{|X|X|X|X|}
\hline \rowcolor{popblueL} Caractères & Pinyin & Nature & Traduction \\ \hline
如果 & rúguǒ & conj. & si, au cas où \\ \hline
那么 & nàme & adv. & alors, dans ce cas \\ \hline
就 & jiù & adv. & alors, tout de suite, dès que \\ \hline
有的 & yǒude & pron. & certains, quelques-uns \\ \hline
是 & shì & verbe & être (marqueur d'emphase) \\ \hline
的 & de & part. & particule de structure / emphase finale \\ \hline
零 & líng & num. & zéro \\ \hline
万 & wàn & classif. & dix mille ($10^4$) \\ \hline
亿 & yì & classif. & cent millions ($10^8$) \\ \hline
两 & liǎng & num. & deux (devant classif./grands nombres) \\ \hline
\end{tabularx}
\end{center}

\courssub{Règles de grammaire à connaître par cœur}
\begin{center}
\begin{tabularx}{\linewidth}{|X|X|X|}
\hline \rowcolor{popblueL} Structure & Exemple (Caractères + Pinyin) & Traduction \\ \hline
\zh{如果 + Condition, 那么/就 + Conséquence} & \zh{如果下雨，那么我不去。} (Rúguǒ xiàyǔ, nàme wǒ bù qù.) & S'il pleut, alors je n'y vais pas. \\ \hline
\zh{如果 + Condition, 就 + Conséquence (certaine/immédiate)} & \zh{如果你去，我就去。} (Rúguǒ nǐ qù, wǒ jiù qù.) & Si tu y vas, j'irai (moi aussi). \\ \hline
Nombres : groupement par 4 (\zh{万}, \zh{亿}) & \zh{一万两千零五} (yī wàn liǎng qiān líng wǔ) & 12 005 (\cle{两} pour 2000, \cle{零} unique) \\ \hline
\zh{有的 + V1, 有的 + V2} (sujet commun omis) & \zh{有的吃饭，有的睡觉。} (Yǒude chīfàn, yǒude shuìjiào.) & Certains mangent, d'autres dorment. \\ \hline
\zh{有的 + Suj A + V, 有的 + Suj B + V} & \zh{有的同学看书，有的老师喝茶。} & Certains élèves lisent, certains profs boivent du thé. \\ \hline
\zh{是 + [Temps/Lieu/Manière] + Verbe + 的} (passé accompli) & \zh{我是坐飞机去的北京。} (Wǒ shì zuò fēijī qù de Běijīng.) & C'est en avion que je suis allé à Pékin. \\ \hline
Négation \zh{是\ldots 的} & \zh{我不是坐火车去的。} (Wǒ bùshì zuò huǒchē qù de.) & Ce n'est pas en train que j'y suis allé. \\ \hline
\end{tabularx}
\end{center}

\courssub{Méthodes express}
\begin{enumerate}
\item \textbf{Conditionnel (\zh{如果\ldots})} : Identifier la proposition condition (après \cle{如果}) et la conséquence (après \cle{那么} ou \cle{just}). \textbf{Choix} : \cle{那么} = conséquence logique/générale (\og dans ce cas\fg{}) ; \cle{就} = conséquence certaine, immédiate ou inévitable (\og du coup\fg{}). Ne pas mettre \cle{会} (huì) après \cle{就}.
\item \textbf{Nombres} : Séparer en groupes de 4 chiffres (unités \zh{万}, \zh{亿}). Lire \zh{零} (líng) une seule fois pour des zéros consécutifs au milieu. Utiliser \zh{两} (liǎng) pour 2 devant \zh{百/千/万/亿} et classif. ; \zh{二} (èr) pour compter, maths, 2e étage, février.
\item \textbf{Répartition (\zh{有的\ldots{}有的\ldots})} : Exprime \og les uns\ldots les autres\fg{}. Si le sujet est le même (groupe), omettre le sujet : \zh{有的V, 有的V}. Si sujets différents, les expliciter : \zh{有的S1 V, 有的S2 V}.
\item \textbf{Emphase passée (\zh{是\ldots 的})} : Souligne \textbf{comment, quand, où, pourquoi, avec qui} une action \textbf{terminée} s'est déroulée. \cle{的} obligatoire en fin de phrase. Négation : \zh{不是\ldots 的}. Ne pas utiliser pour une action future ou habituelle.
\end{enumerate}


\begin{aretenir}[Checklist avant le Bac]
$\square$ Je construis \zh{如果\ldots, 那么/就\ldots} sans inverser l'ordre (condition d'abord).\\
$\square$ Je distingue \cle{那么} (logique) et \cle{就} (immédiat/certain) et j'évite \zh{会} après \cle{就}.\\
$\square$ J'écris et lis les nombres jusqu'aux \zh{亿} en groupant par 4 chiffres (\zh{万}, \zh{亿}).\\
$\square$ J'insère un seul \cle{零} pour des zéros consécutifs au milieu (ex: 10 005 \textrightarrow \zh{一万零五}).\\
$\square$ J'utilise \zh{两} (liǎng) devant \zh{百/千/万/亿} et classif., \zh{二} (èr) sinon.\\
$\square$ J'alterne \zh{有的\ldots 有的\ldots} pour répartir sans répéter le sujet commun.\\
$\square$ Je place \cle{有的} devant le verbe (sujet commun) ou devant le sujet (sujets différents).\\
$\square$ Je forme l'emphase passée \zh{是\ldots 的} uniquement sur action \textbf{accomplie}.\\
$\square$ Je place \cle{的} \textbf{à la fin} de la phrase \zh{是\ldots 的} (jamais après le verbe seul).\\
$\square$ Je ne confonds pas \cle{的} (nominal), \cle{得} (complément de degré), \cle{地} (adverbial).\\
$\square$ Je traduis \og Si j'ai le temps, je viendrai \fg{} par \zh{如果我有时间，我就来} (pas \zh{我会来}).\\
$\square$ Je sais écrire dans l'ordre les traits des clés : \zh{女} (nǚ), \zh{子} (zǐ), \zh{宀} (mián), \zh{辶} (chuò).
\end{aretenir}

\findecours

\end{document}
